\newcommand{\ogdcC}{\mathcal C}
\newcommand{\ogdDMod}{\operatorname{DMod}}
\newcommand{\ogdOrd}{\operatorname{Ord}}
\newcommand{\ogdOSurj}{\mathsf{OSurj}}
\newcommand{\ogdAlg}{\operatorname{Alg}}
\newcommand{\ogdCoalg}{\operatorname{Coalg}}
\newcommand{\ogdMod}{\operatorname{Mod}}
\newcommand{\ogdComod}{\operatorname{Comod}}
\newcommand{\ogdSym}{\operatorname{Sym}}
\newcommand{\ogdCE}{\operatorname{CE}}
\newcommand{\ogdTW}{\operatorname{TW}}
\newcommand{\ogdid}{\mathrm{id}}
\newcommand{\ogdone}{K}
\newcommand{\ogdUU}{\mathscr U}
\newcommand{\ogdotex}{\mathbin{\otimes^!}}
\newcommand{\ogdotexd}{\mathbin{\otimes^{\mathbb D}}}
\newcommand{\ogdciteFG}{\cite{chhop:FG}}
\newcommand{\ogdciteAKY}{\cite{S4aAlfonsiKimYoung}}

\chapter{Ordered collisions and associative chiral bar--cobar duality}
\label{ogd:foundation}

How should ordered field insertions combine when nonadjacent positions
coincide? Let\phantomsection\label{ogd:def:initial-curve-base} $k$ be a field of characteristic zero and initially let
$X$ be a smooth curve over $k$. A word of three positions can have the form
$(x,y,x)$ with $x\ne y$. Its two consecutive cuts are
\[
 (x)\mid(y,x),\qquad (x,y)\mid(x).
\]
Neither cut separates disjoint positions: $x$ occurs on both sides.
The middle insertion still occupies its place in the word.

An operator calculation detects its loss. On $k^2$, let $E_{ij}$ send
the $j$th basis vector to the $i$th and annihilate the other basis vector.
Then
\[
 E_{12}E_{21}E_{12}=E_{12},\qquad E_{12}^2=0.
\]
Identifying the outer positions does not permit deleting the intervening
operation. The tensor product must retain the ordered inputs and every
equality that crosses a cut.

We first construct this cut product on word coefficients and compute
its first bar and cobar terms. We then add the order in which coincident
points collide by passing to the full ordered collision atlas.
That atlas has a symmetric convolution summed over arbitrary ordered
decompositions. Projection onto consecutive decompositions gives the
explicit comparison with the cut product.

\section{Ordered words and their cut product}
\label{ogd:sec:cut-support}

The coefficients must record position derivatives and must restrict
coherently when adjacent positions coincide. Differential modules supply
these operations. Their derived restriction also retains the additional
cohomology of a collision.

\subsection{Coefficient complexes and geometric operations}

For\phantomsection\label{ogd:def:right-differential-module-conventions} a smooth scheme \(Y\), let \(\ogdDMod(Y)\) be the stable enhancement of the unbounded derived category of right algebraic \(\mathcal D_Y\)-modules.
We require their underlying \(\mathcal O_Y\)-modules to be quasi-coherent.
An enhancement retains mapping complexes and coherent homotopies, in addition to the triangulated category.
We use extraordinary inverse image \(f^!\), differential-module direct image \(f_+\), and external product \(\boxtimes\).
For an open immersion \(j\), write \(j_*\) for \(j_+\).
The notation \(Rq_*^{\mathcal O}\) will mean ordinary derived sheaf direct image.

On\phantomsection\label{oco:def:differential-operator-sheaf} a coordinate chart of a smooth curve, the algebra \(\mathcal D\)
is generated by regular functions and the coordinate derivative, with
\(\partial_z f-f\partial_z=f'\).
A right differential module is a sheaf with a right action of this algebra.
The actions agree under coordinate changes.
External\phantomsection\label{ogd:def:external-coefficient-product} products put such coefficients at independent positions.
Diagonal direct images describe coefficients supported at collisions.

Restricting to a diagonal can produce additional cohomology.
The inverse image must therefore be derived.
The direct image, in the other direction, must include normal
derivatives of distributions supported on the diagonal.

Transfer\phantomsection\label{ogd:def:transfer-dimensions-and-form-lines} modules define these two operations. For a map
$f:Z\to Y$ between smooth spaces, let $d_Z,d_Y$ be their dimensions
and let $\omega_Z,\omega_Y$ be their top differential-form sheaves.
The notation $f^{-1}$ means inverse image of sheaves, obtained by
pulling back germs of sections. The line $\omega_Y^{-1}$ is the
dual of $\omega_Y$.

The\phantomsection\label{ogd:def:right-transfer-bimodule} transfer module
$\mathcal D_{Z\to Y}=\mathcal O_Z\otimes_{f^{-1}\mathcal O_Y}
 f^{-1}\mathcal D_Y$ has a left $\mathcal D_Z$-action
and a right $f^{-1}\mathcal D_Y$-action. Put
$\mathcal D_{Y\leftarrow Z}
=\omega_Z\otimes_{\mathcal O_Z}\mathcal D_{Z\to Y}
 \otimes_{f^{-1}\mathcal O_Y}f^{-1}\omega_Y^{-1}$.
In local volume coordinates, transposition fixes functions, negates
coordinate derivatives, and reverses products.
It gives this module a left $f^{-1}\mathcal D_Y$-action and
a right $\mathcal D_Z$-action.

For complexes $M$ on $Z$ and $N$ on $Y$, the direct and
exceptional inverse images are
\[
 f_+M=Rf_*^{\mathcal O}\bigl(M\otimes^{\mathbb L}_{\mathcal D_Z}
                         \mathcal D_{Z\to Y}\bigr),\qquad
 f^!N=\bigl(f^{-1}N\otimes^{\mathbb L}_{f^{-1}\mathcal D_Y}
                  \mathcal D_{Y\leftarrow Z}\bigr)[d_Z-d_Y].
\]
The superscript $\mathbb L$ denotes a derived tensor product.
One factor is replaced by a resolution whose tensor product preserves
quasi-isomorphisms in the other factor. A bounded free resolution is
one example. The notation $Rf_*^{\mathcal O}$ denotes the derived sheaf direct image.
Ordinary restriction alone does not compute these derived operations.

For an open inclusion $j$, $j^!$ is restriction. For an open complement
of diagonals, $j_*j^!$ permits poles along those diagonals.
For\phantomsection\label{ogd:def:closed-differential-direct-image} a closed inclusion $i$, write $i_*=i_+$. This direct image includes
the normal derivatives of distributions supported on the image,
and $i^!i_*\simeq1$.

For the structure map $a:X\to\operatorname{Spec}k$,
the inverse-image formula gives the dualizing complex
$a^!k=\omega_X[1]$.
Let $i_x$ include a $k$-rational point $x\in X(k)$.
For a complex $N$ on $X$, its exceptional fiber at $x$ is $i_x^!N$.
Choose\phantomsection\label{ogd:def:regular-germ-ring} a local coordinate $u$ vanishing at $x$, and write
$\mathcal O_{X,x}$ for the ring of germs of regular functions at $x$.

Using $du$ to trivialize the form lines, the point transfer module
has a two-term free resolution with differential given by right
multiplication by $u$. Tensoring with $i_x^{-1}\omega_X$ gives
$\mathcal O_{X,x}\xrightarrow{u}\mathcal O_{X,x}$ in degrees $-1,0$.
Its kernel is zero and its cokernel is $k$.
The shift $[-1]$ in $i_x^!$ cancels the shift $[1]$ in $\omega_X[1]$,
so $i_x^!(\omega_X[1])=k$ in degree zero.

These functors carry their standard coherent composition maps.
For the structure morphism \(p_Y:Y\to\operatorname{Spec}k\), the smooth right dualizing object is
\[
 J_Y=p_Y^!k=\omega_Y[\dim Y],
\]
where \(\omega_Y\) is the canonical line with its right differential-module action.
Our normalization fixes all extraordinary shifts by this identity.
External products satisfy \(J_Y\boxtimes J_Z=J_{Y\times Z}\).
The coherent functorial framework is the one described in \ogdciteFG*{\S1.4}.

\subsection{The open base-change identity}

The product will extend sections from a disjoint locus.
Its collision compatibility requires one specific base-change identity.
For\phantomsection\label{ogd:def:closed-complement-inclusion} an open immersion \(j:U\hookrightarrow Y\), let \(i:Z\hookrightarrow Y\) be the reduced closed complement.
For possibly singular \(Z\), define \(\ogdDMod(Z)\) by descent from differential modules supported on local smooth embeddings of \(Z\), using Kashiwara's equivalence.
Open--closed recollement gives
\begin{equation}\label{ogd:eq:recollement}
 i_*i^!N\longrightarrow N\longrightarrow j_*j^!N,
 \qquad i^!j_*=0,\qquad j^!j_*=\ogdid.
\end{equation}
The same triangle descends over smooth chart covers, including those of
nonseparated schemes.

\begin{lemma}\label{ogd:lem:basechange}
For a Cartesian square with \(j\) open,
\[
\begin{tikzcd}
 U'\arrow[r,"f_U"]\arrow[d,"j'"']&U\arrow[d,"j"]\\
 Y'\arrow[r,"f"']&Y,
\end{tikzcd}
\qquad
 f^!j_*M\simeq j'_*f_U^!M.
\]
These equivalences are coherent under composition of the maps \(f\).
\end{lemma}
\begin{proof}
Let \(i':Z'\hookrightarrow Y'\) be the complementary closed immersion.
The map \(fi'\) factors through \(i\), say \(fi'=ih\).
Thus
\[
 i'^!f^!j_*M=h^!i^!j_*M=0.
\]
The triangle~\eqref{ogd:eq:recollement} on \(Y'\) identifies \(f^!j_*M\) with
\(j'_*j'^!f^!j_*M\).
Transitivity of extraordinary inverse image identifies the latter object with \(j'_*f_U^!M\).

For a composite \(f\), every comparison lies in the fully faithful image of \(j'_*\).
Restriction to \(U'\) reduces the comparison to transitivity of extraordinary inverse image.
Full faithfulness lifts the same coherent comparison uniquely, including higher homotopies.
\end{proof}

\subsection{Word coefficients and consecutive cuts}

Write\phantomsection\label{ogd:def:ground-derived-category}\phantomsection\label{ogd:def:derived-vector-space-category} $D(k)$ for complexes of $k$-vector spaces up to quasi-isomorphism.

For\phantomsection\label{ogd:def:word-coefficient-category} the word construction, now allow $X$ to be a smooth separated
finite-type scheme of any positive pure dimension over $k$.
Let $\operatorname{OrdSurj}_+$ have
nonempty finite ordered sets and monotone surjections as arrows, and put
\[
 \mathcal W_X=D(k)\times
 \varprojlim_{\operatorname{OrdSurj}_+}\ogdDMod(X^n).
\]
The transition for $\alpha:[n]\twoheadrightarrow[m]$ is extraordinary
inverse image along $d_\alpha(x)_i=x_{\alpha(i)}$.
Objects include these identifications and their coherences. The empty
coefficient is isolated. We use unbounded right differential modules with
quasi-coherent underlying modules and derived external tensor products.

For\phantomsection\label{ogd:def:consecutive-cut-open} $1\leq p<n$, define the cut open and its inclusion by
\[
 U_p=\{x_i\ne x_j\text{ for all }i\leq p<j\}\subset X^n,
 \qquad j_p:U_p\hookrightarrow X^n.
\]

The cut product is
\begin{equation}\label{ogd:eq:cut-product}
 (F\star_{\mathrm{cut}}G)_n
 =\bigoplus_{p=0}^{n}j_{p,*}
       (F_p\boxtimes G_{n-p})|_{U_p}.
\end{equation}
An empty block means scalar tensoring on the whole space.
The unit $K$ is $k$ at the empty configuration and zero elsewhere.

Every cross-block diagonal is removed in $U_p$, including nonadjacent
pairs. Under a monotone collapse, a cut inside a fiber has empty inverse
image. A cut between fibers descends to the corresponding quotient cut.
Lemma~\ref{ogd:lem:basechange} therefore makes
\eqref{ogd:eq:cut-product} a product of coefficient systems.

Its iterates are indexed by consecutive partitions, with all blocks
pairwise disjoint. This common expression supplies the associativity
coherences. No symmetry of the cut product is assumed.

A coaugmented cut coalgebra has a coproduct
$C\to C\star_{\mathrm{cut}}C$, a counit $C\to K$ and a
coaugmentation $K\to C$. These maps satisfy the coalgebra identities
and their coherent homotopies in $\mathcal W_X$.

\phantomsection\label{ogd:def:strong-cut-coalgebra}A cut coalgebra is \emph{strong} if every supplied map
\[
 C_n|_{U_p}\longrightarrow(C_p\boxtimes C_{n-p})|_{U_p}
\]
is an equivalence. Inverting these maps gives associative multiplications
on the disjoint correspondences. Conversely, open adjunction turns the
inverse of such a multiplication into a coproduct component.

On every final consecutive partition, these two constructions invert the
same coherent diagram, so they are inverse on mapping spaces as well.
This is the associative version of strong nonunital factorization in
\cite{interfaces-Lat}*{Definition~1.1.6}. Remark~A.1.4 there states the
$E_\infty$ version. The isolated empty unit above is additional to the
nonempty index, and is not a local field vacuum.

Define the reduced closed cut support
\begin{equation}\label{ogd:eq:Z-support}
 Z_1=X,\qquad
 Z_n=X^n\setminus\bigcup_{p=1}^{n-1}U_p
     =\bigcap_{p=1}^{n-1}\ \bigcup_{i\leq p<j}\Delta_{ij}.
\end{equation}
A tuple lies in $Z_n$ exactly when every cut has an equality crossing it.
In particular,
\begin{equation}\label{ogd:eq:Z-small}
 Z_2=\Delta_{12},\qquad Z_3=\Delta_{13},\qquad
 Z_4=\Delta_{14}\cup(\Delta_{13}\cap\Delta_{24}).
\end{equation}

For three positions, failure of both cuts forces $x_1=x_3$.
For four, $x_1=x_4$ blocks every cut.
If $x_1\ne x_4$, failure of the first and last cuts leaves the possibilities
$12,34$ or $13,24$: the other two choices force $14$.
The choice $12,34$ leaves the middle cut disjoint unless it too forces
$14$. Hence only $13,24$ remains. These are identities of reduced supports.

The same cut calculation proves
\begin{equation}\label{ogd:eq:Z-pullback}
 d_\alpha^{-1}Z_n=Z_m
 \quad\text{for every monotone }\alpha:[n]\twoheadrightarrow[m].
\end{equation}
Indeed, the surviving cut opens under pullback are exactly the quotient
cut opens. Thus local cohomology on $Z_n$ is compatible with the word
coefficient identifications. We write\phantomsection\label{ogd:def:supported-local-cohomology} $R\Gamma_Z M=i_*i^!M$ for the
local-cohomology term of the localization triangle for a reduced closed
support $i:Z\hookrightarrow Y$.

\subsection{Cohomological bounds on word coefficients}

Cohomological truncation must respect collision identifications.
It cannot be performed separately on every coefficient of a word
object. The standard bounds on those coefficients nevertheless
determine one half of a t-structure.
Accessibility means that the subcategory of nonpositive objects
is presentable; our indexing is cohomological.

\begin{proposition}\label{ogd:prop:word-t-structure}
Suppose that $X$ is a nonempty smooth separated finite-type curve over $k$.
There is an accessible t-structure on $\mathcal W_X$ for which
\[
 \mathcal W_X^{\geq0}
 =\{F:F_n\in\ogdDMod(X^n)^{\geq0}\text{ for every }n\geq0\},
\]
with $\ogdDMod(X^0)=D(k)$ for the isolated empty coefficient.
Each evaluation preserves the nonnegative half of the t-structure.
If the coefficients of $F$ and $G$ have lower bounds $a$ and $b$,
those of $F\star_{\mathrm{cut}}G$ have lower bound $a+b$.
The analogous statement holds for pointwise upper bounds.

Componentwise standard truncation does not define truncation on
$\mathcal W_X$. In particular, there is no t-structure on this
category for which every evaluation is t-exact with respect to
the standard t-structures of the coefficients.
\end{proposition}

\begin{proof}
The closed-diagonal transitions $d_\alpha^!$ preserve colimits and
are right adjoint to the corresponding direct images
\cite{chhop:FG}*{\S\S1.4.1--1.4.6}.
They therefore preserve limits as well. The word category is
presentable and stable, and its evaluations preserve limits and
colimits. Hence evaluation at $n$ has a left adjoint $L_n$;
this is the limit-category construction of
\cite{chhop:FG}*{\S2.1.3}, applied to monotone surjections.

Choose a small set $\mathcal S_n$ generating
$\ogdDMod(X^n)^{\leq0}$ under colimits and extensions, including
$n=0$. Generate an aisle $\mathcal U$ from all $L_nE$ with
$E\in\mathcal S_n$. The presentable stable generation theorem
\cite{cke:Lurie}*{Proposition~1.4.4.11} supplies an accessible
t-structure with $\mathcal W_X^{\leq0}=\mathcal U$.
By adjunction,
\[
 F\in\mathcal U^\perp
 \quad\Longleftrightarrow\quad
 \operatorname{Map}(E,F_n)\simeq *
 \ \text{for all }n\text{ and }E\in\mathcal S_n
 \quad\Longleftrightarrow\quad F_n\in\ogdDMod(X^n)^{\geq1}.
\]
This identifies the asserted nonnegative half after shifting.
External tensor products over the field $k$ add cohomological
bounds. On a curve each $j_{p,*}j_p^!$ is exact localization at
the product of the cross-pair differences. The finite sum in
\eqref{ogd:eq:cut-product} proves both tensor-bound assertions.
No description of $\mathcal W_X^{\leq0}$ by pointwise upper
bounds is used.

For the truncation assertion, let $p_n:X^n\to\operatorname{Spec}k$
be the structure map and take
\[
 F_0=0,\qquad F_n=p_n^!k[-1]=\omega_{X^n}[n-1]\quad(n>0).
\]
Transitivity of extraordinary inverse image gives all collision
identifications. Componentwise $\tau^{\geq0}$ retains
$F_1=\omega_X$ and kills $F_2=\omega_{X^2}[1]$.
These values cannot satisfy $d_{[2]\to[1]}^!F_2\simeq F_1$.
If every evaluation were t-exact, it would commute with
truncation and give precisely those incompatible values.
\end{proof}

\subsection{Bar words and a second algebra}
\label{ogd:sec:first-cut-words}

The cut product determines which tensor powers enter a bar construction.
Let\phantomsection\label{ogd:def:augmented-cut-algebra} $A$ be an augmented $E_1$ algebra in
$(\mathcal W_X,\star_{\mathrm{cut}},K)$, with $A_\varnothing=k$.
Its augmentation ideal is $a=\operatorname{fib}(A\to K)$, so
$a_\varnothing=0$ and the unit splits $A\simeq K\oplus a$.

Write\phantomsection\label{ogd:def:cut-bar-suspension} $s=[1]$ for the cohomological shift, with $|sa|=|a|-1$.
On a chain model, the normalized bar has underlying graded word object
\begin{equation}\label{ogd:eq:first-cut-bar}
 B_{\mathrm{cut}}A
   =K\oplus\bigoplus_{r\geq1}(sa)^{\star_{\mathrm{cut}}r}.
\end{equation}
At arity $n$ it contains at most $n$ positive factors.
Its coefficients are therefore finite sums at each configuration.

On\phantomsection\label{ogd:def:cut-bar-differential} a strict differential graded model, the unary and binary parts of
the bar differential are
\[
 b_1(sa)=-s(da),\qquad b_2(sa,sb)=(-1)^{|a|}s(ab).
\]
The full differential applies these degree-one maps at all possible
positions, with the sign from passing the preceding suspended elements.
For closed elements of degree zero, the first two calculations are
\begin{equation}\label{ogd:eq:first-cut-bar-faces}
 b[a|b]=[ab],\qquad
 b[a|b|c]=[ab|c]-[a|bc].
\end{equation}
Applying $b$ again to the word of length three gives
$[(ab)c]-[a(bc)]=0$. This is where associativity enters the bar
differential. On a general $E_1$ model, higher multiplication faces
carry its associativity homotopies and their compatibilities.

Each word denotes its localized geometric coefficient, so these
identities describe coefficient maps, not a choice of global pure tensors.

The bar is a coalgebra. Its coproduct retains every consecutive cut,
including either empty end. On a word of length two it is
\begin{equation}\label{ogd:eq:first-cut-deconcatenation}
 \Delta[a|b]
 =[]\otimes[a|b]+[a]\otimes[b]+[a|b]\otimes[].
\end{equation}
Here $\otimes$ denotes the corresponding summand of the cut product.
The empty word is the coefficient $K$.

There are now two different ways to proceed. Cobar forms the associative
algebra $\Omega_{\mathrm{cut}}B_{\mathrm{cut}}A$, with its natural map
to $A$. The finite-word proof below establishes this map as a resolution.

To\phantomsection\label{ogd:def:dual-cut-tensor} exchange $A$ for a second algebra, one must also supply a coefficient
duality. Suppose a contravariant equivalence $\sharp$ is defined on
a monoidal subcategory containing $K$, the bar and its tensor powers.
Transport the cut tensor through $\sharp$, and denote the resulting
tensor by $\star_{\mathrm{cut},\sharp}$. Thus
\[
 F^\sharp\star_{\mathrm{cut},\sharp} G^\sharp
       =(F\star_{\mathrm{cut}}G)^\sharp.
\]

Transposing the coproduct then defines a multiplication on
\[
 A^!_\sharp=(B_{\mathrm{cut}}A)^\sharp,
 \qquad
 A^!_\sharp\star_{\mathrm{cut},\sharp} A^!_\sharp
       \xrightarrow{\ \Delta^\sharp\ } A^!_\sharp.
\]
Coassociativity supplies associativity. The bar counit supplies the unit
$K^\sharp\to A^!_\sharp$, and its coaugmentation supplies the augmentation
in the opposite direction. The domain of the duality and the transported
tensor are part of this construction.

The three arrows consequently have different targets:
\[
 A\longmapsto B_{\mathrm{cut}}A,\qquad
 \Omega_{\mathrm{cut}}B_{\mathrm{cut}}A\longrightarrow A,\qquad
 A\longmapsto A^!_\sharp.
\]
They give a coalgebra, a resolution and, under the stated duality
hypotheses, an exchanged algebra. A description by
$R\operatorname{Hom}_A(P,P)$ requires a specified $A$-module $P$
and a comparison with its endomorphisms. A vacuum expectation alone
does not specify that module or comparison. The four bar constructions
of Table~\ref{ogd:tab:bars} keep these coefficient questions separate.

\subsection{The first two cobar complexes}
\label{ogd:sec:first-cut-cobar}

Let\phantomsection\label{ogd:def:strong-cut-cobar-input} $C$ be a strong cut coalgebra with $C_\varnothing=k$.
The reduced cobar uses words of desuspended positive coefficients,
with $t=[-1]$ and $|tc|=|c|+1$. Its differential splits a block by
the reduced coproduct $\bar\Delta c=\sum c'\otimes c''$.
On a strict model the splitting of $tc$ is
\[
 tc\longmapsto
 \sum (-1)^{|c'|+1}(tc')(tc'').
\]
This notation abbreviates the coproduct on its localized coefficient
objects. The internal term is $-t(d_Cc)$.

At\phantomsection\label{ogd:def:binary-cut-localization} two positions there is one gap. Put
$L_1M=j_{1,*}j_1^!M$.
Strongness identifies the coefficient of length two with $L_1C_2$.
Regroup the desuspensions by
\[
 (tc_1)\otimes(tc_2)\longmapsto
       (-1)^{|c_1|}(c_1\otimes c_2)[-2].
\]

The splitting sign becomes $-1$, and the reduced cobar is
\begin{equation}\label{ogd:eq:two-cut-cobar}
 \overline{\Omega_{\mathrm{cut}}C}_2
  \simeq\operatorname{Tot}\bigl[C_2[-1]
                 \xrightarrow{-\mathrm{loc}}L_1C_2[-2]\bigr]
  \simeq R\Gamma_{\Delta_{12}}C_2[-1].
\end{equation}
The last equivalence is the localization triangle.
Internal differentials carry the signs of the displayed shifts.

At\phantomsection\label{ogd:def:ternary-cut-localizations} three positions there are two gaps. Write
$L_pM=j_{p,*}j_p^!M$ for $p=1,2$ and
$L_{12}M=j_{12,*}j_{12}^!M$, where $j_{12}$ includes $U_1\cap U_2$.
The consecutive partitions are
\[
 123,\qquad 1|23,\quad 12|3,\qquad 1|2|3.
\]

For three desuspended coefficients the regrouping sign is
$(-1)^{|c_2|}$. The two ways of adding the missing gap therefore
have opposite signs, giving
\begin{equation}\label{ogd:eq:three-cut-cobar}
 \begin{aligned}
 \overline{\Omega_{\mathrm{cut}}C}_3
 &\simeq\operatorname{Tot}\Bigl[
 C_3[-1]\xrightarrow{(-\mathrm{loc},-\mathrm{loc})}
 (L_1C_3\oplus L_2C_3)[-2]\\
 &\hspace{40mm}\xrightarrow{\ (u,v)\mapsto u|_{12}-v|_{12}\ }
 L_{12}C_3[-3]\Bigr].
 \end{aligned}
\end{equation}
The two composites from $C_3$ cancel. The last two terms compute
restriction to $U_1\cup U_2$, so the augmented complex is
$R\Gamma_{Z_3}C_3[-1]$.

Since $Z_3=\Delta_{13}$, its allowed support includes the outer collision
with the middle point distinct.
The general comparison in Theorem~\ref{ogd:thm:cut-local-cohomology}
uses the same finite cube of cuts and proves compatibility with every
monotone collision.

For\phantomsection\label{ogd:def:word-dualizing-coalgebra} example, put $J_n=p_{X^n}^!k=\omega_{X^n}[n\dim X]$ and
$J_\varnothing=k$. Localization to cut opens defines a strong cut
coalgebra, and transitivity of $!$-inverse image gives its collision maps.
Writing $i_{13}(x,y)=(x,y,x)$, the two preceding localization complexes give
\begin{equation}\label{ogd:eq:cut-dualizing}
 \overline{\Omega_{\mathrm{cut}}J}_2
    \simeq\Delta_{12,*}J_X[-1],\qquad
 \overline{\Omega_{\mathrm{cut}}J}_3
    \simeq i_{13,*}J_{X^2}[-1].
\end{equation}
The second object has a nonzero coefficient where $x_1=x_3\ne x_2$.
It is not a small-diagonal field coefficient. The general criterion in
Corollary~\ref{ogd:cor:cut-diagonal} will isolate this extra stratum.

The word tensor retains consecutive cuts and their obstructing equalities.
Recording the order of a collision adds a separate geometric choice.
For that construction, return to a smooth curve $X$ over $k$.
The next example takes $X=\mathbb A^1_k$.

\section{The binary collision problem}

Section~\ref{e1f:sec:tensor} required a tensor product that retains
collisions and the order of inputs. The word construction supplies
consecutive cuts. A full collision atlas also records which coincident
point came first. Its coefficients must be kept as complexes with their
operations, because their cohomology can lose singular products.

A simultaneous field product retains every Laurent coefficient.
To apply bar--cobar duality geometrically, those coefficients must also
survive restriction to collision strata. The two-point model shows why
a homotopy of regular products is too small for this purpose.

Two labels at distinct points have no collision order to record.
When their positions coincide, there are two possible orders.
Over a field \(k\) of characteristic zero, this gives a scheme made from two copies of
\(\operatorname{Spec}k[z_1,z_2]\), glued where \(z_1-z_2\) is invertible.
It has two diagonal branches and one common disjoint locus.
Keeping these branches requires a sheaf on this scheme, together with maps that describe subsequent collisions.

There is a small complex that displays the algebraic consequence.
Set\phantomsection\label{ogd:def:binary-chart-coefficients} \(R=k[z_1,z_2]\), \(u=z_1-z_2\), and \(L=R[u^{-1}]\).
Consider
\begin{equation}\label{ogd:eq:opening-binary}
 R\xi\xrightarrow{\ d\ }Le_{12}\oplus Le_{21},
 \qquad d\xi=e_{12}-e_{21},\qquad |\xi|=-1.
\end{equation}
The regular difference of the two products is a boundary.
The localized difference \((e_{12}-e_{21})/u\) is not a boundary, since its primitive would require \(\xi/u\).
The whole cohomology says which classes survive: \(H^{-1}=0\), and
\([(a,b)]\mapsto(a+b,[a])\) is an isomorphism \(H^0\cong L\oplus L/R\), under which
\((e_{12}-e_{21})/u^m\) goes to \((0,[u^{-m}])\ne0\) for every \(m\ge1\).
Thus a homotopy between regular products can coexist with a nonzero singular operation.
Section~\ref{ogd:sec:example} derives this complex from an associative algebra on the ordered collision schemes.

The construction must combine coefficients from disjoint subsets and retain their behavior along every collision.
For this purpose, ordinary tensor products of separate coordinate algebras are insufficient data.
Cohomology is also insufficient data.

The two-branch scheme is the product of the line with coordinate \(z_2\) and the two-branch line in \(u\).
As associative algebras, the gluing algebra \(D_1\) of \cite[Chapter~15]{S:R} models the derived functions on the two-branch line, and \(D_1\otimes_kk[z_2]\) those on the scheme.
The algebra \(D_1\) is not formal over \(k\), even for quasi-isomorphisms that ignore the action of the positions \cite[Theorem~15.5.4]{S:R}.
Setting \(z_2=0\) retracts \(D_1\otimes_kk[z_2]\) onto \(D_1\) by algebra maps.
Composed with these maps, an \(A_\infty\)-quasi-isomorphism from the cohomology algebra of \(D_1\otimes_kk[z_2]\) to \(D_1\otimes_kk[z_2]\) would give one from the cohomology algebra of \(D_1\) to \(D_1\).
Hence \(D_1\otimes_kk[z_2]\) is not formal either.

The\phantomsection\label{ogd:def:chamber-differential} two chambers of a collision in real time, $t>0$ and $t<0$, teach the same
lesson with an explicit algebra. For a cospan of unital algebras
$B_0\xrightarrow{\alpha_0}C\xleftarrow{\alpha_1}B_1$, such as $R\to K\leftarrow R$ with
$R=k[[u]]$ and $K=k((u))$, the complex $B_0\oplus B_1\to C$ with
$d(a,b)=\alpha_1(b)-\alpha_0(a)$ is a unital dg algebra under
\[
 (a,b)(a',b')=(aa',bb'),\qquad (a,b)c=\alpha_0(a)c,\qquad c(a,b)=c\,\alpha_1(b),\qquad cc'=0 .
\]
The two mixed products use different endpoints, and the Leibniz rule is the
identity $\alpha_1(bb')-\alpha_0(aa')=(\alpha_1(b)-\alpha_0(a))\alpha_1(b')+\alpha_0(a)(\alpha_1(b')-\alpha_0(a'))$.
For the collision cospan the cohomology is $R$ in degree zero and $K/R$ in degree
one: the pole class sits in degree one. Polynomial forms on an auxiliary interval
with endpoint values in $B_0$ and $B_1$ form a quasi-isomorphic model, but
integration over the interval is not an algebra map. With $f=x$ and $g=u^{-1}dx$,
$\int_0^1fg=1/2u$, while $f(0)\int_0^1g=0$.

\begin{question}\label{ogd:qn:chamber}
Do the ordered integrals $\int_{0\le x_1\le\cdots\le x_n\le1}g_1(x_1)\cdots g_n(x_n)$, with
signs fixed by suspension, form a strictly unital $A_\infty$ quasi-isomorphism from the path
model to the chamber algebra in every arity, compatibly with modules and a common
vacuum vector? The identity $\int_{y_{j-1}}^{y_j}f'=f(y_j)-f(y_{j-1})$ produces exactly the
adjacent products and the two endpoint actions that such a proof must match.
\end{question}

For polynomial paths over a characteristic-zero field, the answer
is given by Theorem~\ref{ogd:thm:path-chamber} and
Proposition~\ref{ogd:prop:path-modules}.
The suspended Taylor maps have coefficient $+1$ on tensors of
one-forms. The signs occur in the bar differential and its insertion
after a suspended prefix.
The module comparison retains a chosen common vacuum vector.
These statements concern polynomial paths with degree-zero
coefficient algebras and modules, as specified below.

\subsection{Polynomial paths and their contraction}
\label{ogd:sec:polynomial-paths}

For polynomial paths, the comparison can be constructed in every arity.
The endpoint conditions determine its first map.
The higher maps then follow from integrating one derivative at a time.
Work over the characteristic-zero field \(k\), and take tensor products over \(k\).
We keep the coefficient algebras \(B_0,B_1,C\) in degree zero.
They need not be commutative, and the homomorphisms
\(\alpha_i:B_i\to C\) need not be injective.

Write\phantomsection\label{ogd:def:named-chamber-algebra} \(A_{\mathrm{ch}}\) for the chamber algebra above.
Its polynomial path algebra \(\mathcal P\) has components
\begin{equation}\label{ogd:eq:path-carrier}
\begin{split}
 \mathcal P^0
 &=\{(a,b,f)\in B_0\oplus B_1\oplus C[x]:
       f(0)=\alpha_0(a),\ f(1)=\alpha_1(b)\},\\
 \mathcal P^1&=C[x]\,dx,
 \qquad d(a,b,f)=f'(x)\,dx .
\end{split}
\end{equation}
There are no other degrees.
The variable \(x\) is central.

Multiplication in degree zero is componentwise, the mixed products are
\(f(g\,dx)=fg\,dx\) and \((g\,dx)f=gf\,dx\), and the product of two
one-forms is zero.
The unit is \((1_{B_0},1_{B_1},1_C)\).
These rules make \(\mathcal P\) a unital differential graded algebra.
The derivative rule for a product in \(C[x]\) proves the Leibniz identity
without a commutativity assumption on \(C\).

All integrals below are algebraic integrals of polynomials.
For example, \(\int_0^x c y^r\,dy=c x^{r+1}/(r+1)\).
Characteristic zero supplies the denominators.
Define a projection \(p:\mathcal P\to A_{\mathrm{ch}}\) and an inclusion
\(i:A_{\mathrm{ch}}\to\mathcal P\) by
\begin{equation}\label{ogd:eq:path-linear-maps}
\begin{aligned}
 p(a,b,f)&=(a,b),
 &p(g\,dx)&=\int_0^1g(x)\,dx,\\
 i(a,b)&=(a,b,(1-x)\alpha_0(a)+x\alpha_1(b)),
 &i(c)&=c\,dx .
\end{aligned}
\end{equation}
The inclusion is linear and is not asserted to preserve multiplication.

\begin{samepage}
\begin{lemma}\label{ogd:lem:path-contraction}
The maps \(p,i\) are chain maps and \(pi=1\).
The degree-minus-one map
\[
 h(g\,dx)=
 \left(0,0,\int_0^x g(y)\,dy-x\int_0^1g(y)\,dy\right),
 \qquad h|_{\mathcal P^0}=0,
\]
satisfies
\[
 dh+hd=1-ip,\qquad h^2=0,\qquad hi=0,\qquad ph=0.
\]
In particular, \(p\) is a quasi-isomorphism of complexes.
These maps are natural in homomorphisms of the coefficient cospan.
\end{lemma}
\end{samepage}

\begin{proof}
The identity
\(\int_0^1f'(x)\,dx=\alpha_1(b)-\alpha_0(a)\)
shows that \(p\) is a chain map.
The derivative of the linear interpolation is the same difference,
so \(i\) is a chain map.
Integration of a constant gives \(pi=1\).

The polynomial in the definition of \(h\) vanishes at both endpoints,
so \(h\) has its stated target.
On a degree-zero path, \(hd\) replaces \(f(x)\) by
\[
 f(x)-f(0)-x(f(1)-f(0)).
\]
This is its difference from the linear interpolation.

On \(g(x)\,dx\), the map \(dh\) gives
\((g(x)-\int_0^1g)\,dx\).
These two identities prove the contraction.

The map \(h\) vanishes on degree zero and on constant one-forms.
Its output has zero endpoint components.
This proves the three side conditions.
Every formula commutes with the maps of a coefficient cospan.
\end{proof}

For \(R=k[[u]]\to K=k((u))\leftarrow R\), this contraction identifies
\(H^0(\mathcal P)\) with the diagonal copy of \(R\) and
\(H^1(\mathcal P)\) with \(K/R\).
The form \(u^{-m}dx\) represents \([u^{-m}]\ne0\) for \(m\geq1\).
The path variable is polynomial even though its coefficients lie in
the Laurent series field \(K\).

\subsection{Ordered integrals and higher comparison maps}
\label{ogd:sec:path-higher-maps}

The projection \(p\) does not preserve multiplication on cochains.
Its higher components retain the order of the factors.
For \(g_1,\ldots,g_n\in C[x]\), put
\begin{equation}\label{ogd:eq:path-ordered-integral}
 J_n(g_1,\ldots,g_n)=
 \int_{0\leq x_1\leq\cdots\leq x_n\leq1}
 g_1(x_1)\cdots g_n(x_n)\,dx_1\cdots dx_n .
\end{equation}
The coefficients are multiplied in the displayed order.

For \(g_j(x)=c_jx^{r_j}\), with \(r_j\geq0\), direct integration gives
\begin{equation}\label{ogd:eq:path-monomial-integral}
 J_n(c_1x^{r_1},\ldots,c_nx^{r_n})
 =
 \frac{c_1\cdots c_n}
 {(r_1+1)(r_1+r_2+2)\cdots(r_1+\cdots+r_n+n)}.
\end{equation}
In particular, constant coefficients give \(c_1\cdots c_n/n!\).
The integral is a finite \(k\)-linear combination of ordered coefficient
products.

To fix all signs, use the cochain suspension \(sA=A[1]\), with
\(|sa|=|a|-1\).
For any differential graded algebra \(A\), the tensor coalgebra
\[
 T^c(sA)=k\oplus\bigoplus_{n\geq1}(sA)^{\otimes n}
\]
has the deconcatenation coproduct.
The full underlying algebra is used here, so an augmentation of \(A\)
is not required.
Its degree-one bar coderivation \(b_A\) is specified by
\[
 b_1(sa)=-s(da),\qquad
 b_2(sa\otimes sb)=(-1)^{|a|}s(ab).
\]
When either map is inserted after a prefix, it acquires the Koszul sign
given by the degree of that suspended prefix.
These are the signs of~\eqref{ogd:eq:bar-signs}.

Define\phantomsection\label{ogd:def:path-taylor-maps} degree-zero Taylor maps
\[
 F_n:(s\mathcal P)^{\otimes n}\longrightarrow sA_{\mathrm{ch}}
\]
as follows.
The linear map is \(F_1(sv)=s\,p(v)\).
For \(n\geq2\), set
\begin{equation}\label{ogd:eq:path-taylor}
 F_n\bigl(s(g_1dx)\otimes\cdots\otimes s(g_ndx)\bigr)
   =sJ_n(g_1,\ldots,g_n),
\end{equation}
and set \(F_n\) to zero on every homogeneous tensor containing a
degree-zero path.
A one-form has suspended degree zero, as does the output in
\(sA_{\mathrm{ch}}^1\).
Thus the displayed maps have the stated degree.

Extend the Taylor maps to a coaugmentation-preserving coalgebra map
\(F:T^c(s\mathcal P)\to T^c(sA_{\mathrm{ch}})\) by consecutive subdivisions:
\[
 F(w_1\otimes\cdots\otimes w_n)
 =
 \sum_{r=1}^n\ \sum_{0=i_0<i_1<\cdots<i_r=n}
 F_{i_1-i_0}(w_1,\ldots,w_{i_1})\otimes\cdots\otimes
 F_{i_r-i_{r-1}}(w_{i_{r-1}+1},\ldots,w_n).
\]
Each word has finitely many subdivisions.

\begin{theorem}\label{ogd:thm:path-chamber}
For every cospan of unital associative degree-zero algebras over \(k\),
the map \(F\) satisfies
\[
 b_{A_{\mathrm{ch}}}F=F b_{\mathcal P}.
\]
It defines a natural strictly unital \(A_\infty\) quasi-isomorphism from
\(\mathcal P\) to \(A_{\mathrm{ch}}\), with linear term \(p\).
With the displayed suspension convention, every integral in
\eqref{ogd:eq:path-taylor} has coefficient \(+1\).
\end{theorem}

\begin{proof}
The difference \(b_{A_{\mathrm{ch}}}F-Fb_{\mathcal P}\) is a coderivation
along \(F\).
It is determined by its projection to \(sA_{\mathrm{ch}}\).
Indeed, if it vanishes on shorter input words, deconcatenation makes its
value on the next word primitive.
A primitive in the tensor coalgebra has length one, where the projection
detects it.
Thus induction reduces the proof to the length-one outputs.
The one-input identity is the chain identity for \(p\).

First take a word consisting entirely of one-forms.
Their differentials and pairwise products vanish.
Every Taylor output lies in \(A_{\mathrm{ch}}^1=C\), whose differential
and products with another degree-one output also vanish.
Both sides are zero.

Next suppose that exactly one input is a degree-zero path
\((a,b,f)\), in position \(j\).
All other suspended inputs have degree zero.
Write \(I_j\) for the ordered integral obtained by replacing this input
by \(f'(x_j)\,dx_j\).
When \(1<j<n\), write \(J_{j,-}\) for the integral in which
\(g_{j-1}f\) replaces the two adjacent factors on the left.
Write \(J_{j,+}\) for the integral with the adjacent product \(fg_{j+1}\)
on the right.

Integration in \(x_j\), between \(x_{j-1}\) and \(x_{j+1}\), gives
\[
 I_j=J_{j,+}-J_{j,-}.
\]
The internal differential contributes \(-sI_j\).
The two bar faces contribute \(-sJ_{j,-}\) and \(sJ_{j,+}\), respectively.
Their sum is zero.

On the other side, a nonzero Taylor output containing this input must be
a singleton.
Two blocks can isolate it only at an endpoint.

At the first position, the same integration has lower endpoint zero:
\[
 I_1=J_{1,+}-\alpha_0(a)J_{n-1}(g_2,\ldots,g_n).
\]
The surviving term is the left action of \((a,b)\) on the degree-one
Taylor output.
At the last position the upper endpoint is one, and
\[
 I_n=J_{n-1}(g_1,\ldots,g_{n-1})\alpha_1(b)-J_{n,-}.
\]
The surviving term is minus the right action on that Taylor output.
These are exactly the two signs in \(b_2\).
The argument includes two-input words.

A word with at least two degree-zero inputs contributes nothing when
\(n\geq3\).
In every term of \(F b_{\mathcal P}\), a differential or a multiplication
leaves at least one degree-zero input.
The remaining Taylor map has arity greater than one and therefore
vanishes.
A subdivision into two blocks also leaves such an input in a nonsingleton
block.
For \(n=2\), the remaining identity is that evaluation at the two
endpoints preserves the product of degree-zero paths.
This exhausts all homogeneous words and proves the coalgebra chain
identity.

The map \(F_1\) preserves the unit.
Every higher \(F_n\) vanishes when an input is the unit, since the unit
has degree zero.
The comparison is therefore strictly unital.
Lemma~\ref{ogd:lem:path-contraction} proves that its linear term is a
quasi-isomorphism.
Naturality follows from the ordered products and polynomial integrations.
\end{proof}

The two-form term already measures the failure of strict multiplicativity.
For a degree-zero path \(v=(a,b,f)\) and a one-form \(\omega=g\,dx\),
the endpoint identities are
\[
 p(v\omega)-p(v)p(\omega)=J_2(f',g),\qquad
 p(\omega v)-p(\omega)p(v)=-J_2(g,f').
\]
Taking \(v=(0,1,x)\) and \(\omega=u^{-1}dx\) gives the differences
\(1/(2u)\) and \(-1/(2u)\).
The path \(v\) is not closed.
The higher component accounts for its multiplicative defect on cochains.

\subsection{Modules and a common vacuum}
\label{ogd:sec:path-modules}

To\phantomsection\label{ogd:def:path-module-inputs} compare states, let \(M_i\) be a unital left \(B_i\)-module and \(N\)
a unital left \(C\)-module, all in degree zero.
Suppose that \(\beta_i:M_i\to N\) satisfies
\(\beta_i(bm)=\alpha_i(b)\beta_i(m)\).
The chamber module \(\mathcal M_{\mathrm{ch}}\) is
\[
 M_0\oplus M_1\xrightarrow{\ \beta_1-\beta_0\ }N
\]
in degrees zero and one.
Its degree-zero action on \(M_0\oplus M_1\) is componentwise.
The other actions are
\[
 (a,b)n=\alpha_0(a)n,\qquad
 c(m_0,m_1)=c\,\beta_1(m_1),\qquad cn=0
\]
for \(c\in A_{\mathrm{ch}}^1\) and \(n\in\mathcal M_{\mathrm{ch}}^1\).

The\phantomsection\label{ogd:def:module-path-projection} path module \(\mathcal P_M\) has degree-zero elements
\((m_0,m_1,v(x))\) with
\(v(0)=\beta_0(m_0)\) and \(v(1)=\beta_1(m_1)\).
Its degree-one part is \(N[x]\,dx\), and its differential is
\(v\mapsto v'\,dx\).
Multiplication of a coefficient polynomial by a module polynomial,
together with the endpoint actions, gives its \(\mathcal P\)-action.
Projection to endpoints in degree zero and integration in degree one
define \(p_M:\mathcal P_M\to\mathcal M_{\mathrm{ch}}\).

The module words form the comodules
\[
 T^c(s\mathcal P)\otimes s\mathcal P_M,\qquad
 T^c(sA_{\mathrm{ch}})\otimes s\mathcal M_{\mathrm{ch}}.
\]
The coaction cuts their algebra prefix.
Their bar differentials use the algebra faces, the internal module term
\(-s(dm)\), and the last action face \((-1)^{|a|}s(am)\).
The same suspended-prefix signs apply.

A module comparison over \(F\) is a map between these comodules covering
\(F\) and commuting with their differentials.
Its degree-zero Taylor maps have source
\((s\mathcal P)^{\otimes r}\otimes s\mathcal P_M\) and target
\(s\mathcal M_{\mathrm{ch}}\).
They extend to module words by consecutive subdivisions, with the module
input in the last block.

\begin{proposition}\label{ogd:prop:path-modules}
The map \(p_M\) extends naturally to a strictly unital \(A_\infty\)
module quasi-isomorphism over \(F\).
For \(r\geq1\), its suspended Taylor component on \(r\) algebra inputs
and one module input is
\begin{equation}\label{ogd:eq:path-module-taylor}
\begin{split}
 &F^M_r\bigl(s(g_1dx),\ldots,s(g_rdx);s(n\,dx)\bigr)\\
 &\qquad=
 s\int_{0\leq x_1\leq\cdots\leq x_{r+1}\leq1}
 g_1(x_1)\cdots g_r(x_r)n(x_{r+1})\,dx_1\cdots dx_{r+1}.
\end{split}
\end{equation}
It vanishes on a homogeneous tensor containing a degree-zero input,
and \(F^M_0=s p_M s^{-1}\).

If \(v_i\in M_i\) have a common image
\(\beta_0(v_0)=\beta_1(v_1)=v\), then the constant path
\((v_0,v_1,v)\) maps to the closed chamber vector \((v_0,v_1)\).
Every higher module component vanishes when its last input is this
constant path.
\end{proposition}

\begin{proof}
Encode a module as the upper-right entry of a triangular algebra:
\[
 \begin{pmatrix}B_i&M_i\\0&k\end{pmatrix}
 \longrightarrow
 \begin{pmatrix}C&N\\0&k\end{pmatrix}.
\]
The map is \(\alpha_i\) on the upper-left entry, \(\beta_i\) on the
upper-right entry and the identity on \(k\).
It is an algebra homomorphism by the compatibility assumed for
\(\beta_i\).

Apply Theorem~\ref{ogd:thm:path-chamber} to this cospan.
In a product of \((1,1)\) algebra entries and one \((1,2)\) module entry,
consider the \((1,2)\) component.
It vanishes whenever an algebra entry follows the module entry.
The resulting Taylor maps are precisely
\eqref{ogd:eq:path-module-taylor}.
The coalgebra chain identity in that matrix entry is the module identity.
Its last multiplication face has sign
\((-1)^{|a|}s(am)\), which is the same bar sign used above.

Replace \(B_i,C,\alpha_i\) in \eqref{ogd:eq:path-linear-maps} by
\(M_i,N,\beta_i\).
The same polynomial integrations and homotopy give a contraction of the
module complexes.
Thus \(p_M\) is a quasi-isomorphism.

The higher components vanish on a degree-zero unit input, which proves
strict unitality.
The matrix construction gives naturality.
The common-image condition makes the displayed constant path a closed
element, and its chamber image is closed for the same reason.
Every higher component has only degree-one inputs, proving the final
assertion.
\end{proof}


\subsection{Differential graded coefficients and marked bimodules}

The interval comparison also allows a differential on the coefficient
algebras. Its scalar iterated integrals are the interval comparison
of Bandiera--Sch\"atz\footnote{R. Bandiera and F. Sch\"atz,
\emph{How to discretize the differential forms on the interval},
\href{https://arxiv.org/abs/1607.03654}{arXiv:1607.03654},
Theorem 1.3, Remark 1.4 and \S3.1.}.
Let \(k\) now be a commutative \(\mathbb Q\)-algebra and
let \(B_0\xrightarrow{\alpha_0}C\xleftarrow{\alpha_1}B_1\)
be a cospan of unital differential graded associative \(k\)-algebras.
Place interval forms before coefficients. The polynomial path algebra is
\[
 \mathcal P_C=
 \{(a_0,a_1,f(u)+du\,g(u)):
       f(0)=\alpha_0a_0,\ f(1)=\alpha_1a_1\},
\]
with the total grading and differential
\[
 d(f+du\,g)=d_Cf+du(\partial_uf-d_Cg).
\]
An endpoint of a homogeneous path has the degree of its zero-form part;
the coefficient of its one-form part has degree one less.

Write \(T_C=B_0\oplus B_1\oplus eC\), where \(e\) has degree one.
For homogeneous endpoint pairs \(a=(a_0,a_1)\), define
\begin{align}
 d_Ta&=(d a_0,d a_1)+e(\alpha_1a_1-\alpha_0a_0),&
 d_T(ec)&=-e\,d_Cc,\notag\\
 ab&=(a_0b_0,a_1b_1),&
 a(ec)&=(-1)^{|a|}e(\alpha_0a_0c),\label{ogd:eq:dg-chamber}\\
 (ec)b&=e(c\alpha_1b_1),&
 (ec)(ec')&=0.\notag
\end{align}
The endpoint difference telescopes under multiplication, proving the
Leibniz identity. The other algebra identities follow from those of the
cospan. Thus \(T_C\) is a differential graded algebra.

Use suspension of degree minus one and the bar differential
\[
 b_1(sa)=-s(da),\qquad b_2(sa,sb)=(-1)^{|a|}s(ab).
\]
The prefix sign of either operation at position \(j\) is
\((-1)^{\sum_{i<j}(|a_i|-1)}\).

For marked coefficients, let
\(D_0\xrightarrow{\delta_0}E\xleftarrow{\delta_1}D_1\)
be another dg algebra cospan. Let \(M_i\) be a
\((B_i,D_i)\)-bimodule, \(N\) a \((C,E)\)-bimodule,
and let \(\beta_i:M_i\to N\) be cochain maps satisfying
\(\beta_i(bmd)=\alpha_i(b)\beta_i(m)\delta_i(d)\).
Define \(\mathcal P_M\) by polynomial paths with these endpoints,
and put \(T_M=M_0\oplus M_1\oplus eN\). Its differential and
endpoint actions are the off-diagonal components of
\eqref{ogd:eq:dg-chamber} for the corresponding triangular algebras.

\begin{theorem}\label{ogd:thm:dg-path-chamber}
The linear map
\[
 F_1(a_0,a_1,f+du\,g)
       =(a_0,a_1,e\!\int_0^1g(u)\,du)
\]
extends to a natural strictly unital \(A_\infty\) quasi-isomorphism
\(\mathcal P_C\to T_C\). For \(r>1\) its suspended Taylor map
vanishes if any input has interval-form degree zero. On one-form inputs,
\begin{equation}\label{ogd:eq:dg-iterated-integrals}
 \overline F_r\bigl(s(du\,g_1),\ldots,s(du\,g_r)\bigr)
 =s\left(e\int_{0<u_1<\cdots<u_r<1}
           g_1(u_1)\cdots g_r(u_r)\,du_1\cdots du_r\right).
\end{equation}
Every Taylor integral has coefficient \(+1\), and its coefficient
factors remain in their input order. The same integrals, with one
marked input retained in its position, define a strictly unital
\(A_\infty\) bimodule quasi-isomorphism \(\mathcal P_M\to T_M\)
over the comparisons for the two diagonal algebra cospans.
\end{theorem}

\begin{proof}
First check the internal coefficient differentials. For separated
homogeneous inputs \(\omega_i\otimes c_i\), moving coefficients
through the suspended interval factors contributes
\[
 (-1)^{\sum_{i<j}|c_i|(|\omega_j|-1)}.
\]
When every \(\omega_j\) is a one-form this sign is one.
The differential of their ordered coefficient product has prefix
\(\sum_{i<j}|c_i|\), which is exactly the suspended input prefix.
The minus sign in \(b_1\), together with the form-first differential,
therefore carries the coefficient differential to the differential of
the output edge. For a singleton zero-form the same assertion is the
chain identity for endpoint evaluation.

The remaining terms differentiate only interval coordinates or multiply
adjacent inputs. The proof of Theorem~\ref{ogd:thm:path-chamber}
applies with the displayed coefficient-transport sign. More explicitly,
with exactly one zero-form in position \(j\), integration of its
derivative gives its value at the next variable minus its value at the
previous variable. These are the two adjacent multiplication faces.
At the first or last position, the remaining value is its endpoint
action in \eqref{ogd:eq:dg-chamber}. No coefficient is commuted.
With only one-forms, all interval products and all products of output
edges vanish. With at least two zero-forms and arity greater than two,
each term still has a zero-form in a higher Taylor map and vanishes.
In arity two the last case is endpoint multiplicativity. This proves
the bar equation in every arity.

The linear section is
\[
 i(a_0,a_1,ec)=
 (a_0,a_1,(1-u)\alpha_0a_0+u\alpha_1a_1+du\,c).
\]
On one-forms set
\[
 h(du\,g)=\int_0^u g(v)\,dv-u\int_0^1g(v)\,dv,
\]
with zero endpoint components, and set \(h=0\) on zero-forms.
Direct differentiation gives \(dh+hd=1-iF_1\); the internal
differential terms cancel with the tensor signs. Thus \(F_1\)
is a quasi-isomorphism. The unit is a constant zero-form, so every
higher component containing it vanishes.

For the bimodule assertion, take cospans of upper triangular differential
graded algebras whose two diagonal entries are the left and right
coefficient algebras and whose off-diagonal entry is the bimodule.
Retain exactly one off-diagonal input in the preceding identities.
Matrix multiplication preserves its position and supplies the left
and right endpoint actions. Two off-diagonal entries multiply to zero.
The resulting identities are precisely the bimodule identities, and
the same linear contraction proves quasi-isomorphy.
\end{proof}

The rational hypothesis is necessary for polynomial paths. Over
\(\mathbb Z\), the class of \(u\,du\) is nonzero in
\(\mathbb Z[u]du/d\mathbb Z[u]\), since a primitive would have
quadratic coefficient \(1/2\). Twice this class is zero. Even the
identity coefficient cospan therefore cannot have the displayed
polynomial integration comparison over every characteristic-zero ring.

Finite holomorphic pole bounds behave additively in
\eqref{ogd:eq:dg-iterated-integrals}: integration introduces rational
scalars and no new holomorphic denominators. Each fixed Taylor map
can consequently be evaluated in a compatible coefficient tower on
which the indicated products are defined. The derived inverse limit
of its linear quasi-isomorphisms is again an equivalence. This is a
coefficientwise construction; it imposes no uniform pole bound or
convergence across the tower.

A vacuum vector is not a vacuum functional. If the first variation of a
Yangian state exchange $S(u)$ on the first-current states has vacuum
coefficient $2c^2/u^3$ on $e\otimes f$ at level $c$, then no regular scalar
functionals $\pi_\pm=\pi_0+\hbar\pi_{\pm,1}+\cdots$ on the two chambers satisfy
$\pi_-S=\pi_+$: at order $\hbar$ the difference $\pi_{+,1}-\pi_{-,1}$ would be regular and
equal to $2c^2/u^3$.

These comparisons retain the differential and the two endpoint actions
of the full chamber algebra.
They do not replace that algebra by its cohomology.
They compare a fixed coefficient cospan and its module cospans.
A geometric field comparison must additionally supply compatible maps
as the points and collision strata vary.

\subsection{Compact currents and their ordered envelope}

There is a second interval comparison in which the input is a Lie
algebra of compact currents. Let \(\mathfrak g\) be a differential
graded Lie algebra over \(\mathbb C\). For an open interval
\(J\subset\mathbb R\), put
\begin{equation}\label{ogd:eq:compact-current-envelope}
 \mathcal I_{\mathfrak g}(J)=
 C_*\bigl(\Omega_c^\bullet(J)\otimes\mathfrak g\bigr)
 =\bigoplus_{n\geq0}
   \operatorname{Sym}^n\bigl((\Omega_c^\bullet(J)
                                      \otimes\mathfrak g)[1]\bigr).
\end{equation}
All tensors are algebraic. Forms precede Lie coefficients, so
\([\omega x,\eta y]=(-1)^{|x||\eta|}\omega\eta[x,y]\).
The Chevalley differential has linear term \(sa\mapsto-s(da)\)
and binary term
\[
 d_{\mathrm{CE}}(sa\,sb)\big|_{\text{length }1}
                         =(-1)^{|a|}s[a,b].
\]
Extension by zero and multiplication for disjoint supports give the
prefactorization maps. On the target, let
\(\operatorname{Loc}U(\mathfrak g)\) be the locally constant
factorization algebra on the line with interval value
\(U(\mathfrak g)\); disjoint intervals multiply in their real order.

\begin{theorem}\label{ogd:thm:current-time-ordering}
There is a natural strict unital comparison on the interval basis
\[
 \Pi_J:\mathcal I_{\mathfrak g}(J)\longrightarrow U(\mathfrak g)
\]
which is a quasi-isomorphism and commutes with all disjoint-interval
products. It sends a symmetric monomial containing a time zero-form
to zero. On one-forms it is
\begin{equation}\label{ogd:eq:current-chambers}
 \begin{split}
 &\Pi_J\left(\prod_{i=1}^n s(f_i(t)\,dt\otimes x_i)\right)\\
 &\quad=\sum_{\sigma\in S_n}\epsilon(\sigma;x)
   \left(\int_{t_{\sigma(1)}<\cdots<t_{\sigma(n)}}
             \prod_i f_i(t_i)\,dt_1\cdots dt_n\right)
     x_{\sigma(1)}\cdots x_{\sigma(n)}.
 \end{split}
\end{equation}
Here each integral uses positive product Lebesgue measure, and
\(\epsilon(\sigma;x)\) is the Koszul sign for the degrees of the
\(x_i\). The empty input maps to one. Weiss extension gives an
equivalence of factorization algebras on the line.
\end{theorem}

\begin{proof}
The degree of \(s(dt\otimes x_i)\) is \(|x_i|\), so the
permutation sign makes the formula independent of a representative
of the symmetric monomial. Compact support makes all integrals finite.

For the chain identity, first take only time one-forms. Their real
differentials and pairwise brackets vanish. The coefficient differential
commutes with the ordered integral, with the suspended prefix signs,
and becomes the differential of the ordered product in
\(U(\mathfrak g)\). Next take exactly one zero-form \(f(t)x\).
Its real derivative is integrated between its two neighboring times.
The outer boundary contributes zero. An internal boundary belongs to
two chambers differing by the exchange of the adjacent inputs. Their
ordered products differ by the graded commutator, which in the
enveloping algebra is their Lie bracket. This cancels the corresponding
binary Chevalley term. For even \(x,y\), the two-input identity is
\[
 \int_{s<t}f'(s)g(t)\,ds\,dt\,xy
 +\int_{t<s}f'(s)g(t)\,ds\,dt\,yx
                  =\int fg\,dt\,[x,y].
\]
For homogeneous coefficients, the exchange contributes
\((-1)^{|x||y|}\) on both the chamber and Chevalley sides.
With at least two zero-forms, each linear or binary differential
leaves a zero-form and its image under \(\Pi_J\) vanishes.
These cases prove the chain identity in every word length.

For disjoint intervals, their real order fixes every chamber involving
points in different intervals. The integrals factor, and the resulting
product is the ordered enveloping product. Extension by zero preserves
each integral. This proves the strict prefactorization identities and
the unit assertion.

Filter the source by symmetric length and the target by PBW length.
In the associated graded, commute the \(x_i\) using the graded
commutative product. Its permutation sign cancels
\(\epsilon(\sigma;x)\), and all \(n!\) chambers together give
the integral over \(J^n\). Thus the associated-graded map is the
symmetric algebra of
\[
 (\Omega_c^\bullet(J)\otimes\mathfrak g)[1]
       \longrightarrow\mathfrak g,\qquad
 s(f\,dt\otimes x)\longmapsto\left(\int_J f\,dt\right)x.
\]
Compact de Rham integration is a quasi-isomorphism, with the bump
contraction used above. Symmetric powers preserve it over
\(\mathbb C\). A cycle in the cone has a largest word length.
Associated-graded acyclicity removes that part by a boundary; repeating
removes the cycle after finitely many steps. This proves quasi-isomorphy
without a boundedness assumption on \(\mathfrak g\).

For descent, a Weiss cover \((U_a)\) of an open \(U\) has
\((U_a^n)\) as a cover of \(U^n\). Every finite sum of algebraic
tensors has compact product support. A finite refinement by products
of sufficiently small opens, and a separate smooth partition of unity
in each variable, expresses it in the augmented \v{C}ech complex of this
cover. The same refinement contracts its positive \v{C}ech degrees.
Consequently descent holds in each form degree. At fixed length there
are only finitely many time-form degrees, and finite permutation
coinvariants are exact over \(\mathbb C\). This proves descent for
the associated graded with arbitrary dg Lie coefficients. The binary
differential lowers length, so the preceding finite induction proves
descent for \(\mathcal I_{\mathfrak g}\). The basis comparison
therefore extends to the asserted factorization equivalence.
\end{proof}

For identical unit-integral profiles, each chamber contributes
\(1/n!\). Their sum gives the PBW leading map with coefficient one.
The comparison is natural in dg Lie maps and can be applied levelwise
to coefficient towers. Applied to a sphere mode algebra, it constructs
the real-line envelope of those modes. A comparison with the local
fields on \(\mathbb R\times\mathbb C^d\) still requires their
supported field-operation maps.

The constructions below therefore keep the coefficient complexes, with all their operations.
We first construct the schemes and their collision maps.
Extraordinary inverse image then supplies the compatibility condition on differential modules.
Open direct image supplies the singular coefficients in the product.

Let\phantomsection\label{ogd:def:atlas-curve-and-grading} \(X\) be a smooth separated finite-type \(k\)-scheme of pure dimension one.
All complexes are cochain complexes.
The shift is \(V[m]^i=V^{i+m}\), and tensor products obey the Koszul rule.
Thus \(s=[1]\) lowers the degree of an element by one, while \(t=[-1]\) raises it by one.
All tensor products and functors between derived categories are derived.

The affine-line chart construction and its derived functions occur in the ravioli configuration spaces of Alfonsi--Kim--Young \ogdciteAKY*{\S4}.
The general nilpotent bar--cobar mechanism appears in Francis--Gaitsgory \ogdciteFG*{\S\S4.1--4.2}.
The proofs below construct the ordered collision diagram and use finite configuration arity directly.

\section{Ordered collision schemes}\label{ogd:sec:geometry}

For\phantomsection\label{ogd:def:finite-set-total-orders} a finite nonempty set \(I\), let \(\ogdOrd(I)\) be the set of its total orders.
Write \(\Delta_{ab}\subset X^I\) for the diagonal \(x_a=x_b\).
For \(\sigma,\rho\in\ogdOrd(I)\), define
\[
 V_{\sigma\rho}=X^I\setminus
 \bigcup_{\substack{\{a,b\}\subset I\\
       \sigma,\rho\text{ give opposite orders to }a,b}}\Delta_{ab}.
\]
Take one copy \(X^I_\sigma\) of \(X^I\) for each order.
Identify their opens \(V_{\sigma\rho}\) by the identity on position coordinates.

\begin{proposition}\label{ogd:prop:glue}
These identifications define a smooth, quasi-compact, quasi-separated scheme \(\mathscr R_I\).
The chart maps define a morphism \(q_I:\mathscr R_I\to X^I\) that is locally an isomorphism.
Over a geometric position tuple with equality blocks \(B\), the fibre of \(q_I\) has \(\prod_B |B|!\) points.
\end{proposition}
\begin{proof}
If two orders disagree on a pair, at least one disagrees with any third order on that pair.
Consequently,
\[
 V_{\sigma\rho}\cap V_{\sigma\eta}
   =V_{\sigma\rho}\cap V_{\rho\eta}.
\]
This is the triple-overlap condition for gluing schemes along open immersions.
The charts are smooth and the cover is finite.
Each intersection is an open subscheme of the noetherian scheme \(X^I\), so it is quasi-compact.
The resulting scheme is therefore quasi-separated.
The identities to \(X^I\) agree on overlaps and define \(q_I\).

At a fixed geometric tuple, two charts are identified precisely when their orders agree inside each equality block.
An arbitrary order on each block extends to an order on \(I\).
The distinct points of the fibre therefore correspond to these independent block orders.
\end{proof}

In particular, the triple diagonal has six branches for three labels.
These branches record transitive total orders.
They are not independent choices of three pairwise signs.
For \(|I|\geq2\), the scheme is generally nonseparated.

A\phantomsection\label{ogd:def:ordered-fibre-surjection-category} collision also needs an order within each collapsed fibre.
Let \(\ogdOSurj_+\) be a skeleton of the category of nonempty finite sets and surjections with totally ordered fibres.
For \(\alpha:I\twoheadrightarrow J\), expand each label of an order \(\sigma\) on \(J\) into its ordered fibre.
Denote the resulting order on \(I\) by \(e_\alpha(\sigma)\).
For composable surjections, order a composite fibre by concatenating the ordered fibres in the prescribed order.
Concatenation is associative and defines composition in \(\ogdOSurj_+\).

The affine-line model of this atlas is due to Alfonsi--Kim--Young
\cite[\S4.1]{S4aAlfonsiKimYoung}. The overlap equations give the following collision maps over the curve $X$.

\begin{proposition}\label{ogd:prop:collision}
The diagonal substitution \(x_i=y_{\alpha(i)}\) defines a morphism
\[
 \delta_\alpha:\mathscr R_J\longrightarrow \mathscr R_I.
\]
It sends the \(\sigma\)-chart into the \(e_\alpha(\sigma)\)-chart.
For \(I\xrightarrow\alpha J\xrightarrow\beta H\),
\[
 \delta_\alpha\delta_\beta=\delta_{\beta\alpha}.
\]
\end{proposition}
\begin{proof}
Let \(d_\alpha:X^J\to X^I\) be the usual diagonal map.
The expanded orders agree inside each fibre.
For distinct fibres, they disagree precisely when the corresponding labels of \(J\) disagree.
It follows that
\[
 d_\alpha^{-1}V_{e_\alpha(\sigma),e_\alpha(\rho)}=V_{\sigma\rho}.
\]
The chartwise diagonal maps therefore glue.
Both composite maps use the substitution \(x_i=z_{\beta\alpha(i)}\) and the same concatenated order.
They agree on every chart, which proves the identity.
\end{proof}

An ordered collision map need not be a closed immersion.
For example, collapse the ordered fibre \(\{1,3\}\), with \(1<3\), and leave \(2\) separate.
In the chart \(123\), its image is
\[
 \{x_1=x_3,\ x_1\ne x_2\}.
\]
The selected triple-collision branches lie in charts where \(1,3\) are consecutive.
The constructions below use extraordinary inverse image along the actual morphism \(\delta_\alpha\).

\section{Differential modules and disjointness convolution}\label{ogd:sec:category}

\subsection{The enhanced coefficient category}

On \(\mathscr R_I\), define differential modules and these operations by Zariski descent from the finite chart cover.
Every chart intersection is smooth and separated.
The descent data retain all transition equivalences and their homotopies.
Thus this definition does not impose separatedness on \(\mathscr R_I\).
The functors used below are local on these charts.

Define
\begin{equation}\label{ogd:eq:category}
 \ogdcC_+=\varprojlim_{\ogdOSurj_+}\ogdDMod(\mathscr R_I),
 \qquad
 \ogdcC=D(k)\times\ogdcC_+,
\end{equation}
where the arrow \(\alpha:I\twoheadrightarrow J\) acts by
\(\delta_\alpha^!:\ogdDMod(\mathscr R_I)\to\ogdDMod(\mathscr R_J)\).
An object \(F\) consists of \(F_\varnothing\in D(k)\), objects \(F_I\), and equivalences
\[
 \delta_\alpha^!F_I\simeq F_J
\]
compatible with all compositions and higher homotopies.
The empty component is isolated: no nonempty configuration has a collapse arrow to it.
Call \(F\) positive if \(F_\varnothing=0\).
The component functors are exact and jointly conservative, since equivalences in a limit category are detected componentwise.

Extraordinary inverse image preserves colimits in this setting.
On smooth charts it is a derived transfer-module tensor functor, with the fixed shifts and canonical-line conversions.
Tensor products preserve colimits, and the assertion descends through the chart cover.
Consequently, colimits in \(\ogdcC\) are computed componentwise.

\subsection{The convolution and its coherence}

For \(I=A\sqcup B\), restriction of total orders defines a morphism
\(q_{A,B}:\mathscr R_I\to \mathscr R_A\times \mathscr R_B\).
Its definition on charts respects gluing because a disagreement inside either subset was already a disagreement on \(I\).
Let \(U_{A,B}\subset \mathscr R_I\) impose \(x_a\ne x_b\) for all \(a\in A,b\in B\).
Restriction of orders gives an isomorphism
\begin{equation}\label{ogd:eq:disjoint-open}
 U_{A,B}\simeq(\mathscr R_A\times \mathscr R_B)_{\mathrm{disj}}.
\end{equation}
Indeed, concatenate orders on \(A\) and \(B\) to define the inverse on product charts.
Any other shuffle differs only on cross pairs, whose diagonals have been removed.
The inverses therefore agree on overlaps as scheme maps.

Write\phantomsection\label{ogd:def:atlas-disjoint-open-inclusion} \(j_{A,B}\) for the open inclusion into \(\mathscr R_I\).
Define
\begin{equation}\label{ogd:eq:convolution}
 (F\star G)_I=
 \bigoplus_{I=A\sqcup B}
 j_{A,B,*}(F_A\boxtimes G_B)|_{U_{A,B}}.
\end{equation}
The pairs \((A,B)\) are ordered, and either subset may be empty.
An empty subset contributes scalar tensor product with the corresponding object of \(D(k)\).
At the empty configuration the formula is \(F_\varnothing\otimes_kG_\varnothing\).

\begin{theorem}\label{ogd:thm:convolution}
Formula~\eqref{ogd:eq:convolution} defines a colimit-preserving symmetric monoidal product on \(\ogdcC\).
Its unit \(K\) has \(K_\varnothing=k\) and \(K_I=0\) for nonempty \(I\).
For an ordered partition \(P=(I_1,\ldots,I_r)\), every parenthesized product has the common form
\begin{equation}\label{ogd:eq:many-convolution}
 (F_1\star\cdots\star F_r)_I=
 \bigoplus_{I=I_1\sqcup\cdots\sqcup I_r}
 j_{P,*}(F_{1,I_1}\boxtimes\cdots\boxtimes F_{r,I_r})|_{U_P},
\end{equation}
where \(U_P\) removes every cross-block diagonal.
\end{theorem}
\begin{proof}
First check collision compatibility for \(\alpha:I\twoheadrightarrow J\).
If a fibre meets both \(A\) and \(B\), its collision lies outside \(U_{A,B}\).
Lemma~\ref{ogd:lem:basechange} makes the extraordinary pullback of this whole summand zero.
Otherwise there is a unique decomposition \(J=C\sqcup D\) with
\(A=\alpha^{-1}C\) and \(B=\alpha^{-1}D\).
The restricted ordered surjections \(\alpha_A,\alpha_B\) satisfy
\[
 \delta_\alpha^{-1}U_{A,B}=U_{C,D},\qquad
 q_{A,B}\delta_\alpha=(\delta_{\alpha_A}\times\delta_{\alpha_B})q_{C,D}.
\]
Base change and external-product compatibility identify the surviving summand with
\[
 j_{C,D,*}(\delta_{\alpha_A}^!F_A\boxtimes
                     \delta_{\alpha_B}^!G_B)|_{U_{C,D}}
 \simeq j_{C,D,*}(F_C\boxtimes G_D)|_{U_{C,D}}.
\]
Their sum gives \(\delta_\alpha^!(F\star G)_I\simeq(F\star G)_J\).
A partition survives two successive collisions exactly when it is saturated for the composite surjection.
The coherent base-change identifications then give all collision coherences.

For associativity, restrict an inner open direct image to the outer disjoint locus.
Lemma~\ref{ogd:lem:basechange} replaces it with direct image from the intersection of the two open conditions.
Open direct image also commutes with external products here.
Locally, the opens are complements of products of diagonal equations, so this statement is the identity for localization and external tensor product over \(k\).
Composing the open immersions imposes exactly the cross-block conditions of \(U_P\).
This proves~\eqref{ogd:eq:many-convolution} by induction.

Every associativity diagram has the same final partition and the same coefficient object on \(U_P\).
The coherent compositions of open direct image identify all its arrows.
This proves the pentagon and its higher coherence equations.
Permutation of the ordered parts, with the symmetry and Koszul signs of external products, gives the symmetric structure.
The same description proves its compatibility with associativity.
For the unit, only the decomposition \(\varnothing\sqcup I\), or \(I\sqcup\varnothing\), survives.
It gives the identity on \(F_I\), compatibly with every partition refinement.
Finally, the sums are finite at fixed \(I\), while localization and external tensor product preserve colimits.
Hence convolution preserves colimits separately.
\end{proof}

The order resides in \(\mathscr R_I\) and its collision maps.
Convolution itself is symmetric.
An associative algebra for this product can nevertheless have noncommutative multiplication.

\subsection{Projection onto consecutive decompositions}
\label{ogd:sec:atlas-cut-comparison}

The cut product on words and the atlas convolution retain different
decompositions. Their comparison can be written on each chart.
Let\phantomsection\label{ogd:def:increasing-chart-restriction} $s_n:X^n\hookrightarrow\mathscr R_{[n]}$ be the chart with order
$1<\cdots<n$. For an atlas coefficient system $F$, define
\[
 (\operatorname{res}_{\mathrm{ord}}F)_n=s_n^!F_{[n]},\qquad
 (\operatorname{res}_{\mathrm{ord}}F)_\varnothing=F_\varnothing.
\]
For a monotone surjection $\alpha:[n]\twoheadrightarrow[m]$,
$\delta_\alpha s_m=s_nd_\alpha$.
The collision identifications of $F$ therefore give a functor
\[
 \operatorname{res}_{\mathrm{ord}}:\ogdcC\longrightarrow\mathcal W_X.
\]

\begin{proposition}[The colax cut comparison]
\label{ogd:prop:atlas-cut-projection}
Projection onto decompositions
$[n]=[1,p]\sqcup[p+1,n]$ defines natural maps
\begin{equation}\label{ogd:eq:atlas-cut-projection}
 \pi_{F,G}:\operatorname{res}_{\mathrm{ord}}(F\star G)
 \longrightarrow
 (\operatorname{res}_{\mathrm{ord}}F)
       \star_{\mathrm{cut}}(\operatorname{res}_{\mathrm{ord}}G).
\end{equation}
Together with $\operatorname{res}_{\mathrm{ord}}K=K$, these maps make
$\operatorname{res}_{\mathrm{ord}}$ a colax monoidal functor.
In general it is not strong monoidal.
For an atlas coalgebra $C$, its induced cut coproduct is
\[
 \operatorname{res}_{\mathrm{ord}}C
 \xrightarrow{\operatorname{res}_{\mathrm{ord}}\Delta}
 \operatorname{res}_{\mathrm{ord}}(C\star C)
 \xrightarrow{\pi_{C,C}}
 (\operatorname{res}_{\mathrm{ord}}C)
       \star_{\mathrm{cut}}(\operatorname{res}_{\mathrm{ord}}C).
\]
\end{proposition}

\begin{proof}
Restrict each summand of \eqref{ogd:eq:convolution} to the increasing
chart. Lemma~\ref{ogd:lem:basechange} identifies it with the open
direct image on $X^n$ for the same decomposition $A\sqcup B$.
For a prefix $A=[1,p]$ and suffix $B=[p+1,n]$, this is precisely
the $p$th summand of \eqref{ogd:eq:cut-product}.
The finite direct-sum projection gives \eqref{ogd:eq:atlas-cut-projection}.

Under a monotone collision, a summand vanishes if one fiber meets
both $A$ and $B$. Otherwise the decomposition descends to the quotient.
A subset is saturated if it is a union of fibers.
A saturated subset is a prefix before collapse exactly when its
quotient is a prefix. Thus the projection commutes with the collision
maps and their coherent base-change identifications.

For three or more factors, either parenthesization retains exactly
the decompositions into final consecutive blocks.
All retained coefficient maps are the same restriction and open
extension on the common disjoint locus.
The coherent functoriality of these operations proves the colax
associativity equations and their higher compatibilities.
An empty factor leaves the unique unit summand, so the unit equations
hold as well. The displayed composite then transports every
coalgebra coherence.

To see that the comparison need not be an equivalence, take the
positive dualizing atlas system $J^+$.
Its empty coefficient is zero, and its nonempty coefficient is
$p_{\mathscr R_I}^!k$.
Assume $X$ has two distinct geometric points.
On $U=X^2\setminus\Delta$, the comparison for $F=G=J^+$ is
\[
 J_U\oplus J_U\longrightarrow J_U.
\]

The source has the decompositions $(\{1\},\{2\})$ and
$(\{2\},\{1\})$. Only the first is consecutive in the increasing
word, so the map is projection onto that summand.
Its kernel is the nonzero second copy of $J_U$.
\end{proof}

\input{manuscript/constructions/ordered-cut-operations.tex}

Coalgebra transport thus uses the specified projection.
It does not identify the symmetric atlas convolution with the cut
tensor. The following finite-word argument applies to each product
with its own coefficient category and unit.
\section{Finite words and associative duality}\label{ogd:sec:duality}

\subsection{The arity filtration}

A word made from positive factors assigns at least one label to each factor.
This observation controls every infinite expression used below.
For \(q\geq0\), define
\[
 \ogdcC_{\geq q}=\{F:F_I=0\text{ for }|I|<q\}.
\]
Thus \(\ogdcC_{\geq1}=\ogdcC_+\), with the empty component understood as zero.

\begin{lemma}\label{ogd:lem:finite}
Convolution satisfies
\[
 \ogdcC_{\geq p}\star\ogdcC_{\geq q}\subseteq\ogdcC_{\geq p+q}.
\]
If \(F_1,\ldots,F_r\) are positive and \(M\in\ogdcC\) is arbitrary, then
\begin{equation}\label{ogd:eq:mixed-bound}
 (F_1\star\cdots\star F_r\star M)_I=0\qquad(r>|I|).
\end{equation}
If \(W_r\in\ogdcC_{\geq r}\), the canonical map \(\bigoplus_rW_r\to\prod_rW_r\) is an equivalence.
Its component at \(I\) is the finite biproduct \(\bigoplus_{r=0}^{|I|}(W_r)_I\).
\end{lemma}
\begin{proof}
In a nonzero summand of~\eqref{ogd:eq:convolution}, the two parts have sizes at least \(p,q\).
The first assertion follows by addition of these sizes.
For~\eqref{ogd:eq:mixed-bound}, each of the first \(r\) parts is nonempty.
The part carrying \(M\) may be empty, so the bound is \(r\), without an additional unit of arity.

Set \(W_I=\bigoplus_{r\leq|I|}(W_r)_I\).
For a collision \(I\twoheadrightarrow J\), exactness gives
\[
 \delta_\alpha^!W_I\simeq\bigoplus_{r\leq|I|}(W_r)_J
  =\bigoplus_{r\leq|J|}(W_r)_J=W_J.
\]
The vanished terms and the component equivalences give coherent descent.
Colimit continuity identifies \(W\) with the coproduct.
A family of maps \(T\to W_r\) gives, at each \(I\), a map to a finite product.
These maps commute with collisions because they do so after every projection to \((W_r)_J\).
The same construction applies to mapping spaces and their coherent homotopies.
It proves the product universal property for \(W\).
\end{proof}

This lemma asserts no continuity for arbitrary products under extraordinary inverse image.
It constructs the particular products of increasingly positive words that occur in bar and cobar.
Any compatible finite expression on each configuration therefore defines a global map.

Restriction to configurations of size at most \(n\) defines a category \(\ogdcC_{\leq n}\).
The same partition formula makes restriction strong symmetric monoidal.
There are equivalences
\[
 \ogdcC\simeq\varprojlim_n\ogdcC_{\leq n},
 \qquad (\ogdcC_{\leq n,+})^{\star(n+1)}=0.
\]
The first follows because every finite set and every finite diagram occur in a finite truncation.
The second is exactly the nonempty-part bound.

\subsection{Bar and cobar differentials}

Let\phantomsection\label{ogd:def:augmented-atlas-algebra} \(A\) be an augmented associative algebra in \((\ogdcC,\star,K)\).
We require that its augmentation \(A\to K\) be an equivalence at the empty component.
The augmentation ideal \(a=\operatorname{fib}(A\to K)\) is positive, and the unit gives \(A\simeq K\oplus a\).
Define the normalized bar by
\begin{equation}\label{ogd:eq:bar}
 BA=K\oplus\bigoplus_{r\geq1}(sa)^{\star r}.
\end{equation}
The coproduct cuts a word into an initial and a final segment, including the two empty segments.
On a strict differential graded representative the differential is the coderivation determined by
\begin{equation}\label{ogd:eq:bar-signs}
 b_1(sa)=-s(da),\qquad b_2(sa,sb)=(-1)^{|a|}s(ab).
\end{equation}
For a word \([a_1|\cdots|a_r]\), its multiplication term is
\[
 \sum_{i=1}^{r-1}(-1)^{\sum_{j<i}(|a_j|-1)+|a_i|}
 [a_1|\cdots|a_ia_{i+1}|\cdots|a_r].
\]
Thus the bar retains ordered words, with their geometric convolution coefficients.

Let\phantomsection\label{ogd:def:reduced-atlas-coalgebra} \(C\) be a coaugmented counital coassociative coalgebra whose coaugmentation is an equivalence at the empty component.
Write \(c=\operatorname{cofib}(K\to C)\) and let \(\bar\Delta:c\to c\star c\) be its reduced coproduct.
Define
\begin{equation}\label{ogd:eq:cobar}
 \Omega C=K\oplus\bigoplus_{r\geq1}(tc)^{\star r}
\end{equation}
with concatenation product.
Its differential is the derivation determined by
\begin{equation}\label{ogd:eq:cobar-signs}
 d(tc)=-t(d_Cc)+
 \sum_{\bar\Delta c=\sum c'\otimes c''}
          (-1)^{|c'|+1}(tc')(tc'').
\end{equation}
Sweedler notation here abbreviates a morphism in \(\ogdcC\), including its localized coefficients.
It does not assert an elementwise decomposition of every differential module.

\begin{lemma}\label{ogd:lem:squares}
The differentials in~\eqref{ogd:eq:bar-signs} and~\eqref{ogd:eq:cobar-signs} square to zero.
They preserve the coalgebra and algebra structures, respectively.
\end{lemma}
\begin{proof}
For the bar, two internal differentials give zero.
Contractions at disjoint adjacent pairs cancel because both contractions have degree one.
For an overlapping triple, the remaining coefficient is
\[
 (-1)^{|b|}s\bigl((ab)c-a(bc)\bigr)=0.
\]
The mixed internal and multiplication terms cancel by the Leibniz identity for \(d_A\).
These cases exhaust \(b^2\).
The differential is a coderivation by its construction on the tensor coalgebra.

For cobar, the internal square is zero and the mixed terms cancel because \(\bar\Delta\) is a chain map.
The two orders of splitting a factor into three parts have opposite suspension signs.
Coassociativity identifies their underlying maps, so they cancel.
Splittings in different factors cancel by the derivation sign.
The differential extends as a derivation, which proves product compatibility.
\end{proof}

\subsection{Coherent functors and normalized words}
\label{ogd:sec:enhanced-bar}

A\phantomsection\label{ogd:def:coherent-coalgebra-category} comparison of algebra structures must also compare their higher homotopies.
We use associative algebra objects in the stable monoidal enhancement, and define
\[
 \ogdCoalg_{E_1}^{\mathrm{coaug}}(\ogdcC,\star)
 =\bigl(\ogdAlg_{E_1}^{\mathrm{aug}}
       (\ogdcC^{\mathrm{op}},\star)\bigr)^{\mathrm{op}}.
\]
Thus coalgebra morphisms preserve the counit, coproduct and all their coherences.
This definition imposes no choice of a strict coalgebra model.

Each differential-module category in~\eqref{ogd:eq:category} is presentable.
The transfer-module functors are accessible and preserve colimits.
Limits of presentable categories along such functors remain presentable, with the same underlying limit of categories
\cite{wdesc:lurie-htt}*{Proposition~5.5.3.13}.
Consequently, $\ogdcC$ admits geometric realizations and cosimplicial totalizations.
Colimits compute componentwise, and exactness gives the same assertion for finite limits.
We use no componentwise formula for arbitrary products.

\begin{proposition}\label{ogd:prop:enhanced-bar}
Geometric bar and cosimplicial cobar define an adjunction
\[
 B:\ogdAlg_{E_1}^{\mathrm{aug}}(\ogdcC,\star)
 \ \rightleftarrows\
 \ogdCoalg_{E_1}^{\mathrm{coaug}}(\ogdcC,\star):\Omega,
 \qquad B\dashv\Omega.
\]
Its derived mapping-space equivalence is
\begin{equation}\label{ogd:eq:forward-adjunction}
 \operatorname{Map}_{\ogdCoalg}(BA,C)
 \simeq\operatorname{Map}_{\ogdAlg}(A,\Omega C),
\end{equation}
natural in $A$ and $C$.
For positive reduced parts, these functors have the normalized expressions
\eqref{ogd:eq:bar} and~\eqref{ogd:eq:cobar}.
Their unit and counit are coherent maps
\[
 u_A:A\longrightarrow\Omega BA,
 \qquad v_C:B\Omega C\longrightarrow C.
\]
They satisfy the triangle identities with all higher coherences.
\end{proposition}
\begin{proof}
Apply the associative bar theorem of Lurie
\cite{LurieHA}*{Theorem~5.2.2.17 and Remark~5.2.2.19}.
For the augmented version, use objects $K\to T\to K$ whose composite is the identity.
The theorem identifies both mapping spaces in~\eqref{ogd:eq:forward-adjunction}
with the same space of compatible algebra structures on arrows.
Representability therefore supplies the adjunction, including its unit, counit and coherent triangle identities.
Its hypotheses require the indicated realizations and totalizations.
They do not require tensor product to preserve both
\cite{LurieHA}*{Warning~5.2.2.16}.

The underlying bar is the realization with simplicial terms $A^{\star q}$.
Faces multiply adjacent factors or apply an augmentation.
Degeneracies insert the unit.
Splitting $A=K\oplus a$ removes the degenerate summands and leaves $a^{\star q}[q]$ in normalized degree $q$.
The skeletal filtration and its attaching maps retain the complete simplicial structure
\cite{LurieHA}*{Theorem~1.2.4.1 and Remarks~1.2.4.2--1.2.4.3}.
They give~\eqref{ogd:eq:bar}, including its differential.

The coalgebra structure also agrees with the normalized formula.
The enhanced coproduct is the map
\[
 K\star_A K\simeq K\star_A A\star_A K
 \longrightarrow K\star_A K\star_A K.
\]

The normalized two-sided-bar comparison cuts a word at every position:
\[
 [a_1|\cdots|a_r]\longmapsto
 \sum_{i=0}^r[a_1|\cdots|a_i]\otimes[a_{i+1}|\cdots|a_r].
\]
This is the Alexander--Whitney comparison, with both empty cuts retained.
Its iterates sum over the same weakly increasing lists of cut positions.
Every parenthesization therefore retains the same coefficient maps and simplicial face identities.
This identifies the strict normalized bar with the enhanced coalgebra, including its coassociative coherences.

Dually, the underlying cobar totalizes the cosimplicial terms $C^{\star q}$.
Its cofaces use the coproduct and coaugmentation, and its codegeneracies use the counit.
Normalization leaves $c^{\star q}[-q]$.
At configuration $I$, this object vanishes for $q>|I|$.
The normalized totalization tower therefore stabilizes after $|I|$ steps.

Each preceding step uses finite limits.
Exact collision pullback preserves those steps and their comparison maps.

The resulting finite totalizations form a compatible family $H_I$.
They compute the totalization in $\ogdcC$, not merely its proposed components.
Indeed, maps into a compatible family form a limit of component mapping spaces.
For every $T\in\ogdcC$, commuting this limit with the cosimplicial limit gives
\[
 \operatorname{Map}_{\ogdcC}(T,H)
 \simeq\varprojlim_{[q]\in\Delta}
        \operatorname{Map}_{\ogdcC}(T,C^{\star q}).
\]
This is the required universal property.

Concatenation and its coherent operations follow from the same cosimplicial construction.
Its normalized differential is~\eqref{ogd:eq:cobar-signs}.

Every finite diagram of collision maps admits a common finite arity bound.
Hence these identifications retain the functors, their natural transformations and their higher coherences.
They bound normalized word length, not the cohomological degrees of a coefficient resolution.
\end{proof}

On strict derived presentations, normalization gives exactly the signs above.
For an \(A_\infty\) presentation, replace \(b_2\) by the degree-one Taylor maps
\(b_r:(sa)^{\star r}\to sa\).
The full equation is
\[
 \sum_{u+v+w=n}b_{u+1+w}
 (1^{\star u}\star b_v\star1^{\star w})=0,
\]
with tensor-composition signs.
A degree-zero coalgebra morphism between tensor bars has Taylor maps \(F_r\), extends by consecutive subdivisions, and obeys \(b'F=Fb\).
These are the normalized coordinates of the coherent construction in
Proposition~\ref{ogd:prop:enhanced-bar}.
All normalized higher terms remain finite at each arity by Lemma~\ref{ogd:lem:finite}.

\subsection{The comparison maps and their proof}

Write\phantomsection\label{ogd:def:twisting-cochain-data} \(\mu=\mu_A:a\star a\to a\) for the multiplication on the augmentation ideal.
A degree-one map \(\tau:c\to a\) is a twisting cochain if
\begin{equation}\label{ogd:eq:twist}
 d_A\tau+\tau d_C+\mu(\tau\star\tau)\bar\Delta=0.
\end{equation}
The tensor sign in \(\tau\star\tau\) contributes \((-1)^{|c'|}\) on \(c'\otimes c''\).
Thus~\eqref{ogd:eq:twist} is exactly the chain-map condition for \(\Omega C\to A\).
Iterated reduced coproduct gives the equivalent coalgebra map \(C\to BA\).
Only finitely many terms survive at each configuration.

These formulas identify strict twisting cochains with strict comparison maps.
Proposition~\ref{ogd:prop:enhanced-bar} supplies the enhanced functors independently of that identification.
The reverse adjunction will follow from their equivalence below.

The strict comparison
\[
 \epsilon_A:\Omega BA\longrightarrow A
\]
sends \(t(sa)\) to \(a\), sends a generator represented by a longer bar word to zero, and extends multiplicatively.
The strict comparison \(\eta_C:C\to B\Omega C\) is the identity on \(K\), and on \(c\) it is
\begin{equation}\label{ogd:eq:unit}
 \eta_C(c)=\sum_{r\geq1}(s\tau)^{\star r}\bar\Delta^{(r-1)}c,
 \qquad \tau(c)=tc.
\end{equation}

The degree-zero map \(s\tau\) contributes no further global sign.
For two outputs \(c',c''\), the splitting term of \(b_1s\tau(c)\) has coefficient \(-(-1)^{|c'|+1}\).
The multiplication term of the length-two part has coefficient \((-1)^{|c'|+1}\).
They cancel, and the same cancellation at every adjacent pair proves the chain identity in all lengths.
Deconcatenation matches iterated coproduct, so \(\eta_C\) is a coalgebra map.
The corresponding twisting equation proves that \(\epsilon_A\) is an algebra map.

\begin{theorem}\label{ogd:thm:duality}
The functors \(B\) and \(\Omega\) induce inverse equivalences
\[
 \ogdAlg^{\mathrm{aug},\,A_\varnothing=k}_{E_1}(\ogdcC,\star)
 \simeq
 \ogdCoalg^{\mathrm{coaug},\,C_\varnothing=k}_{E_1}(\ogdcC,\star).
\]
The coalgebras are homotopy-coherent counital coassociative coalgebras.
No additional global elementwise conilpotence assumption is needed.
\end{theorem}
\begin{proof}
We apply the gap calculation to the coherent maps $u_A$ and $v_C$ of
Proposition~\ref{ogd:prop:enhanced-bar}.
For an increasing filtration, the tensor filtration in degree $n$ combines terms whose indices sum to at most $n$.
For a decreasing filtration, it combines terms whose indices sum to at least $n$.
These constructions give monoidal categories of filtered objects.
The enhanced adjunction applies in those categories as well.

Give $A$ the increasing filtration $F_0A=K$ and $F_nA=A$ for $n\geq1$.
Its associated graded has reduced part $a$ in weight one and zero reduced multiplication.
The induced filtration on $\Omega BA$ counts original $a$-factors.

Give $C$ the decreasing filtration $F^0C=C$, $F^1C=c$ and $F^2C=0$.
Its associated graded has primitive reduced part $c$ in weight one.
The induced filtration on $B\Omega C$ counts original $c$-factors.
The maps $u_A$ and $v_C$ preserve these filtrations by naturality of the enhanced construction.

Flatten a term of \(\Omega BA\) into its \(n\) original \(a\)-factors.
Each of the \(n-1\) intervening gaps is either inside a bar block or between two cobar factors.
Filter increasingly by \(n\).
The coefficient differential and block splitting preserve \(n\), while multiplication of adjacent \(a\)-factors decreases it.
On the associated graded, each gap contributes a two-term identity complex.
More explicitly, the gap states form the cube \(\{0,1\}^{n-1}\), and the differential changes an inside gap to an outside gap.
Every edge map is an identity on the common shifted coefficient \(a^{\star n}\), with its suspension sign.
The square-zero identity gives opposite signs around each square.
Starting at one vertex fixes compatible signs identifying this cube with the tensor product of \(n-1\) identity complexes.
That product is contractible for \(n>1\).
For \(n=1\), its surviving term is \(a\), and the empty term is \(K\).

In weights zero and one, normalized bar and cobar are the identity and the inverse shifts $s,t$.
Thus the associated graded unit $u_A$ is the identity on the surviving terms.
In higher weights it maps zero to the contractible gap complex.
At a configuration $I$, all terms with $n>|I|$ vanish.
The filtration is finite, so its cofiber sequences prove that $(u_A)_I$ is an equivalence.

For \(B\Omega C\), flatten to the original \(c\)-factors and use the decreasing filtration by their number \(n\).
The coefficient differential and bar multiplication preserve \(n\).
The cobar splitting increases it.
The associated graded again has the gap cube, now changing an outside bar cut into an inside product gap.
It is contractible for \(n>1\), with the singleton term \(c\) surviving.

The associated graded counit $v_C$ is the identity on weights zero and one.
In higher weights it maps the contractible gap complex to zero.
Normalization and associated graded commute here because each evaluated filtration is finite.
After evaluation at $I$, the filtration ends at $|I|$.
Hence $(v_C)_I$ is an equivalence.

The component functors are jointly conservative.
They therefore prove both equivalences in \(\ogdcC\), hence in the algebra and coalgebra categories.
For coherent presentations, the additional higher products lower the first filtration, and higher coproducts raise the second.
Their associated graded comparisons are unchanged.
The same finite argument applies, and the adjunction supplies the triangle coherences.
\end{proof}

\input{manuscript/constructions/benchmark-chiral-koszul.tex}

The equivalence pairs algebras with coalgebras: \(\Omega BA\) resolves \(A\) and is not a dual of it, while for holonomic coefficients Verdier duality turns \(BA\) into an algebra dual to \(A\) (Theorem~\ref{ogd:thm:polarity}).

Inverting $u$ and $v$ gives the adjunction $\Omega\dashv B$ as well.
Its enhanced counit and unit are $u_A^{-1}$ and $v_C^{-1}$.
For positive reduced parts, its mapping-space equivalence is
\begin{equation}\label{ogd:eq:reverse-adjunction}
 \operatorname{Map}_{\ogdAlg}(\Omega C,A)
 \simeq\operatorname{Map}_{\ogdCoalg}(C,BA),
 \qquad f\longmapsto Bf\circ v_C^{-1}.
\end{equation}
Its inverse sends $g$ to $u_A^{-1}\circ\Omega g$.
These maps are natural in both objects and retain the triangle coherences of the equivalence.

We\phantomsection\label{ogd:def:strict-cobar-comparisons} identify the strict comparison formulas by an explicit structured comparison.
Let $\Omega_{\mathrm s}$ denote the displayed strict cobar on strict derived presentations.
Write $\epsilon^{\mathrm s}$ and $\eta^{\mathrm s}$ for its displayed comparison maps.

The decreasing gap filtration also applies directly to
$\eta^{\mathrm s}_C:C\to B\Omega_{\mathrm s}C$.
Its singleton component is the identity, and its other components increase the number of original factors.
The associated graded higher-weight targets are the same contractible gap cubes.
Thus $\eta^{\mathrm s}_C$ is an equivalence of coalgebras.
Here $\eta^{\mathrm s}_C$ is already a strict coalgebra map.

The two strict triangle identities follow by consecutive subdivisions:
\[
 B\epsilon^{\mathrm s}_A\circ\eta^{\mathrm s}_{BA}=1,
 \qquad
 \epsilon^{\mathrm s}_{\Omega_{\mathrm s}C}
       \circ\Omega_{\mathrm s}\eta^{\mathrm s}_C=1.
\]
In the first identity, every block of length greater than one vanishes under $\epsilon^{\mathrm s}_A$.
Only the completely cut subdivision survives.
In the second, only the one-block subdivision survives on a cobar generator.
Multiplicativity gives the identity on every cobar word.

Define\phantomsection\label{ogd:def:strict-enhanced-cobar-equivalence} the natural equivalence
\[
 \theta_C=(\Omega\eta^{\mathrm s}_C)^{-1}
                 \circ u_{\Omega_{\mathrm s}C}:
 \Omega_{\mathrm s}C\longrightarrow\Omega C.
\]
Both factors are equivalences by the preceding result and Theorem~\ref{ogd:thm:duality}.
Naturality of $u$, together with the first strict triangle, gives
$\theta_{BA}=u_A\circ\epsilon^{\mathrm s}_A$.
Naturality of $v$ and its enhanced triangle identity give
\[
 v_C\circ B\theta_C\circ\eta^{\mathrm s}_C=1.
\]
Consequently $\theta$ identifies the strict formulas with the inverse enhanced unit and counit.

The enhanced functors remain defined for all coherent coalgebras.
This argument requires strict models only where it displays strict formulas.

The proof uses precisely the local conilpotence of reduced configuration words.
It also explains why unnormalized unit factors must first be removed: their number has no arity bound.
The ordinary degeneracy normalization removes those factors before the finite argument applies.

\section{The cobar algebra of the dualizing coalgebra}\label{ogd:sec:example}

The\phantomsection\label{ogd:def:atlas-dualizing-coalgebra} geometry provides a coalgebra before any algebra is chosen.
Put \(J_I=p_I^!k=\omega_{\mathscr R_I}[|I|]\) and \(J_\varnothing=k\).
Transitivity gives \(\delta_\alpha^!J_I=J_J\), so \(J\in\ogdcC\).
For every ordered decomposition \(P\), localization gives
\begin{equation}\label{ogd:eq:J-coprod}
 J_I\longrightarrow j_{P,*}J_{U_P}
   =j_{P,*}(J_{I_1}\boxtimes\cdots\boxtimes J_{I_r})|_{U_P}.
\end{equation}
The sum of these maps is the iterated coproduct.
Successive localizations impose the same final disjointness conditions, which proves coassociativity.
Permutation of the blocks proves cocommutativity with the symmetry of differential modules.
The counit \(J\to K\) and coaugmentation \(K\to J\) are the identity at the empty component and zero elsewhere.
The counit identities select the decompositions \((I,\varnothing)\) and \((\varnothing,I)\), whose localization maps are identities.
Thus \(J\) is a connected coaugmented counital \(\star\)-coalgebra.

Define the associative chiral algebra
\begin{equation}\label{ogd:eq:U}
 \ogdUU=\Omega J.
\end{equation}
Theorem~\ref{ogd:thm:duality} gives a coalgebra equivalence \(J\xrightarrow{\sim}B\ogdUU\).
The explicit differential of \(\ogdUU\) permits direct computation of its singular multiplication.

\subsection{Ordered partitions and all signs}

For \(|I|=n>0\), the underlying complex is
\begin{equation}\label{ogd:eq:U-partitions}
 \ogdUU_I=
 \bigoplus_{r=1}^{n}\ \bigoplus_{\substack{I=B_1\sqcup\cdots\sqcup B_r\\B_i\ne\varnothing}}
 j_{P,*}\omega_{U_P}[n-r],
 \qquad P=(B_1|\cdots|B_r).
\end{equation}
The term with \(r\) blocks has degree \(r-n\).
Its differential splits one block into two and localizes to the refined disjoint locus.
Multiplication concatenates the block lists and uses the common disjoint-locus identification.
These are the actual multiplication and differential of the tensor cobar construction.

To\phantomsection\label{ogd:def:block-density-generators} specify the signs, choose a reference order on the labels and local coordinates on the curve.
For a block \(B=\{b_1<\cdots<b_m\}\), write
\[
 c_B=dz_{b_1}\wedge\cdots\wedge dz_{b_m},
 \qquad |c_B|=-m\text{ in }J_B,
 \qquad \xi_B=tc_B.
\]
For \(B=A\sqcup C\), define \(\varepsilon(A,C)\in\{1,-1\}\) by
\(c_A\wedge c_C=\varepsilon(A,C)c_B\).
The density component of the coproduct is
\(c_B\mapsto\varepsilon(A,C)c_A\otimes c_C\).
Equation~\eqref{ogd:eq:cobar-signs} becomes
\begin{equation}\label{ogd:eq:partition-sign}
 d\xi_B=\sum_{\substack{B=A\sqcup C\\A,C\ne\varnothing}}
       (-1)^{|A|+1}\varepsilon(A,C)\xi_A\xi_C.
\end{equation}
On a word \(\xi_{B_1}\cdots\xi_{B_r}\), the split \(B_i=A\sqcup C\) has sign
\begin{equation}\label{ogd:eq:word-sign}
 (-1)^{\sum_{j<i}(1-|B_j|)+|A|+1}\varepsilon(A,C).
\end{equation}
Each coefficient map is the localization map in~\eqref{ogd:eq:J-coprod}.
The density identifications are ordered external-product identifications.
Replacing them with a common coordinate volume form must retain the displayed permutation signs.

The square-zero identity has a direct partition interpretation.
Refinements of two different blocks cancel by the odd derivation rule.
For one block split into \((A|B|C)\), the two intermediate partitions are \((A\sqcup B|C)\) and \((A|B\sqcup C)\).
Wedge associativity gives the same total density sign on both routes.
The exponents from~\eqref{ogd:eq:word-sign} differ by one.
Thus each triple refinement cancels, proving the identity in every arity.
All coordinate changes preserve the canonical-line maps, so the formula glues globally.

Writing \(a_i=\xi_{\{i\}}\), the first differentials are
\begin{align}
 d\xi_{12}&=a_1a_2-a_2a_1,\label{ogd:eq:binary-sign}\\
 d\xi_{123}&=a_1\xi_{23}-a_2\xi_{13}+a_3\xi_{12}
              -\xi_{12}a_3+\xi_{13}a_2-\xi_{23}a_1.\label{ogd:eq:ternary-sign}
\end{align}
The ternary complex has one, six, and six partition summands in degrees \(-2,-1,0\).
Substituting~\eqref{ogd:eq:binary-sign} in~\eqref{ogd:eq:ternary-sign} yields
\begin{align*}
 &a_1(a_2a_3-a_3a_2)-a_2(a_1a_3-a_3a_1)+a_3(a_1a_2-a_2a_1)\\
 &\quad -(a_1a_2-a_2a_1)a_3+(a_1a_3-a_3a_1)a_2
              -(a_2a_3-a_3a_2)a_1=0.
\end{align*}
Every ordered word of length three occurs twice with opposite signs.
On the completely disjoint locus, all coefficient localizations become identities.
The ordered partitions then give the oriented cellular chain complex of the permutahedron, placed in degrees \(1-n,\ldots,0\).
Along a collision, the same differential retains the localizations in~\eqref{ogd:eq:U-partitions}.

\subsection{The binary principal part and its bar cycle}

Take \(X=\mathbb A^1_k\) and restrict to either chart of \(\mathscr R_{\{1,2\}}\).
Use \(R,L,u\) from~\eqref{ogd:eq:opening-binary}.
In ordered external-density bases \(e_{12}=a_1a_2\), \(e_{21}=a_2a_1\), formula~\eqref{ogd:eq:U-partitions} is exactly
\[
 [R\xi\longrightarrow Le_{12}\oplus Le_{21}],\qquad r\xi\longmapsto(r,-r).
\]
All three modules carry the corresponding right density action.
It follows that \(H^{-1}(\ogdUU_2)=0\) and
\[
 H^0(\ogdUU_2)=(L\oplus L)/\{(r,-r):r\in R\}.
\]
An explicit decomposition is
\begin{equation}\label{ogd:eq:binary-split}
 [(f,g)]\longmapsto(f+g,[f])\in L\oplus L/R,
 \qquad
 (h,[p])\longmapsto[(p,h-p)].
\end{equation}
Changing the lift \(p\) by \(r\in R\) changes the inverse by \((r,-r)\).
This proves that both maps are well-defined and mutually inverse.

For\phantomsection\label{ogd:def:right-density-coordinate-action} a constant density generator \(\nu\), the right action is
\[
 (f\nu)\cdot\partial_{z_i}=-(\partial_{z_i}f)\nu.
\]
Localization, quotient, and the constant matrices in~\eqref{ogd:eq:binary-split} commute with this action.
Therefore the decomposition is differential-module linear:
\[
 H^0(\ogdUU_2)\simeq\omega_{X^2}(*\Delta)
 \oplus\bigl(\omega_{X^2}(*\Delta)/\omega_{X^2}\bigr)
\]
on the chosen chart.
The injection \([p]\mapsto[(p,-p)]\) identifies the second summand.

The\phantomsection\label{ogd:def:fixed-coefficient-exchange} chosen splitting also determines how coefficient exchange acts.
Hold the variables and density bases fixed, and set $E(f,g)=(g,f)$.
Equation~\eqref{ogd:eq:binary-split} gives
\[
 E(h,[p])=(h,[h]-[p]),\qquad E^2=1.
\]
For $(f,g)=(u^{-1},0)$, the exchanged pair has splitting coordinates $(h,[p])=(u^{-1},0)$.
Thus this action differs from $(h,[p])\mapsto(h,-[p])$.

The half-difference coordinate $[q]=[(f-g)/2]$ instead gives $(h,[q])\mapsto(h,-[q])$.
Its inverse sends $(h,[q])$ to $[(h/2+q,h/2-q)]$.
Changing the lift of $[q]$ changes this pair by $(r,-r)$ for $r\in R$.

These formulas describe the fixed-coefficient involution $E$.
A geometric label permutation also pulls back the coordinates and permutes the density factors.
For example, interchanging $z_1,z_2$ sends $u$ to $-u$ and negates $dz_1\wedge dz_2$.
Differential-module linearity of the splitting does not remove those operations.

The class
\begin{equation}\label{ogd:eq:singular}
 \beta=\left[\frac{e_{12}-e_{21}}{u}\right]
\end{equation}
corresponds to \((0,[u^{-1}])\), hence is nonzero.
In coordinates \(u=z_1-z_2,v=z_2\), its density is the residue generator \([du\wedge dv/u]\).
Normal differentiation generates the higher principal parts.
The quotient is the diagonal-supported right differential module \(\Delta_+\omega_X\).

The\phantomsection\label{ogd:def:binary-bar-cycle} binary bar contains the relation and its primitive together.
With \([a]=sa\), set
\[
 Z=[\xi]+[a_1|a_2]-[a_2|a_1].
\]
It has degree \(-2\) and
\[
 bZ=-[e_{12}-e_{21}]+[e_{12}]-[e_{21}]=0.
\]
In fact, its complete reduced binary complex is
\begin{equation}\label{ogd:eq:binary-bar}
 R\oplus L\oplus L\longrightarrow L\oplus L,
 \qquad (r,f,g)\longmapsto(-r+f,r+g),
\end{equation}
in degrees \(-2,-1\).
The map is surjective, and its kernel consists of \((r,r,-r)\).
Restoring the density gives \(H^{-2}(B\ogdUU)_2=\omega_{X^2}\), with every other cohomology group zero on this chart.
This is the component \(J_2=\omega_{X^2}[2]\).
The length-one suspension of~\eqref{ogd:eq:singular} is killed by
\[
 b\bigl(u^{-1}([a_1|a_2]-[a_2|a_1])\bigr)
       =[u^{-1}(e_{12}-e_{21})].
\]
The bar comparison consequently preserves both the singular multiplication and the higher differential that resolves it.

\section{The full derived coefficient complex}\label{ogd:sec:cech}

The coalgebra \(J\) lies on \(\mathscr R_I\).
To express its coefficients on the ordinary position space \(X^I\), apply the explicit morphism \(q_I:\mathscr R_I\to X^I\).
Its full derived direct image retains contributions from every intersection of order charts.

Choose\phantomsection\label{ogd:def:ordered-chart-cech-data} an auxiliary total order on the finite set \(\Sigma_I=\ogdOrd(I)\).
For nonempty \(S\subseteq\Sigma_I\), let \(V_S\subset X^I\) remove every diagonal whose pair has varying order among the elements of \(S\).
It is the image of the corresponding chart intersection.
Let \(j_S:V_S\hookrightarrow X^I\) be the open inclusion.
Define the complex of left differential modules
\begin{equation}\label{ogd:eq:cech}
 \check C_I^p=\bigoplus_{\sigma_0<\cdots<\sigma_p}
      j_{\{\sigma_0,\ldots,\sigma_p\},*}\mathcal O_{V_{\{\sigma_0,\ldots,\sigma_p\}}},
 \qquad d=\sum_{a=0}^{p+1}(-1)^a\operatorname{res}_a.
\end{equation}
Different subsets of charts remain different summands, even when their geometric intersection opens happen to coincide.
There are \(|\Sigma_I|\) columns, in degrees zero through \(|\Sigma_I|-1\).

\begin{proposition}\label{ogd:prop:directimage}
For \(n=|I|\), there are equivalences of right differential modules on \(X^I\),
\begin{equation}\label{ogd:eq:realization}
 q_{I,+}(B\ogdUU)_I\simeq q_{I,+}J_I
 \simeq\bigl(\omega_{X^I}\otimes_{\mathcal O_{X^I}} \check C_I\bigr)[n].
\end{equation}
The tensor on the right converts the left connection on \(\check C_I\) into a right differential module.
\end{proposition}
\begin{proof}
The first equivalence is the direct image of the unit \(J\to B\Omega J\).
For the second, differential-module direct image is computed by the transfer module:
\[
 q_{I,+}M=Rq_{I,*}^{\mathcal O}
  \bigl(M\otimes^L_{\mathcal D_{\mathscr R_I}}
                    \mathcal D_{\mathscr R_I\to X^I}\bigr).
\]
Since \(q_I\) is locally an isomorphism, its transfer module is \(\mathcal D_{\mathscr R_I}\), with the induced base differential-operator action.
The transfer tensor contributes neither a relative de Rham differential nor a dimension shift.
The remaining derived sheaf direct image is computed by the finite chart cover.

Locally on \(X^I\), each \(V_S\) is the complement of a Cartier divisor given by a product of diagonal equations.
Its inclusion is affine, and its quasi-coherent direct image is localization.
Thus the intersections are acyclic for this direct-image calculation.
The augmented sheaf \v{C}ech resolution is exact: near any point, choose a chart containing it and insert that chart as an extra index to contract the augmented complex.
Applying direct image therefore gives~\eqref{ogd:eq:cech}.
The canonical line on \(\mathscr R_I\) is the pullback of \(\omega_{X^I}\).
The projection formula and its right differential-module action give the final expression in~\eqref{ogd:eq:realization}.
\end{proof}

For\phantomsection\label{ogd:def:chart-intersection-coordinate-rings} \(X=\mathbb A^1\), put
\[
 R_S=k[z_i:i\in I]
 [(z_a-z_b)^{-1}:\text{the order of }a,b\text{ varies in }S].
\]
These are the coefficient rings in~\eqref{ogd:eq:cech}.
The complex retains the coefficientwise flat connection \(\partial_{z_i}\).
For two labels, it is
\[
 \check C_2=[R\oplus R\xrightarrow{(f,g)\mapsto g-f}L],
 \qquad H^0(\check C_2)=R,\quad H^1(\check C_2)=L/R.
\]
This is a derived function complex on a scheme with two diagonal branches.
Equation~\eqref{ogd:eq:realization} realizes it by the associative chiral bar with its density and shift.

\subsection{A commutative Thom--Whitney model}

The \v{C}ech differential is explicit, but a strictly commutative product is more convenient on polynomial differential forms.
For\phantomsection\label{ogd:def:polynomial-simplex-coefficients} a finite nonempty set \(S\), define
\[
 \Omega_{\mathrm{poly}}(\Delta^S)=
 k[t_\sigma,dt_\sigma:\sigma\in S]/
       (\textstyle\sum t_\sigma-1,\sum dt_\sigma),
\]
where \(|t_\sigma|=0\), \(|dt_\sigma|=1\), and \(d(t_\sigma)=dt_\sigma\).
For \(T\subset S\), face restriction sets \(t_\sigma=dt_\sigma=0\) for \(\sigma\notin T\).
Let \(L_I=R_{\Sigma_I}\), the ring in which all pairwise differences are invertible.
The Thom--Whitney algebra is
\begin{equation}\label{ogd:eq:TW}
 \ogdTW_I=
 \left\{\theta\in L_I\otimes\Omega_{\mathrm{poly}}(\Delta^{\Sigma_I}):
 \theta|_{\Delta^S}\in R_S\otimes\Omega_{\mathrm{poly}}(\Delta^S)
 \text{ for all }\varnothing\ne S\subseteq\Sigma_I\right\}.
\end{equation}
Its differential acts only on the simplex variables.
The base connection differentiates the \(z_i\)-coefficients and preserves every \(R_S\).
Equivalently, \(\ogdTW_I\) is the compatible family of forms on all faces with coefficients in their assigned rings.
The inclusions \(R_S\hookrightarrow L_I\) identify that family with~\eqref{ogd:eq:TW}.

\begin{proposition}\label{ogd:prop:TW}
Integration over oriented faces defines a connection-preserving quasi-isomorphism of complexes
\[
 \mathcal I:\ogdTW_I\longrightarrow \check C_I.
\]
The algebra \(\ogdTW_I\) is a commutative differential graded model of \(Rq_{I,*}^{\mathcal O}\mathcal O_{\mathscr R_I}\).
After left-to-right conversion, it may replace \(\check C_I\) in~\eqref{ogd:eq:realization}.
\end{proposition}
\begin{proof}
On an ordered \(p\)-simplex, integration is the algebraic operation determined by
\[
 \int_{\Delta^p}t_0^{a_0}\cdots t_p^{a_p}\,dt_1\cdots dt_p
       =\frac{a_0!\cdots a_p!}{(a_0+\cdots+a_p+p)!}.
\]
It vanishes on forms of degree different from \(p\).
Characteristic zero makes these rational coefficients available in \(k\).
Algebraic Stokes gives the alternating face differential, so \(\mathcal I\) is a chain map.
It commutes with the base connection because it acts only on the polynomial simplex factors.

For completeness, the quasi-isomorphism follows from the finite skeleton filtration of the simplex.
Restriction of polynomial forms to its boundary is surjective onto compatible boundary forms.
One can extend one face at a time in barycentric coordinates, subtracting the previous face restrictions before extending the next.
The kernel consists of relative forms on the simplex and its boundary.
Polynomial contraction to a vertex proves the polynomial Poincar\'e lemma: pull back along \(H(s,t)=(1-s)v+st\), take its \(ds\)-coefficient, and integrate in \(s\).
It gives \(dh+hd=\ogdid-\operatorname{ev}_v\).
The exact sequence for the boundary, followed by induction on dimension, shows that relative polynomial forms have one cohomology class in the top degree.
Integration identifies that class with \(k\).

Apply this argument successively to every face, tensoring its relative forms with the assigned coefficient ring \(R_S\).
Restriction to a newly attached face is unrestricted in \(R_S\), while its boundary values have already been fixed.
Thus the associated graded of the skeleton filtration is the direct sum of the relative simplex complexes with coefficients \(R_S\).
The integration map on it is the top-degree isomorphism just proved.
The filtration is finite, so \(\mathcal I\) is a quasi-isomorphism.

The compatible-form algebra computes the homotopy limit of the full chart diagram.
Indeed, the surjective boundary restriction maps give its successive homotopy matching maps, and the polynomial contraction gives the required resolutions of constants.
The \v{C}ech descent calculation then identifies this homotopy limit with the derived structure-sheaf direct image.
Its multiplication is the pointwise wedge product of compatible forms.
\end{proof}

The affine-line algebra~\eqref{ogd:eq:TW} is the derived-section model of \ogdciteAKY*{\S4.2 and Appendix A}.
The all-face condition specifies its full domain.
For a general curve, the same construction uses the sheaves \(j_{S,*}\mathcal O_{V_S}\) and glues over coordinate opens.
Direct image to a point is a further operation and introduces the de Rham complex for differential modules.
It is not part of~\eqref{ogd:eq:realization}.

\section{Diagonal support and ordered words}
\label{ogd:sec:recognition}

Which coalgebras have a cobar concentrated on the diagonal branches?
The answer depends on which ordered coproduct is required to be invertible.
At two distinct points the bar of a diagonal algebra has two words,
$[1|2]$ and $[2|1]$. Each coproduct component selects one word.
Their joint map retains both. This distinction determines the restriction
of Theorem~\ref{ogd:thm:duality}.

For\phantomsection\label{ogd:def:ordered-small-diagonal-insertion}\phantomsection\label{ogd:def:h373p04-ordered-diagonal-insertion}\phantomsection\label{ogd:def:ordered-diagonal-coefficient-system} a nonempty set $I$, let $U_I\subset\mathscr R_I$ be the fully distinct
locus and let $D_I$ be the union of its small-diagonal branches.
Write $\delta_\sigma:X\hookrightarrow\mathscr R_I$ for the branch with
order $\sigma\in\ogdOrd(I)$. For a right differential module $N$ on $X$, put
\[
 (\iota_{\mathrm{ord}}N)_I
   =\bigoplus_{\sigma\in\ogdOrd(I)}\delta_{\sigma,*}N,
 \qquad (\iota_{\mathrm{ord}}N)_\varnothing=0.
\]
An ordered-fibre collapse pulls back precisely the branches whose orders
are expanded quotient orders. Proper base change therefore makes this a
coefficient system. Kashiwara's equivalence and the collision identifications
give a natural equivalence $\iota_{\mathrm{ord}}(F_1)\simeq F$ whenever
$F_\varnothing=0$ and $F_I$ is supported on $D_I$.

For\phantomsection\label{ogd:def:coalgebra-singleton-coefficient} a coaugmented coalgebra $C$ with $C_\varnothing=k$, set $N=C_1$.
Its iterated coproducts give the joint map
\begin{equation}\label{ogd:eq:joint-recognition}
 \Theta_I:C_I|_{U_I}\longrightarrow
 \bigoplus_{\sigma\in\ogdOrd(I)}
 (N_{\sigma(1)}\boxtimes\cdots\boxtimes N_{\sigma(n)})|_{U_I},
 \qquad n=|I|.
\end{equation}
Call $C$ \emph{jointly word-factorising} when these maps are equivalences
for every $I$. This is a condition on the supplied coproducts, rather than
the existence of some isomorphism between the displayed objects.

\begin{theorem}\label{ogd:thm:joint-recognition}
Bar and cobar restrict to inverse equivalences between augmented collision
algebras whose positive coefficients are supported on $D_I$ and jointly
word-factorising collision coalgebras. For such a coalgebra the proper
branch counits give the coefficient equivalence
\begin{equation}\label{ogd:eq:diagonal-recognition}
 \iota_{\mathrm{ord}}(C_1[-1])\xrightarrow{\sim}\overline{\Omega C}.
\end{equation}
The equivalence includes the derived mapping spaces of the two full
subcategories.
\end{theorem}

\begin{proof}
Restriction to the fully distinct loci is strong monoidal: every summand
of convolution becomes its external coefficient product there.
If $A=K\oplus a$ is diagonal-supported, a bar factor occupying more than
one label vanishes on $U_I$. The surviving terms are the singleton words,
one for each order of $I$, and every multiplication face has zero target.
The restricted bar is thus the tensor coalgebra on the singleton coefficient
$a_1[1]$, with its actual deconcatenation maps. Its joint map is an equivalence.

Conversely, write $\operatorname{Sep}$ for restriction to the fully distinct
loci, with tensor product summed over all ordered subsets of $I$.
Let $M$ be the separated coefficient system with $M_1=C_1$
and all other positive components zero. Projection to the singleton,
followed by all iterated coproducts, defines a coalgebra map
\[
 \operatorname{Sep}(C)\longrightarrow
 T^c(M)=K\oplus\bigoplus_{r\geq1}M^{\star r}.
\]
Positivity makes the sum finite at each $I$.
Coassociativity identifies this map's deconcatenations
with its iterated coproducts; its component is exactly
\eqref{ogd:eq:joint-recognition}. The hypothesis makes it an equivalence.
Separation commutes with normalized cobar. For a fixed singleton order,
the cobar of $T^c(M)$ has the gap complex used in the proof of
Theorem~\ref{ogd:thm:duality}. It contracts when $|I|>1$ and leaves
$C_1[-1]$ when $|I|=1$.

Every equality stratum with $r$ distinct positions and specified internal
orders is the restriction of an ordered-fibre collision
$\delta_\alpha:\mathscr R_J\to\mathscr R_I$ to $U_J$, where $|J|=r$.
Its extraordinary restriction of $\Omega C$ is therefore zero for $r>1$.
The finite filtration by equality strata and open--closed recollement
put $\overline{\Omega C}$ on $D_I$. The proper counits on its branches
give \eqref{ogd:eq:diagonal-recognition}, compatibly with collisions.
The ambient unit and counit preserve these full subcategories, so the
restricted functors are inverse on objects and mapping spaces.
\end{proof}

The theorem is the collision-space recognition of
Theorem~\ref{interfaces-thm:word}. It uses all ordered-fibre collisions
to pass from distinct points to partially collided points.
Its\phantomsection\label{ogd:def:binary-joint-graph-map} binary form also explains why invertibility of each component is a
different requirement. If both binary components are equivalences, use
the first to identify $C_2|_U$ with
$V=(C_1\boxtimes C_1)|_U$, and the symmetry to identify the reverse target
with $V$. The joint map is $(1,R):V\to V\oplus V$ for an equivalence $R$.
The quotient map $(x,y)\mapsto y-Rx$, with section $v\mapsto(0,v)$,
splits its cofibre as $V$. The two-term cobar consequently has
\begin{equation}\label{ogd:eq:strong-graph}
 (\overline{\Omega C})_2|_U
   \simeq\operatorname{cofib}(1,R)[-2]\simeq V[-2].
\end{equation}
This is nonzero whenever $V$ is nonzero, including for $R=1$.
In the other direction, the two components of the bar of a diagonal
algebra are the two projections from its two singleton words. Thus a
nonzero separated singleton square prevents componentwise strongness
from replacing joint factorisation on the collision carrier.

\section{Local cohomology of the reduced cut cobar}
\label{ogd:sec:cut-cobar-support}

If only consecutive labels may be collapsed, the equality $x_1=x_3$
does not shorten the word $(x_1,x_2,x_3)$ when $x_2$ is distinct.
It nevertheless prevents either cut from separating that word into
disjoint positions. The cobar retains precisely such configurations.
We now determine their support and its algebra structure in every arity.
Return here to the word category $\mathcal W_X$, for a smooth separated
scheme $X$ of any positive pure dimension.

Bar and cobar for this product will be denoted $B_{\mathrm{cut}}$ and
$\Omega_{\mathrm{cut}}$. Their normalized words have at most $n$ positive
factors at arity $n$. The finite gap contraction of
Theorem~\ref{ogd:thm:duality} thus applies to this product as well and gives
an ambient equivalence of augmented algebras and coaugmented coalgebras
with empty coefficient $k$.

\begin{theorem}\label{ogd:thm:cut-local-cohomology}
For a strong cut coalgebra $C$, there is a natural equivalence of positive
word coefficient systems
\begin{equation}\label{ogd:eq:cut-local-cohomology}
 \overline{\Omega_{\mathrm{cut}}C}_n
       \simeq R\Gamma_{Z_n}(C_n)[-1].
\end{equation}
It commutes with monotone collisions and coalgebra morphisms. Its composite
with $R\Gamma_{Z_n}(C_n)[-1]\to C_n[-1]$ is the length-one projection.
\end{theorem}

\begin{proof}
For a subset $S$ of the $n-1$ gaps, let $P_S$ be its consecutive partition,
let $U_S=\bigcap_{p\in S}U_p$, and write $j_S$ for its inclusion.
Take $U_\varnothing=X^n$. The cobar term for this partition, before
desuspension, is
\[
 Q_C(S)=j_{S,*}
       \left(\mathop{\boxtimes}_{B\in P_S}C_B\right)|_{U_S}.
\]
Iterated strong coproducts identify the entire coherent refinement diagram
with $S\mapsto j_{S,*}j_S^!C_n$. Coassociativity gives compatibility with
each face and its higher homotopies. These are equivalences of diagrams,
not just of their vertices.

The nonempty vertices form the finite \v Cech diagram for
$O_n=\bigcup_pU_p$. Its totalization computes $j_{O_n,*}j_{O_n}^!C_n$.
Including the empty vertex takes the fibre of
$C_n\to j_{O_n,*}j_{O_n}^!C_n$, hence computes $R\Gamma_{Z_n}C_n$.
The cobar has one desuspension for every block, whereas a cut set $S$
has cubical degree $|S|$ and $|S|+1$ blocks. This accounts for $[-1]$.

Explicitly, write $t=[-1]$. For homogeneous block coefficients use
\begin{equation}\label{ogd:eq:cut-suspension}
 (tc_1)\otimes\cdots\otimes(tc_r)\longmapsto
 (-1)^{\sum_{j=1}^r(r-j)|c_j|}
 (c_1\otimes\cdots\otimes c_r)[-r].
\end{equation}
Splitting the $i$th block into degrees $a,b$ changes this exponent by
$\sum_{j<i}|c_j|+a$. The cobar splitting sign has exponent
$\sum_{j<i}(|c_j|+1)+a+1$. Their sum is $i$ modulo two.
The resulting sign $(-1)^i$ is the negative of the \v Cech insertion
sign, and the internal differential has sign $(-1)^r$.
These are the signs of the augmented cube shifted by $[-1]$.
For coherent coalgebras, finite total fibres retain the higher faces of
the same diagram. This gives the same equivalence without strictifying
on the original coefficient objects.

Under a monotone collision the cuts inside its fibres disappear, and
all other cuts become quotient cuts. Open base change and
\eqref{ogd:eq:Z-pullback} identify both cubes and their augmentations.
They give collision compatibility and naturality. Projection to the
empty cut set is the asserted length-one projection.
\end{proof}

\subsection{Multiplication on the localization resolution}

The supported object in \eqref{ogd:eq:cut-local-cohomology} retains a
multiplication because $C$ supplies inverse factorisation maps on the
cut opens. The finite resolution makes this multiplication explicit.
Set\phantomsection\label{ogd:def:cut-localization-terms} $L_S(C_n)=j_{S,*}j_S^!C_n$ and let
\begin{equation}\label{ogd:eq:cut-resolution}
 \mathcal K(C)_n=\operatorname{Tot}
   \left(\bigoplus_{S\subseteq[n-1]}L_S(C_n)[-|S|-1]\right).
\end{equation}
The internal differential on the $S$ term is $(-1)^{|S|+1}d$;
adding a gap $p$ has sign $(-1)^{1+\#\{a\in S:a<p\}}$.
For a coherent diagram, \eqref{ogd:eq:cut-resolution} denotes its finite
coherent totalization with these binary sign conventions.

Let\phantomsection\label{ogd:def:cut-resolution-product-data} $m_p$ be the inverse of $C_{p+q}|_{U_p}\to(C_p\boxtimes C_q)|_{U_p}$.
For $S\subseteq[p-1]$ and $T\subseteq[q-1]$, put
\[
 r=|S|+1,\qquad s=|T|+1,\qquad
 G=S\cup\{p\}\cup(p+T).
\]
The common open for these two terms in their cut product is $U_G$.
Open base change identifies the product coefficient with
$j_{G,*}j_G^!(C_p\boxtimes C_q)$. Restrict $m_p$ to this open and
extend by $j_{G,*}$. On unshifted homogeneous coefficients the product is
\begin{equation}\label{ogd:eq:cut-resolution-product}
 u_Sv_T=(-1)^{s|u|}\bigl(m_p(u\boxtimes v)\bigr)_G.
\end{equation}
Here $|u|$ is its degree in $C_p$; its total degree in the resolution is
$|u|+r$. Empty multiplication is the scalar unit action.

\begin{proposition}\label{ogd:prop:cut-resolution-algebra}
These operations and the coherences of $m$ make $K\oplus\mathcal K(C)$
an augmented $E_1$ cut algebra. The cube comparison gives a natural
equivalence of augmented algebras
\[
 \Omega_{\mathrm{cut}}C\simeq K\oplus\mathcal K(C).
\]
\end{proposition}

\begin{proof}
For a third term of block length $h$, put $a=|u|$ and $b=|v|$.
The two associative signs are $sa+h(a+b)$ and $hb+(s+h)a$, which agree.
Both coefficient maps are the inverse factorisation map on the same
final disjoint open. Its associativity homotopies, and their higher
compatibilities for all final partitions, give the $E_1$ coherences.

Adding a cut in the first word has the same position before and after
multiplication. Adding one in the second shifts its insertion position
by $r$. On the Leibniz side its sign changes by the first term's total
degree $a+r$ and by the new product sign $(s+1)a$; their sum is
$r+sa$ modulo two, precisely the required shift.
For the internal differential on the first coefficient, the two signs
are $sa+r+s$ and $r+s(a+1)$; on the second they are
$sa+r+s+a$ and $a+r+s+sa$. They agree in both cases.
This proves the Leibniz identity with coefficient differentials retained.

Finally the exponent in \eqref{ogd:eq:cut-suspension} for a concatenation
differs from the two separate exponents by $s|u|$.
The iterated coproduct identifies its coefficient map with
\eqref{ogd:eq:cut-resolution-product}. The coherent cube comparison
therefore intertwines concatenation, the unit and their higher coherences.
\end{proof}

This product uses the strong structure maps of $C$ on specified opens.
The local-cohomology object alone does not specify it.

\subsection{The restricted equivalence}

Call an augmented cut algebra $Z$-supported if its augmentation ideal
in arity $n$ is supported on $Z_n$.

\begin{theorem}\label{ogd:thm:Z-restriction}
With the empty coefficient fixed to $k$, bar and cobar restrict to
inverse equivalences
\[
 \ogdAlg_{E_1}^{\mathrm{aug},Z}(\mathcal W_X)
 \ \underset{\Omega_{\mathrm{cut}}}{\overset{B_{\mathrm{cut}}}
 {\rightleftarrows}}\
 \ogdCoalg_{E_1}^{\mathrm{coaug},\mathrm{strong}}(\mathcal W_X).
\]
These are full subcategories of the ambient enhanced categories, so the
equivalence includes their mapping spaces.
\end{theorem}

\begin{proof}
The cobar of a strong coalgebra is $Z$-supported by
Theorem~\ref{ogd:thm:cut-local-cohomology}.
Conversely, restrict the bar of $A=K\oplus a$ to a cut open $U_p$.
A bar factor whose interval crosses $p$ has an admissible internal cut
there, so its $Z$-supported coefficient vanishes.
Every surviving bar word therefore has a separator at $p$.
Deconcatenation at that separator identifies these terms with the two
bar words on the prefix and suffix, with the same external coefficients.
A multiplication face crossing the separator has a zero target; every
remaining face belongs to one of the two factors. The same statement
holds for higher multiplication faces and their coherences.
Thus the actual coproduct component is an equivalence on $U_p$.
The ambient bar--cobar unit and counit now restrict to the stated full
subcategories and give the result.
\end{proof}

Small-diagonal support implies $Z$-support, but the converse already
fails at three positions. The exact additional condition can be read
on equality strata. Let\phantomsection\label{ogd:def:equality-stratum-inclusion} $i_\pi:S_\pi\hookrightarrow X^n$ be the locus
whose equalities are exactly the blocks of a set partition $\pi$.
Call $\pi$ cut-indecomposable if each gap is crossed by some block.
Equivalently, $S_\pi\subseteq Z_n$.

\begin{corollary}\label{ogd:cor:cut-diagonal}
For a strong cut coalgebra $C$, its reduced cobar is supported on the
small diagonal in every arity if and only if
\[
 i_\pi^!C_n=0
\]
for every cut-indecomposable partition with more than one equality block.
In that case the proper diagonal counit gives
$\Delta_{n,*}(C_1[-1])\simeq\overline{\Omega_{\mathrm{cut}}C}_n$.
\end{corollary}

\begin{proof}
On a stratum contained in $Z_n$, the extraordinary restriction of
$R\Gamma_{Z_n}C_n$ is $i_\pi^!C_n$: the open direct-image term of
localization has zero extraordinary restriction there.
On a stratum outside $Z_n$ the local-cohomology object vanishes.
The finite equality-stratum filtration proves that the displayed
vanishing is exactly support on the small diagonal.
Kashiwara and the monotone collapse to the singleton identify its
coefficient and give the stated counit.
\end{proof}

\subsection{A coefficient invisible to binary data}

The extra stratum can vary while the entire binary structure stays fixed.
For the dualizing system of \eqref{ogd:eq:cut-dualizing}, define in every positive arity
\begin{equation}\label{ogd:eq:hidden-coefficient}
 P_n=\operatorname{cofib}
 \bigl(R\Gamma_{\Delta_n}J_n\longrightarrow R\Gamma_{Z_n}J_n\bigr)[-1].
\end{equation}
Both supports pull back under monotone collapses, so $P$ is a word
coefficient system, with $P_\varnothing=0$. It is $Z$-supported.
The first two coefficients vanish. For
$v:X^2\setminus\Delta\hookrightarrow X^2$, localization inside
$\Delta_{13}$ gives
\[
 P_1=P_2=0,\qquad
 P_3\simeq i_{13,*}v_*J_{X^2\setminus\Delta}[-1].
\]

\begin{proposition}\label{ogd:prop:binary-invisible}
If $X$ has two distinct geometric points, there are nonequivalent strong
cut coalgebras with identical empty, unary and binary data.
For $X=\mathbb A^1$ they may be taken weakly translation equivariant.
\end{proposition}

\begin{proof}
Set $A^0=\Omega_{\mathrm{cut}}J$. Regard $P$ as an $A^0$-bimodule through
the augmentation $A^0\to K$, and form the square-zero extension
$A'=A^0\ltimes P$. Thus multiplication on $P\star_{\mathrm{cut}}P$ is
zero, and its mixed actions factor through $K$.
Both positive summands are $Z$-supported. Theorem~\ref{ogd:thm:Z-restriction}
makes $C'=B_{\mathrm{cut}}A'$ strong in every arity.
Projection onto $A^0$, followed by $B_{\mathrm{cut}}A^0\simeq J$,
is an equivalence through arity two: a positive bar word involving $P$
requires at least three labels.

At arity three such a word can only be the single factor $P[1]$.
Consequently $C'_3\simeq J_3\oplus P_3[1]$.
On the stratum $x_1=x_3\ne x_2$ its extraordinary restriction has two
copies of the dualizing coefficient, whereas that of $J_3$ has one.
Their finite ranks differ, so the objects cannot be equivalent.
All supports, localizations and structural maps commute with simultaneous
translations on the affine line, giving the equivariant assertion.
\end{proof}

The construction retains an isolated empty unit and any supplied binary
normalization. It shows why these data cannot determine an unrestricted
derived strong word object. Its additional ternary coefficient also
locates the failure of concentration in a heart.

\begin{lemma}\label{ogd:lem:binary-invisible-heart}
For $X=\mathbb A^1$, let
$M=i_{13,*}v_*\omega_{X^2\setminus\Delta}$, an ordinary right
$\mathcal D_{X^3}$-module. Then
\[
 C'_3[-3]\simeq\omega_{X^3}\oplus M[-1],\qquad
 H^1(C'_3[-3])\simeq M\ne0.
\]
Thus $C'$ does not satisfy the arity heart condition that
$C'_n[-n]$ be an ordinary right differential module for every $n$.
\end{lemma}

\begin{proof}
The open immersion $v$ is affine and quasi-finite, and $i_{13}$ is
closed. Their direct images are exact for the right-module heart
\cite[\S\S2.1.2--2.1.3]{chhop:BD}. Since
$J_{X^3}=\omega_{X^3}[3]$ and
$J_{X^2\setminus\Delta}=\omega_{X^2\setminus\Delta}[2]$,
the preceding decomposition of $C'_3$ gives the displayed formula.
The module $v_*\omega_{X^2\setminus\Delta}$ is nonzero on its
defining open set, and closed direct image is faithful. Hence $M$
is nonzero. With cohomological shifts, $M[-1]$ has its cohomology
in degree one.
\end{proof}

The heart condition removes this particular example, but does not
restore binary reconstruction. The particle systems of
Theorem~\ref{hpc:thm:full-faithfulness} are regular holonomic,
satisfy that heart condition and have coherent local vacua in every
arity. For the two rank-one exchanges their common fields are
\[
 B=\mathbb C1\oplus\mathbb C\epsilon,\qquad \epsilon^2=0,
 \qquad T=0,\qquad Y(a,z)b=ab.
\]
Every iterated state-field word is therefore the same in the two
systems. Their normal collision quotients $q_+,q_-$ at $13|2$,
defined in \eqref{hpc:eq:quotients}, instead satisfy
\[
 q_+(\xi)=0,\qquad q_-(\xi)=2\epsilon\otimes\epsilon,
 \qquad
 \xi=\epsilon\otimes\epsilon\otimes1
       -1\otimes\epsilon\otimes\epsilon.
\]
The pole $\xi/(x_1-x_3)$ belongs to one extension and not the other,
so the identity of their binary field algebras has no geometric
lift. Retaining the involutive exchange recovers the fully faithful
particle construction of the cited theorem.
Theorems~\ref{egr:thm:realization} and~\ref{egr:thm:fully-faithful}
extend this positive construction to every finite-dimensional unital
algebra with an involutive exchange satisfying the fusion identities.
The unit filtration proves heart concentration on curves without an
augmentation. The resulting sheaf construction works on every smooth
variety of positive dimension; its higher-dimensional coefficients
need not satisfy the curve heart condition.
In particular, the passage across the nonconsecutive ternary diagonal in
\cite{interfaces-Lat}*{Proposition~A.2.2} requires more than the two cut
factorisation isomorphisms: their common disjoint locus omits this stratum.

\subsection{Central coefficients and field operations}

The support theorem also retains a level parameter at its singular fibre.
Let\phantomsection\label{ogd:def:central-coefficient-ring} $R$ be an ordinary commutative $k$-algebra, replace
$\ogdDMod(Y)$ by $\operatorname{Mod}_R(\ogdDMod(Y))$, and take every
external tensor product over $R$ in the derived sense.
The isolated empty coefficient is now $R$.

\begin{proposition}\label{ogd:prop:relative-Z}
The local-cohomology comparison and restricted equivalence of
Theorems~\ref{ogd:thm:cut-local-cohomology} and~\ref{ogd:thm:Z-restriction}
hold with these central coefficients. They commute with derived scalar
extension along a map of ordinary commutative $k$-algebras $R\to R'$.
In particular they hold over $k[\kappa]$ and at $\kappa=0$.
\end{proposition}

\begin{proof}
The gap contractions use the identity of $R$; they require no field
splitting or flatness of the input coefficients.
At each arity all normalized diagrams remain finite.
Strong coproducts identify the same refinement cube using derived
external products over $R$, and the supported-bar argument is unchanged.
The geometric functors are coefficient-linear. For the quasi-compact
opens, a finite affine \v Cech presentation computes their exchange with
scalar extension; exact derived tensoring commutes with this finite
diagram. The coefficient-linear closed direct images and extraordinary
pullbacks give the other exchanges. Hence local cohomology, finite
normalization and the displayed comparisons commute with derived base
change. No inverse limit over arities or division by $\kappa$ is used.
\end{proof}

The tensor product over $R$ is part of the assertion: tensoring over $k$
would retain two separate copies of the level at a binary cut.
Nor does the proposition replace derived specialization by specialization
of cohomology. For instance the complex
$[k[\kappa]\xrightarrow{\kappa}k[\kappa]]$ in degrees $-1,0$
has a derived fibre at zero with a copy of $k$ in each degree.

These results specify the coefficient category recognized by strong cut
factorisation in Definition~\ref{e1f:def:e1}.
The Lie--commutative support restriction of
\ogdciteFG*{Theorem~5.2.1} uses a different tensor construction;
its ambient equivalence is Theorem~5.1.1 there.
To realize independently given fields, one must additionally construct
maps between their simultaneous operation complexes and those of the
recognized algebra, preserving substitutions, translation and the local
vacuum. The master-probe inverse of Theorem~\ref{mpr:thm:recognition}
tests such maps on supplied current coefficients. The supported coefficient
equivalence by itself supplies none of these operation maps.

We return to the collision category $\ogdcC$ on the curve, with convolution
$\star$ and empty coefficient $k$. Its algebra and coalgebra objects also
act on states; the same finite-word mechanism compares these actions.

\section{Modules and comodules}\label{ogd:sec:modules}

An algebra acts on states as well as on itself.
For \(A\) as in Theorem~\ref{ogd:thm:duality}, a left module is an object \(M\in\ogdcC\) with a coherent unital action \(A\star M\to M\).
Its empty component may be an arbitrary complex of \(k\)-vector spaces.
A\phantomsection\label{ogd:def:bar-comodule-coalgebra} left comodule over \(C=BA\) is an object \(N\) with a coherent counital coaction \(\rho:N\to C\star N\).
Both definitions use the enhanced categories, so quasi-isomorphic underlying objects are equivalent when their structure maps agree coherently.

Let\phantomsection\label{ogd:def:module-twisting-cochain} \(\tau:C\to A\) be zero on \(K\), equal to \(sa\mapsto a\) on the length-one term, and zero on longer bar words.
It has degree one and satisfies~\eqref{ogd:eq:twist}.
Define\phantomsection\label{ogd:def:module-comodule-twisting-functors}
\[
 F(M)=C\star_\tau M,\qquad G(N)=A\star_\tau N.
\]
Their underlying objects are \(C\star M\) and \(A\star N\).
The first has coaction \(\Delta_C\star1\), and the second has left multiplication on its \(A\)-factor.
For strict representatives their differentials are
\begin{align}
 D_F&=d_{C\star M}+
 (1\star\operatorname{act})(1\star\tau\star1)(\Delta_C\star1),\label{ogd:eq:module-twist}\\
 D_G&=d_{A\star N}-
 (\mu_A\star1)(1\star\tau\star1)(1\star\rho).\label{ogd:eq:comodule-twist}
\end{align}
The tensor-composition signs are part of these formulas.
For example, a summand of the added term in~\eqref{ogd:eq:module-twist} is
\((-1)^{|c'|}c'\otimes\tau(c'')m\).
A summand of the subtracted term in~\eqref{ogd:eq:comodule-twist} is
\(-(-1)^{|a|}a\tau(n_{-1})\otimes n_0\).
Expanding the square groups its terms into the left side of~\eqref{ogd:eq:twist}, followed by the action or coaction.
Thus both differentials square to zero.

The coherent constructions can retain a module as a matrix entry.
Let\phantomsection\label{ogd:def:matrix-coefficient-category} $\operatorname{Mat}_2(\ogdcC)$ be the category of two-by-two matrices of coefficient objects, with product
\[
 (P\star Q)_{ij}=\bigoplus_{k=1}^2P_{ik}\star Q_{kj},
 \qquad \mathbf1=\operatorname{diag}(K,K).
\]
An augmented upper triangular algebra
\[
 T_A(M)=\begin{pmatrix}A&M\\0&K\end{pmatrix}
\]
is exactly an augmented algebra $A$ together with a coherent left $A$-module $M$.
Matrix multiplication records the action and all its associativity data.
An upper triangular coaugmented coalgebra with diagonal $(C,K)$ and entry $N[1]$ in position $(1,2)$ records a coherent $C$-comodule.

Apply Proposition~\ref{ogd:prop:enhanced-bar} to this monoidal matrix category.
The $(1,2)$ entry of its bar is $(C\star_\tau M)[1]$.
The corresponding cobar entry is $\Omega C\star_\tau N$.
A composable reduced matrix word contains at most one off-diagonal factor.
Every other factor is positive, so the word has at most $|I|+1$ factors at configuration $I$.
The finite gap proof therefore applies even when $M_\varnothing$ or $N_\varnothing$ is nonzero.

The matrix construction also fixes the differential signs.
The identification is
\[
 [a_1|\cdots|a_r|m]\longleftrightarrow
 (-1)^{\sum_i(|a_i|-1)}s([a_1|\cdots|a_r]\otimes m).
\]
Its last bar face gives~\eqref{ogd:eq:module-twist}.

On the suspended coalgebra entry, the coaction sends
$sn$ to $\sum(-1)^{|c'|}c'\otimes sn_0$.
The cobar sign then gives $-(tc')n_0$, with the tensor-prefix sign from preceding factors.
This is~\eqref{ogd:eq:comodule-twist}.
Restricting the enhanced matrix comparisons to these entries supplies their units, counits and triangle coherences.

\begin{theorem}\label{ogd:thm:modules}
The functors \(F,G\) give inverse equivalences
\[
 \ogdMod_A(\ogdcC)\simeq\ogdComod_{BA}(\ogdcC).
\]
All modules and comodules in the stated configuration category satisfy the needed convergence condition.
\end{theorem}
\begin{proof}
The augmentation comparison \(GF(M)\to M\) is induced by
\[
 A\star C\star M\xrightarrow{1\star\epsilon_C\star1}
 A\star M\xrightarrow{\operatorname{act}}M.
\]
The two twisting terms cancel after this map, by their opposite signs in
\eqref{ogd:eq:module-twist}--\eqref{ogd:eq:comodule-twist}.
The complex is the normalized two-sided bar resolution of \(M\).
Before normalization, its augmented simplicial object has an extra degeneracy that inserts the algebra unit at the free end.
The identities between faces and this degeneracy give a contraction: adjacent terms cancel, and the remaining term is the augmentation.
Normalization removes degenerate unit factors without changing the realization.
At \(I\), only \(|I|\) positive bar factors can remain by~\eqref{ogd:eq:mixed-bound}.
Thus the normalized realization and its comparison are finite there.

For the other comparison, first use the universal twisting cochain
\(C\to\Omega C\).
Its comodule comparison has underlying map
\[
 N\xrightarrow\rho C\star N\longrightarrow C\star\Omega C\star N,
\]
where the second arrow inserts the algebra unit.
Coassociativity and~\eqref{ogd:eq:twist} make it a chain map to the twice-twisted complex.
Filter this target decreasingly by the total number of positive \(c\)-factors.
Count the first \(C=K\oplus c\) factor if it is reduced, and count every factor of \(\Omega C\).
Reduced coproduct and reduced comodule coaction strictly increase this count.
On the associated graded, the only additional differential transfers the first reduced factor into the first cobar factor.
For each positive count, it is an identity between the two possible forms
\[
 c\star(tc)^{\star(r-1)}\star N
 \quad\text{and}\quad
 (tc)^{\star r}\star N,
\]
with their degree-one suspension sign.
This two-term complex is contractible.
The count-zero term is \(N\), and the comparison has identity component there.
The mixed bound~\eqref{ogd:eq:mixed-bound} makes this filtration finite after evaluation at \(I\).
The comparison is therefore an equivalence.

For \(C=BA\), the algebra equivalence \(\Omega C\to A\) intertwines the universal twisting cochain with \(\tau\).
Extension and restriction along this equivalence identify their module categories and the displayed twisted constructions.
Hence the comodule comparison remains an equivalence for \(A\).
All comparison maps come from the twisting adjunction, so they obey its triangle identities in the enhancement.
\end{proof}

In particular, \(F(M)\) computes \(K\star_A^L M\) with its bar coaction.
The equivalence identifies external derived mapping complexes and their compositions:
\[
 R\operatorname{Hom}_A(M,M')\simeq
 R\operatorname{Hom}_{BA\text{-comod}}(F(M),F(M')).
\]
This supplies an operation-preserving comparison of representation theories on the ordered geometric carrier.

The isolated unit has a useful precise consequence.
For a positive object \(P\), the external \(k\)-valued complex \(R\operatorname{Hom}_{\ogdcC}(P,K)\) vanishes.
Apply derived module Hom into \(K\) to its normalized free bar resolution.
Free-module adjunction turns every positive term into
\(R\operatorname{Hom}_{\ogdcC}((sa)^{\star r},K)=0\) for \(r>0\).
The remaining term is \(R\operatorname{Hom}_{\ogdcC}(K,K)=k\).
Therefore
\begin{equation}\label{ogd:eq:empty-end}
 R\operatorname{Hom}_{\ogdMod_A(\ogdcC)}(K,K)=k.
\end{equation}
An endomorphism algebra associated with a local vacuum instead requires a specified module \(P\) and the complex \(R\operatorname{Hom}_A(P,P)\).
Equation~\eqref{ogd:eq:empty-end} concerns the external mapping complex, not an unspecified internal endomorphism object.

An\phantomsection\label{ogd:def:internal-hom-to-empty-unit} internal mapping object into \(K\) does no better.
For \(F\in\ogdcC\), write \(\underline{R\operatorname{Hom}}(F,K)\) for an object representing the functor \(V\mapsto\operatorname{Map}_{\ogdcC}(V\star F,K)\).

\begin{proposition}\label{ogd:prop:internal-empty}
For every \(F\in\ogdcC\), the object \(\underline{R\operatorname{Hom}}(F,K)\) exists, and
\[
 \underline{R\operatorname{Hom}}(F,K)\simeq\bigl(R\operatorname{Hom}_k(F_\varnothing,k),0\bigr).
\]
In particular, \(\underline{R\operatorname{Hom}}(BA,K)\simeq K\) for every augmented algebra \(A\) with \(A_\varnothing=k\).
\end{proposition}
\begin{proof}
Since \(K\) is supported at the empty configuration, adjunction in \(D(k)\) gives, naturally in \(V\),
\[
 \operatorname{Map}_{\ogdcC}(V\star F,K)
 \simeq\operatorname{Map}_{D(k)}(V_\varnothing\otimes_kF_\varnothing,k)
 \simeq\operatorname{Map}_{D(k)}\bigl(V_\varnothing,R\operatorname{Hom}_k(F_\varnothing,k)\bigr).
\]
The right side is \(\operatorname{Map}_{\ogdcC}\bigl(V,(R\operatorname{Hom}_k(F_\varnothing,k),0)\bigr)\).
For the bar, \((BA)_\varnothing=k\).
\end{proof}

\begin{corollary}\label{ogd:cor:dualizable}
An object of \(\ogdcC\) is dualizable for \(\star\) if and only if it is supported at the empty configuration and its empty component is a perfect complex.
\end{corollary}
\begin{proof}
The functor \(M\mapsto(M,0)\) from \(D(k)\) is symmetric monoidal, so it carries a perfect complex to a dualizable object.
Conversely, let \(F\) be dualizable.
Its dual represents \(V\mapsto\operatorname{Map}_{\ogdcC}(V\star F,K)\), so it is supported at the empty configuration by Proposition~\ref{ogd:prop:internal-empty}.
Since \(\star\) is symmetric, \(F\) is the dual of its dual, and it is supported there as well.
The symmetric monoidal functor \(F\mapsto F_\varnothing\) makes \(F_\varnothing\) dualizable, that is, perfect.
\end{proof}

The argument uses only the support of the unit \(K\): holonomic components, or finite-dimensional fibres, do not make a positive object dualizable.

\section{Bimodules and their compositions}\label{ogd:sec:bimodules}

The module equivalence compares one action at a time.
Composing a right action with a left action also requires their common algebra to be retained between the two modules.
Its bar supplies the intervening relations and homotopies.
A common multibar computes composition before and after the two actions are replaced by coactions (Theorem~\ref{ogd:thm:bimodule-composition}).

For\phantomsection\label{ogd:def:relative-bimodule-coefficients} this section, let \(R\) be a commutative \(k\)-algebra and let \(\ogdcC_R\) be the category of \(R\)-module objects in \(\ogdcC\).
Use derived tensor products over \(R\) in the convolution formula; write \(\star\) for this product and \(K_R=R\otimes_kK\) for its unit.
For \(R=k\) these are the preceding category and product.
The coefficient functors, collision comparisons and exact configuration evaluations are their \(R\)-linear versions.
The proof of Lemma~\ref{ogd:lem:finite} still gives
\begin{equation}\label{ogd:eq:bimodule-bound}
 (V_1\star\cdots\star V_r\star M\star N)_I=0
 \qquad(r>|I|)
\end{equation}
for positive \(V_i\) and arbitrary \(M,N\).
Indeed, only the \(V_i\) require nonempty parts of the configuration.
All\phantomsection\label{ogd:def:relative-bar-and-twisting-cochain} algebras below are augmented to \(K_R\) with positive augmentation ideal.
Put \(C_A=BA\), and let \(\tau_A:C_A\to A\) be the degree-one twisting cochain of Section~\ref{ogd:sec:modules}.
The bars use derived representatives, so their normalized formulas retain the coherent operations of the algebras.

\subsection{The two endpoint twists}

An\phantomsection\label{ogd:def:bimodule-and-bicomodule-categories}\phantomsection\label{ogd:def:enhanced-bicomodule-category} \(A\)--\(B\) bimodule is an object \(M\) with coherent unital left \(A\)- and right \(B\)-actions that commute.
A \(C_A\)--\(C_B\) bicomodule has commuting counital left \(C_A\)- and right \(C_B\)-coactions.
Write \(\operatorname{Bimod}_{A,B}(\ogdcC_R)\) and \(\operatorname{Bicomod}_{C_A,C_B}(\ogdcC_R)\) for their enhanced categories.
Define the two-sided transform by
\begin{equation}\label{ogd:eq:double-bar}
 \mathfrak B_{A,B}(M)=C_A\star M\star C_B.
\end{equation}
Its coactions are the two outer coproducts.
For strict representatives, add to the tensor differential the degree-one maps
\begin{align}
 L(c\otimes m\otimes d)
  &=\sum(-1)^{|c'|}c'\otimes\tau_A(c'')m\otimes d,
       \label{ogd:eq:bimodule-left}\\
 T(c\otimes m\otimes d)
  &=-\sum(-1)^{|c|+|m|}c\otimes m\tau_B(d')\otimes d'',
       \label{ogd:eq:bimodule-right}
\end{align}
where \(\Delta c=\sum c'\otimes c''\) and \(\Delta d=\sum d'\otimes d''\).
As in~\eqref{ogd:eq:cobar-signs}, this notation abbreviates maps with their geometric coefficients, not unrestricted elementwise expansions.

The twisting equation~\eqref{ogd:eq:twist} gives \((d+L)^2=(d+T)^2=0\).
The mixed terms \(LT+TL\) vanish because the two actions commute and the contractions have degree one.
Coassociativity makes both twists compatible with the outer coactions.
Thus \(D=d+L+T\) defines a bicomodule differential.
For example, on \([a]\otimes m\otimes[b]\), the two mixed contractions have coefficients
\(-(-1)^{|a|+|m|}\) and \(( -1)^{|a|+|m|}\).
The minus sign in~\eqref{ogd:eq:bimodule-right} also agrees with the right bar differential: for closed degree-zero inputs \(m,a,b\),
\[
 D(m\otimes[a|b])
   =m\otimes[ab]-ma\otimes[b],\qquad
 D^2(m\otimes[a|b])=-m(ab)+(ma)b=0.
\]
These identities include the terms connecting the two ends as well as the one-sided action identities.

\begin{theorem}\label{ogd:thm:bimodules}
Under the enhanced coefficient hypotheses above, the transform induces an equivalence
\[
 \mathfrak B_{A,B}:\operatorname{Bimod}_{A,B}(\ogdcC_R)
 \simeq\operatorname{Bicomod}_{C_A,C_B}(\ogdcC_R).
\]
Its inverse is the composite of the left and right inverse twisting functors; its underlying object on a bicomodule \(P\) is \(A\star P\star B\).
\end{theorem}
\begin{proof}
Apply Theorem~\ref{ogd:thm:modules} to the left action while retaining the commuting right action.
The compatibility just proved makes every comparison right-linear.
Apply its right-handed version to that right action while retaining the left coaction.
The resulting differential is~\eqref{ogd:eq:bimodule-left}--\eqref{ogd:eq:bimodule-right}.
The two inverse comparisons are the corresponding composites of one-sided units and counits.
Their filtrations by the number of positive factors remain finite by~\eqref{ogd:eq:bimodule-bound}, even when the other endpoint has a nonzero empty component.
Their associated graded contractions are therefore the identity-pair contractions of Theorem~\ref{ogd:thm:modules} with the other endpoint retained.
Configuration evaluations detect the resulting equivalences.
The units and counits retain the triangle coherences of the two twisting adjunctions.
For coherent presentations the same argument applies to the derived bar diagrams before normalization.
\end{proof}

\subsection{Balanced tensor and native cotensor}

To compose an \(A\)--\(B\) bimodule \(M\) with a \(B\)--\(D\) bimodule \(N\), form the derived balanced tensor \(M\star_B^{\mathbb L}N\).
By definition it is the realization of the two-sided simplicial bar, whose faces multiply adjacent \(B\)-factors or act on an endpoint, and whose degeneracies insert the unit of \(B\).
Normalization removes unit factors and gives
\begin{equation}\label{ogd:eq:balanced-bar}
 M\star C_B\star N
   =\bigoplus_{r\geq0}M\star(sb)^{\star r}\star N,
 \qquad b=\operatorname{fib}(B\to K_R).
\end{equation}
Its differential consists of the coefficient differential, the bar differential, and the right and left action terms of~\eqref{ogd:eq:bimodule-right} and~\eqref{ogd:eq:bimodule-left}, with their tensor-position signs.

For a right \(C_B\)-comodule \(P\) and a left \(C_B\)-comodule \(Q\), define their native derived cotensor \(P\mathbin{\Box}^{\mathbb R}_{C_B}Q\) as the totalization of the two-sided cosimplicial cobar.
Its cofaces use the two endpoint coactions and the intervening coproducts; its codegeneracies use counits.
If \(\overline C_B=\operatorname{cofib}(K_R\to C_B)\), its normalized term in cosimplicial degree \(q\) is
\begin{equation}\label{ogd:eq:native-cotensor}
 \bigl(P\star\overline C_B^{\star q}\star Q\bigr)[-q].
\end{equation}
The total differential combines the coefficient differential and the alternating cofaces, with the cochain tensor signs.
Equation~\eqref{ogd:eq:bimodule-bound} makes this totalization finite in the normalized cobar direction at each configuration.
Thus this definition uses the coalgebra's own coactions and coproduct, independently of \(\mathfrak B\).

\begin{theorem}[Composition comparison]\label{ogd:thm:bimodule-composition}
For \(M\) and \(N\) as above, the common multibar
\begin{equation}\label{ogd:eq:common-multibar}
 \mathcal Q(M,N)=C_A\star M\star C_B\star N\star C_D
\end{equation}
with all four adjacent twisting terms represents both
\begin{gather*}
 \mathfrak B_{A,D}(M\star_B^{\mathbb L}N),\\
 \mathfrak B_{A,B}(M)\mathbin{\Box}^{\mathbb R}_{C_B}
                   \mathfrak B_{B,D}(N).
\end{gather*}
The comparison to the second expression is induced by the middle coproduct:
\begin{equation}\label{ogd:eq:middle-coproduct}
 c\otimes m\otimes d\otimes n\otimes e
 \longmapsto
 \sum(c\otimes m\otimes d_{(1)})\otimes
          (d_{(2)}\otimes n\otimes e).
\end{equation}
This is a natural equivalence of \(C_A\)--\(C_D\) bicomodules, with the collision and open-locus comparisons of the convolution.
\end{theorem}
\begin{proof}
Substituting~\eqref{ogd:eq:balanced-bar} into the outer transform gives~\eqref{ogd:eq:common-multibar}.
Every face is an internal differential, a multiplication inside a bar, or one of the four adjacent actions.
Thus its differential agrees with the normalized balanced-tensor differential.

For cotensor, first disregard the differential and put \(X=C_A\star M\), \(Y=N\star C_D\).
The ordinary equalizer of the two middle coactions on \((X\star C_B)\star(C_B\star Y)\) is \(X\star C_B\star Y\).
Its inclusion is \(1\star\Delta\star1\), the map~\eqref{ogd:eq:middle-coproduct}; a counit on one adjacent coalgebra factor gives the inverse on the equalizer.
Coassociativity proves both identities.
Colinearity of the two endpoint differentials makes this equalizer a subcomplex, with precisely the four twists of~\eqref{ogd:eq:common-multibar}.
Consequently the inclusion has zero alternating first coface.
It defines a map to the normalized cobar totalization, supported in cobar degree zero.

To prove that this map computes derived cotensor, filter both complexes increasingly by the total number of original \(b\)-factors in all their copies of \(C_B=BB\).
The coefficient differentials and coproducts preserve that number.
Bar multiplications and endpoint \(B\)-actions lower it.
On the associated graded, the two ends are cofree for the retained differential.
The augmented cofree cobar has an extra codegeneracy: apply the counit to the leftmost cofree factor and remove it.
The cosimplicial identities give a contraction of its augmentation; it descends through normalization and preserves the number of original factors.
Hence the associated graded comparison from the equalizer is an equivalence.

At configuration \(I\), the filtration by original factors ends at \(|I|\), and every reduced intervening coalgebra factor contains at least one such factor.
Both that filtration and the normalized cobar direction are therefore finite.
Finite filtered comparison proves that~\eqref{ogd:eq:middle-coproduct} is an equivalence.
No arbitrary infinite product is interchanged with collision pullback.
All maps use the coproducts, actions and coherent geometric comparisons already defining convolution, so the equivalence is natural on that carrier.
\end{proof}

\subsection{The diagonal and an integral interval contraction}

The identity for balanced composition is the regular \(A\)--\(A\) bimodule \(A\).
The identity for cotensor is \(C_A\), with its two coproduct coactions.
Their comparison must therefore retain every cut of a bar word.
Define
\begin{equation}\label{ogd:eq:diagonal-bimodule}
 \iota_A:C_A\longrightarrow\mathfrak B_{A,A}(A),\qquad
 c\longmapsto\sum c_{(1)}\otimes1_A\otimes c_{(2)}.
\end{equation}
Coassociativity makes this a bicomodule map.
The two twists at the inserted unit cancel, so it is a chain map.

For closed degree-zero \(a,b\), write \(1_C\) for the coaugmentation of \(C_A\).
The formula has three terms:
\[
 \iota_A([a|b])=
 1_C\otimes1_A\otimes[a|b]
 +[a]\otimes1_A\otimes[b]
 +[a|b]\otimes1_A\otimes1_C.
\]
Their differentials, in the same order, are
\begin{align*}
 &1_C\otimes1_A\otimes[ab]-1_C\otimes a\otimes[b],\\
 &1_C\otimes a\otimes[b]+[a]\otimes b\otimes1_C,\\
 &[ab]\otimes1_A\otimes1_C-[a]\otimes b\otimes1_C.
\end{align*}
Their sum is \(\iota_A([ab])\).
The middle cut cancels both action terms; omitting it destroys the chain identity.
For a longer word, the same cancellation takes place between adjacent cuts.

\begin{lemma}[Interval contraction]\label{ogd:lem:interval-contraction}
For \(n\geq0\), let \(V_n\) be the free integral cochain complex with basis \(v_0,\ldots,v_n\) in degree zero and \(e_1,\ldots,e_n\) in degree one, and differential
\[
 dv_j=e_j-e_{j+1},\qquad de_j=0,\qquad e_0=e_{n+1}=0.
\]
Define \(i:\mathbb Z[0]\to V_n\), \(p:V_n\to\mathbb Z[0]\) and \(h:V_n\to V_n[-1]\) by
\begin{gather*}
 i(1)=\sum_{j=0}^nv_j,\qquad p(v_j)=\delta_{j0},\qquad p(e_j)=0,\\
 h(e_k)=\sum_{j=k}^nv_j,\qquad h(v_j)=0.
\end{gather*}
Then \(pi=1\), \(dh+hd=1-ip\), and \(ph=hi=h^2=0\).
\end{lemma}
\begin{proof}
For \(k\geq1\), telescoping gives \(dh(e_k)=e_k\).
For \(j>0\), it gives \(hd(v_j)=v_j\), while \(hd(v_0)=v_0-\sum_{j=0}^nv_j\).
These are exactly the stated homotopy identities on the basis.
The side conditions follow since \(h\) vanishes on every vertex and its image contains no \(v_0\).
For \(n=0\) the maps \(i,p\) are identities and \(h=0\).
\end{proof}

\begin{proposition}\label{ogd:prop:diagonal-bimodule}
The map \(\iota_A\) is an equivalence of \(C_A\)--\(C_A\) bicomodules.
\end{proposition}
\begin{proof}
Write \(A=K_R\oplus a\) and filter the target of~\eqref{ogd:eq:diagonal-bimodule} by the total number of positive bar entries, counting a central entry in \(a\) as one further entry.
A central scalar counts zero.
Bar multiplication and action on a positive central entry lower this number; action on a central scalar preserves it.
The coefficient differential is retained.

For a fixed word of \(n\) original \(a\)-factors, the associated graded has a central unit at one of \(n+1\) cuts, or a positive central entry at one of \(n\) positions.
Denote these states by \(v_j\) and \(f_j\), respectively.
For homogeneous elements put \(\alpha_i=|a_i|-1\).
The endpoint formulas give
\[
 d_{\mathrm{cut}}v_j
 =(-1)^{\sum_{i<j}\alpha_i}f_j
   -(-1)^{\sum_{i\leq j}\alpha_i}f_{j+1}.
\]
Set \(f_0=f_{n+1}=0\) for the missing endpoint terms.
Setting \(e_j=(-1)^{\sum_{i<j}\alpha_i}f_j\) identifies this with the complex \(V_n\) of Lemma~\ref{ogd:lem:interval-contraction}, tensored with the suspended coefficient word.
The suspension identifications include its coefficient differential with the usual tensor signs.
The associated graded source is that coefficient word, and \(\operatorname{gr}(\iota_A)\) is \(i\), the sum of all cuts.
The integral contraction therefore proves that the associated graded map is an equivalence.

Each positive entry consumes a nonempty configuration block, so the filtrations end at \(|I|\) after evaluation at \(I\).
Finite filtered comparison and joint conservativity prove the result.
A central entry in \(a\) was counted once, without choosing generators or a PBW basis for \(A\).
The homotopy uses only coefficients \(0,1,-1\), so it requires no inversion in \(R\).
\end{proof}

\subsection{Successive compositions and the internal centre}

For a sequence of augmented algebras \(A_0,\ldots,A_r\) and \(A_{i-1}\)--\(A_i\) bimodules \(M_{i-1,i}\), every parenthesized composition is represented by the same multibar
\begin{equation}\label{ogd:eq:successive-multibar}
 C_{A_0}\star M_{01}\star C_{A_1}\star M_{12}\star\cdots
 \star M_{r-1,r}\star C_{A_r}.
\end{equation}
Comparison with successive cotensors inserts a coproduct at each intervening coalgebra.
Two insertions at the same coalgebra agree by coassociativity; insertions at different coalgebras commute with the tensor signs.
All coefficient identifications use the common ordered-partition expression~\eqref{ogd:eq:many-convolution}.
Thus these maps already commute in the common multibar diagrams before realization.
They give the associativity, unit and pentagon comparisons, and their coherent refinements for longer strings.
The diagonal maps~\eqref{ogd:eq:diagonal-bimodule} give the units, whose identities reduce to the counit identities in the same diagrams.
Normalization and~\eqref{ogd:eq:bimodule-bound} justify the finite realizations and totalizations at each configuration.

In\phantomsection\label{ogd:def:diagonal-bimodule-notation} particular, write \(\operatorname{Bimod}_A\) for \(\operatorname{Bimod}_{A,A}\), and similarly for bicomodules.
Theorem~\ref{ogd:thm:bimodules}, the composition comparison and the diagonal equivalence give a strong monoidal equivalence of enhanced categories
\begin{equation}\label{ogd:eq:native-monoidal}
 (\operatorname{Bimod}_A(\ogdcC_R),\star_A^{\mathbb L},A)
 \simeq
 (\operatorname{Bicomod}_{C_A}(\ogdcC_R),\mathbin{\Box}_{C_A}^{\mathbb R},C_A).
\end{equation}
An internal centre requires additional enrichment data.
A closed enrichment over \(\ogdcC_R\) has internal Hom objects representing maps with arbitrary coefficient objects: for \(V\in\ogdcC_R\),
\[
 \operatorname{Map}_{\ogdcC_R}
  (V,\underline{R\operatorname{Hom}}_{A-A}(M,N))
 \simeq
 \operatorname{Map}_{\operatorname{Bimod}_A}(V\star M,N).
\]
Use the analogous definition for bicomodules.
These are derived mapping spaces and derived representing objects.
An actual \(\ogdcC_R\)-linear equivalence must respect the actions of every \(V\), with their coherences; naturality under extension of the scalar ring alone does not supply this datum.

\begin{corollary}[Conditional internal centre comparison]\label{ogd:cor:internal-centre}
Assume that \(\ogdcC_R\) is closed symmetric monoidal, that the indicated bimodule and bicomodule derived Homs are represented in \(\ogdcC_R\), and that~\eqref{ogd:eq:native-monoidal} is an actual \(\ogdcC_R\)-linear strong monoidal equivalence for those enrichments.
Then it induces an equivalence of internal \(E_2\)-algebras
\begin{equation}\label{ogd:eq:internal-centre}
 \underline{R\operatorname{Hom}}_{A-A}(A,A)
 \simeq
 \underline{R\operatorname{Hom}}^{C_A-C_A}(C_A,C_A).
\end{equation}
Here \(E_2\) refers to the little-two-disks structure on endomorphisms of the monoidal unit, supplied by coherent interchange between composition and tensor operations.
\end{corollary}
\begin{proof}
For each coefficient object \(V\), linearity and full faithfulness identify the two represented mapping functors after applying \(\mathfrak B_{A,A}\).
The diagonal equivalence replaces \(\mathfrak B_{A,A}(A)\) by \(C_A\).
The enriched Yoneda property therefore gives~\eqref{ogd:eq:internal-centre}.
Its map applies \(\mathfrak B_{A,A}\) to an endomorphism and then uses \(\iota_A\) as an equivalence in the derived bicomodule category; no strict chain inverse is required.
The common multibar comparisons preserve tensor operations as well as composition.
Their interchange and higher coherences are the endomorphism-of-the-unit \(E_2\)-structure, so the equivalence preserves that structure.
\end{proof}

The internal Hom in~\eqref{ogd:eq:internal-centre} is an object on the ordered collision carrier.
It is distinct from the external \(R\)-valued mapping complex and from the Hom of a stalk.
For a map of commutative coefficient algebras \(R\to R'\), the bar transforms, multibar maps and integral contraction commute with derived scalar extension, by their configurationwise finite formulas.
An identification of a scalar-extended internal centre with the centre over \(R'\) additionally requires base change for the representing internal Hom.
For a geometric pullback, the corresponding internal-Hom comparison must likewise be an equivalence; compatibility of the bar maps alone does not establish it.

The unit compared here is the identity bimodule \(A\), not the isolated empty module \(K_R\).
These constructions choose no local physical vacuum and prove no generation statement for a boundary category.
In particular, an equivalence on a small curved Heisenberg carrier at nonzero level remains a separate construction, as does a map from physical local operators to the ordered convolution category.

\input{manuscript/constructions/ordered-relative-empty.tex}

\section{The diagonal tensor and its transported product}\label{ogd:sec:diamond}

Convolution\phantomsection\label{ogd:def:h373p04-derived-diagonal-coefficient-tensor} combines disjoint subsets of labels.
There is also a tensor product on coefficients at the same configuration.
For\phantomsection\label{ogd:def:configuration-diagonal-tensor}\phantomsection\label{ogd:def:h373p04-diagonal-tensor}\phantomsection\label{ogd:def:diagonal-coefficient-tensor} a smooth \(Y\), set
\[
 M\ogdotex N=\Delta_Y^!(M\boxtimes N).
\]
It is symmetric monoidal with unit \(J_Y=p_Y^!k\).
Extraordinary inverse image is strong monoidal for this tensor.
Consequently,
\[
 (F\ogdotex G)_I=F_I\ogdotex G_I
\]
defines\phantomsection\label{ogd:def:h373p04-componentwise-diagonal-tensor} a symmetric monoidal product on \(\ogdcC\) with unit \(J\).
At the empty component this is ordinary tensor over \(k\), and \(K\ogdotex K=K\).

\subsection{Interchange and the coalgebra unit}

Both convolutions in \((F\star G)\ogdotex(H\star L)\) are sums over decompositions.
Project to the summands with the same ordered decomposition in both factors.
For such a decomposition \(P\), the strong monoidal restriction \(j_P^!\) has a lax monoidal right adjoint.
Its structure map is
\[
 j_{P,*}M\ogdotex j_{P,*}N\longrightarrow j_{P,*}(M\ogdotex N).
\]
Together with external products, these maps define
\begin{equation}\label{ogd:eq:interchange}
 \zeta:(F\star G)\ogdotex(H\star L)
       \longrightarrow(F\ogdotex H)\star(G\ogdotex L).
\end{equation}
This direction is sufficient to tensor coalgebras.

\begin{proposition}\label{ogd:prop:coalgtensor}
Connected coaugmented counital \(\star\)-coalgebras form a symmetric monoidal category under \(\ogdotex\).
Its unit is \(J\), with the localization coproduct~\eqref{ogd:eq:J-coprod}.
For coalgebras \(C,D\), their coproduct is
\[
 C\ogdotex D\xrightarrow{\Delta_C\ogdotex\Delta_D}
 (C\star C)\ogdotex(D\star D)
 \xrightarrow\zeta(C\ogdotex D)\star(C\ogdotex D).
\]
Their counit and coaugmentation are the tensor products of the original maps, using \(K\ogdotex K=K\).
\end{proposition}
\begin{proof}
The complete coherence can be checked at once for a rectangular family \(F_{ab}\).
Expand each of the \(s\) rows of
\(\bigotimes_{a=1}^{s,!}(\mathop\star_{b=1}^rF_{ab})\).
Keep a term precisely when its \(s\) ordered \(r\)-partitions all coincide.
On the retained term, apply the iterated lax monoidal map of \(j_{P,*}\).
This defines a comparison to
\(\mathop\star_{b=1}^r(\bigotimes_{a=1}^{s,!}F_{ab})\).

Refining a row in two stages retains a term exactly when the coarse partitions and their corresponding refinements agree.
This is equivalent to equality of the final refined partitions.
Combining rows in two stages retains a term exactly when all row partitions agree.
On each retained term, both routes are the same adjunction-induced localization comparison.
Its associativity and symmetry follow from the strong monoidal restriction functor.
These observations prove all substitution and higher coherence diagrams.
Collision pullback preserves this construction: a split-fibre partition vanishes, while a saturated partition has a unique quotient partition.
Equality of saturated partitions is preserved and reflected by passage to that quotient.

The nullary comparisons supply the units.
An empty row family gives \(J\to J^{\star r}\), whose components are the localization units~\eqref{ogd:eq:J-coprod}.
An empty column family gives the iterated identity \(K\ogdotex\cdots\ogdotex K=K\).
Their compatibility follows from composition of localization units and the identities for empty parts.
Thus the coproduct, counit, and coaugmentation of \(C\ogdotex D\) obey the coalgebra identities.

It remains to check that the pointwise constraint \(C\ogdotex J\simeq C\) is a coalgebra map.
For a fixed \(P\), its coproduct is the \(P\)-component of \(\Delta_C\), tensored with the localization unit of \(J\).
The pointwise unit constraint turns this into exactly the same component of \(\Delta_C\).
The right unit constraint has the identical verification.
The counit and coaugmentation constraints hold at the empty component, where all maps are the scalar identities.
Finally, \((C\ogdotex D)_\varnothing=k\otimes k=k\), so connectedness is preserved.
\end{proof}

\subsection{Transport to associative algebras}

The bar--cobar equivalence now defines a second tensor product on connected augmented associative chiral algebras:
\begin{equation}\label{ogd:eq:diamond}
 A\diamond A'=\Omega(BA\ogdotex BA'),\qquad\mathbf1_\diamond=\Omega J=\ogdUU.
\end{equation}
There are coherent equivalences
\[
 B(A\diamond A')\simeq BA\ogdotex BA'.
\]
To construct an associativity or unit constraint, apply \(B\), use the corresponding constraint in Proposition~\ref{ogd:prop:coalgtensor}, and lift along the equivalence \(B\).
The space of such lifts is contractible.
Every coherence diagram commutes after applying \(B\), and its full faithfulness gives the same coherent diagram before applying \(B\).
This proves the entire monoidal structure, including both unit constraints.

At a singleton, only one positive bar or cobar factor survives.
Thus, for augmentation ideals \(a,a'\),
\[
 (A\diamond A')_{\{1\}}\simeq(a_{\{1\}}\ogdotex a'_{\{1\}})[1].
\]
The suspension and the unit \(\ogdUU\) are part of this product.

A coalgebra \(C\) satisfies factorization if the following holds for every ordered partition \(P=(I_1|\cdots|I_r)\) of \(I\) into nonempty parts.
The \(P\)-component of its iterated coproduct restricts to an equivalence
\begin{equation}\label{ogd:eq:factorization}
 j_P^!C_I\xrightarrow{\sim}
 (C_{I_1}\boxtimes\cdots\boxtimes C_{I_r})|_{U_P}.
\end{equation}
The coalgebra \(J\) satisfies this condition because its maps are localization units.
The product \(\ogdotex\) preserves it: restriction turns its coproduct into the tensor product of the two given equivalences.
Hence \(\diamond\) restricts to the full subcategory of algebras whose bars satisfy factorization.
This is an exact geometric condition on the bar of an associative algebra.

\section{Lie algebras and associative envelopes}\label{ogd:sec:CE}

An associative chiral algebra obtained from a Lie algebra has a further explicit comparison.
Let\phantomsection\label{ogd:def:derived-associative-envelope} \(L\in\ogdcC_+\) be a differential graded Lie object for \(\star\).
Its bracket \([\ ,\ ]:L\star L\to L\) has degree zero, is graded antisymmetric, and satisfies graded Jacobi.
Let \(U_\star L\) be the derived associative envelope: it is the left adjoint to the commutator-Lie functor.
For a cofibrant strict presentation it is the tensor algebra modulo
\[
 xy-(-1)^{|x||y|}yx=[x,y].
\]
Its augmentation sends \(L\) to zero, so it has empty component \(k\).

Define\phantomsection\label{ogd:def:chevalley-coalgebra} the Chevalley--Eilenberg coalgebra
\[
 \ogdCE_\star(L)=K\oplus\bigoplus_{r\geq1}\ogdSym^r_\star(sL).
\]
It has the unshuffle coproduct: split the inputs into a subset and its complement, with their Koszul signs.
Its coderivation is determined by
\[
 Q_1(sx)=-s(dx),\qquad Q_2(sx\,sy)=(-1)^{|x|}s[x,y].
\]
The binary formula is symmetric on the suspended inputs, by graded antisymmetry of the bracket.
The square \(Q^2\) is zero: its unary, binary, and ternary components are respectively \(d^2\), the derivation identity, and graded Jacobi.
Coderivation determines every longer component from these identities.

\begin{theorem}\label{ogd:thm:CE}
For every positive Lie object \(L\), there is an equivalence of coaugmented counital coassociative coalgebras
\begin{equation}\label{ogd:eq:CE-compare}
 \Phi:\ogdCE_\star(L)\xrightarrow{\sim}B(U_\star L).
\end{equation}
On homogeneous inputs it is the unnormalized suspended symmetrization
\begin{equation}\label{ogd:eq:CE-map}
 \Phi(sx_1\cdots sx_r)=\sum_{\sigma\in S_r}
 \varepsilon(\sigma;sx_1,\ldots,sx_r)
          [x_{\sigma(1)}|\cdots|x_{\sigma(r)}].
\end{equation}
Here \(\varepsilon\) is the Koszul sign for permuting the suspended inputs.
\end{theorem}
\begin{proof}
The twisting cochain to \(U_\star L\) sends \(sx\) to \(x\) and vanishes in other symmetric degrees.
Its equation~\eqref{ogd:eq:twist} is the envelope relation.
The induced bar map is~\eqref{ogd:eq:CE-map}.
Deconcatenating its permutation sum amounts to choosing a subset and independently ordering that subset and its complement.
This is exactly the unshuffle coproduct followed by \(\Phi\star\Phi\).
There is no factor \(1/r!\) in this coalgebra map.

For two inputs the differential identity is
\[
 b_2\Phi(sx\,sy)=(-1)^{|x|}
 s\bigl(xy-(-1)^{|x||y|}yx\bigr)=(-1)^{|x|}s[x,y].
\]
Grouping adjacent pairs in the full permutation sum gives this identity in every length.
The unary part commutes with the internal differential.
Thus \(\Phi\) is a chain coalgebra map.

Filter the envelope by the number of original \(L\)-factors.
The Poincar\'e--Birkhoff--Witt identity gives
\begin{equation}\label{ogd:eq:PBW}
 \operatorname{gr}U_\star L\simeq\ogdSym_\star L.
\end{equation}
The categorical form of this identity applies here because the tensor is exact and preserves colimits, and characteristic zero makes finite symmetric-group coinvariants exact.
One way to obtain it is the ordinary operadic PBW identity: ordered associative words, modulo the commutator relation, identify with symmetric words in a Lie input.
Its possible overlap on three inputs is exactly Jacobi.
The symmetrization \(\frac1{r!}\sum_{\sigma}\varepsilon(\sigma)\sigma\) gives the associated graded inverse in degree \(r\).
This is an identity of multilinear operations over \(k\), so it remains valid when tensor powers are interpreted by \(\star\).
For derived inputs, apply it to a free simplicial Lie resolution and realize.
Exactness and colimit preservation retain the filtered identity.

Filter the source of~\eqref{ogd:eq:CE-compare} by symmetric degree.
Filter its target by total PBW weight across all bar factors.
Both filtrations terminate at weight \(|I|\) on a configuration \(I\), since weight \(r\) contains \(r\) positive original \(L\)-factors.
The associated graded map is
\begin{equation}\label{ogd:eq:abelian-Koszul}
 \ogdSym_\star(sL)\longrightarrow B(\ogdSym_\star L)
\end{equation}
with only the internal differential on its source.

Here is a direct proof that~\eqref{ogd:eq:abelian-Koszul} is an equivalence.
Set \(S=\ogdSym_\star L\) and consider the free Koszul complex
\(S\star\ogdSym_\star(sL)\).
In addition to the internal differential, give it the derivation \(d_K(sx)=-x\), \(d_K(x)=0\).
The degree-minus-one derivation \(h(x)=-sx\), \(h(sx)=0\) obeys
\[
 d_Kh+hd_K=N,
\]
where \(N\) counts total polynomial weight.
The internal differential anticommutes with \(h\).
On every positive weight, \(h/N\) contracts the augmentation complex.
Thus this complex resolves \(K\) as an \(S\)-module.
Its filtration by suspended degree has free-module subquotients and is finite at each configuration.
It therefore computes derived tensor product with the augmentation module.
The normalized free bar resolution also computes that tensor product.
The comparison between the two resolutions is the permutation sum~\eqref{ogd:eq:CE-map}, extended \(S\)-linearly.
Its differential compatibility was proved above for the abelian bracket.
It induces the identity on the augmentation and therefore is an equivalence between the resolutions.
Tensoring with \(K\) gives~\eqref{ogd:eq:abelian-Koszul}.

Finite filtered comparison now proves~\eqref{ogd:eq:CE-compare} at every configuration.
Joint conservativity proves the claimed coalgebra equivalence in \(\ogdcC\).
\end{proof}

The PBW symmetrization in~\eqref{ogd:eq:PBW} uses \(1/r!\).
The coalgebra comparison~\eqref{ogd:eq:CE-map} uses the unnormalized sum because its coproducts are unshuffle and deconcatenation.
For degree-zero inputs its binary component is \([x|y]-[y|x]\), whose differential is the bracket relation.

The theorem concerns envelopes of Lie objects in the specified convolution category.
It does not identify the bar of every associative algebra with the Chevalley--Eilenberg coalgebra of its commutator Lie algebra.
The construction is the internal version of the classical enveloping-algebra comparison \cite{ogd:CE}.
The Lie--commutative chiral equivalence and the distinct chiral envelope functors in \ogdciteFG*{\S\S5--6} retain their own tensor conventions.

\section{Scope, units and bars}
\label{ogd:sec:status}

Theorem~\ref{ogd:thm:convolution} constructs the ordered convolution category \((\ogdcC,\star,K)\).
Theorems~\ref{ogd:thm:duality} and~\ref{ogd:thm:modules} give its algebra--coalgebra and module--comodule equivalences.
Theorems~\ref{ogd:thm:bimodules} and~\ref{ogd:thm:bimodule-composition} extend the latter to bimodules and their native compositions, with the diagonal equivalence of Proposition~\ref{ogd:prop:diagonal-bimodule} and the coherent multibar maps~\eqref{ogd:eq:successive-multibar}.
Corollary~\ref{ogd:cor:internal-centre} gives the internal \(E_2\)-centre consequence under its additional closed-enrichment, representability and actual \(\ogdcC_R\)-linearity hypotheses.
Theorem~\textup{R22.9.1} retains an arbitrary empty-configuration algebra, and Theorem~\ref{ogd:thm:CE} identifies the bar of an envelope with a Chevalley--Eilenberg coalgebra.
The bar, module and bimodule comparisons use the positive configuration bound of Lemma~\ref{ogd:lem:finite}, as do the relative-empty and envelope comparisons.
They require no additional cohomological boundedness or global elementwise conilpotence assumption.

The\phantomsection\label{ogd:def:ordinary-ran-coefficient-category} same arguments apply to the ordinary Ran carrier, which records the positions but no order inside a collision.
Let \(\mathsf{FinSurj}_+\) be the category of nonempty finite sets and surjections.
Set
\[
 \mathcal R_X=\varprojlim_{\mathsf{FinSurj}_+}\ogdDMod(X^I),
 \qquad \mathcal R_X^+=D(k)\times\mathcal R_X,
\]
where a surjection \(\alpha:I\twoheadrightarrow J\) acts by exceptional pullback along the diagonal \(d_\alpha:X^J\to X^I\).
A Ran family thus consists of an isolated component \(F_\varnothing\in D(k)\), objects \(F_I\), and coherent equivalences \(d_\alpha^!F_I\simeq F_J\), and it is positive when \(F_\varnothing=0\).
For example, a right differential module \(N\) on \(X\) gives the positive main-diagonal family \(\iota_*N\), with \((\iota_*N)_I=\delta_{I,*}N\) for the small diagonal \(\delta_I:X\to X^I\).
Base change along \(d_\alpha\) gives \(d_\alpha^!\delta_{I,*}N\simeq\delta_{J,*}N\).

The convolution on \(\mathcal R_X^+\) is~\eqref{ogd:eq:convolution}, with \(X^I\) in place of \(\mathscr R_I\) and the same cross-block disjoint opens.
The proof of Theorem~\ref{ogd:thm:convolution} applies to these diagonals without change.
Split-fibre summands vanish, saturated partitions descend, and iterated products localize to the same final disjoint locus.
The unit is again \(K\).

The presentable Ran coefficient category also satisfies Proposition~\ref{ogd:prop:enhanced-bar}.
Its positive normalized words obey the same arity bound.
The coherent adjunction and the finite gap proof therefore give the bar--cobar equivalence for empty component $k$
(Proposition~\textup{R26.3.3}).
Section~\textup{R66.1} describes the same category in the correspondence formalism.

Pullback along the maps \(q_I\) compares the two carriers in one direction:
\[
 q^!:\mathcal R_X^+\longrightarrow\ogdcC,
 \qquad (q^!F)_I=q_I^!F_I,\qquad(q^!F)_\varnothing=F_\varnothing.
\]
The identity \(q_I\delta_\alpha=d_\alpha q_J\) makes these components an ordered collision family.
For each partition, the inverse image under \(q_I\) of the ordinary cross-block disjoint locus is its ordered counterpart.
Lemma~\ref{ogd:lem:basechange} and compatibility with external products identify the corresponding convolution summands.
Transitivity gives the same identifications for every refined partition.
Hence \(q^!\) is strong symmetric monoidal, and it preserves colimits because each \(q_I^!\) does.

Four\phantomsection\label{ogd:def:three-tensor-units} units occur, and each one is the unit of a different structure.
The empty unit \(K\) is the unit of \(\star\).
The dualizing coalgebra \(J\) is the unit of \(\ogdotex\), and \(\ogdUU=\Omega J\) is the unit of the transported product \(\diamond\).
A unital chiral algebra \(A\) has a vacuum, the map \(\omega_X\to A\) from the unit chiral algebra.
In the ordered category, vacuum data enter through a chosen module \(P\) and its external endomorphism algebra \(A^{\mathrm{vac}}=R\operatorname{Hom}_A(P,P)\), a differential graded \(k\)-algebra.
The empty unit detects no positive operation: equation~\eqref{ogd:eq:empty-end} gives \(R\operatorname{Hom}_A(K,K)=k\) for every augmented \(A\) with \(A_\varnothing=k\).

The following are separate construction problems, not hypotheses of the theorems above.
\begin{enumerate}
\item A comparison of geometric operation complexes under \(q^!\), beyond
its compatibility with convolution.
\item A vacuum realization: a module \(P\), the algebra \(A^{\mathrm{vac}}\), and its relation to \(BA\).
For the Heisenberg envelope of Section~\textup{R23.3}, the choice \(P=K\) gives \(A^{\mathrm{vac}}=k[\kappa]\), by the argument of~\eqref{ogd:eq:empty-end} over \(k[\kappa]\).
The same choice gives the same algebra for the envelope of the abelian Lie object on the same states.
This problem includes the construction of a module \(P\) whose algebra \(A^{\mathrm{vac}}\) depends on the double-pole collision~\textup{R23.16}.
Such a module must come with a proof that it generates a specified category of modules.
The internal mapping object into \(K\) retains only the empty configuration (Proposition~\ref{ogd:prop:internal-empty}), so it cannot replace \(P\).
A vacuum module need not generate: for the Weyl algebra \(A_h=\mathbb C[h]\langle x,p\rangle/(px-xp-h)\), the vacuum module \(M_0=A_h/A_hp\) satisfies \(R\operatorname{Hom}(M_0,M_\lambda)=0\) for the nonzero modules \(M_\lambda=A_h/A_h(p-\lambda)\) with \(\lambda\in\mathbb C^\times\) (Proposition~\textup{R40.9.4} and the example following it).
\item Comparisons among the bar constructions of Table~\ref{ogd:tab:bars}.
\item A map from the local operators of a field theory on \(X\times\mathbb R\), holomorphic along \(X\) and topological along \(\mathbb R\), to an algebra in \((\ogdcC,\star)\).
The map must preserve the collision operations of Section~\textup{R12.1}, with the completion of coefficients that their expansions at collisions require.
\item The determination of the algebras whose bars satisfy factorization~\eqref{ogd:eq:factorization}.
The algebra \(\ogdUU\) has this property, and \(\diamond\) preserves it.
Proposition~\ref{ogd:prop:sign-exchange} gives an algebra with this property that is not equivalent to an envelope.
A tensor algebra \(T_\star V=\Omega(K\oplus sV)\) on a positive object \(V\) has bar equivalent to \(K\oplus sV\), with zero reduced coproduct.
Hence \(T_\star V\) lacks the property when \(V_{\{1\}}\boxtimes V_{\{1\}}\) is nonzero on the disjoint locus of \(\mathscr R_{\{1,2\}}\).
\end{enumerate}

\subsection{Four bar constructions}

The bar of an associative algebra depends on its tensor product and its unit.
Four constructions in this book use four different pairs.

\begin{table}[ht]
\centering
\small
\begin{tabular}{@{}>{\raggedright\arraybackslash}p{0.25\linewidth}>{\raggedright\arraybackslash}p{0.18\linewidth}>{\raggedright\arraybackslash}p{0.06\linewidth}>{\raggedright\arraybackslash}p{0.39\linewidth}@{}}
\toprule
Construction & Carrier & Unit & Duality statement\\
\midrule
Transverse bar \(B_{\otimes^!}(J,A,J)\) of an \(\otimes^!\)-algebra
 & \((\ogdcC,\ogdotex)\) or \((\mathcal R_X^+,\otimes^!)\) & \(J\)
 & Constructed chartwise on \(\mathcal R_X^+\) (Proposition~\textup{R69.4.1}). On \(\ogdcC\), \((F\ogdotex G)_I=F_I\ogdotex G_I\) (Section~\ref{ogd:sec:diamond}). Duality needs a completeness hypothesis (Theorem~\textup{R26.2.1}).\\
Associative bar \(BA\) of a \(\star\)-algebra
 & \((\ogdcC,\star)\) & \(K\)
 & Inverse to cobar (Theorem~\ref{ogd:thm:duality}).\\
Chevalley--Eilenberg coalgebra \(\ogdCE_\star(L)\) of a Lie object
 & \((\ogdcC,\star)\) & \(K\)
 & Equivalent to \(B(U_\star L)\) (Theorem~\ref{ogd:thm:CE}). This compares a Lie bar with the bar of an envelope only.\\
Ordered-coefficient bar \(B_{C_I}A_I\) of the state algebra \(A_I\) of Chapter~\textup{R24}
 & \(C_I\)-modules, \(C_I=\ogdTW_I\) & \(C_I\)
 & Weight-completed reconstruction (Theorem~\textup{R24.2.3}). The rooted insertion operad of Section~\textup{R24.7.3} uses the top-degree cohomology of \(C_I\), up to an orientation line. A lift of this operad to \(C_I\otimes\mathbb C[I]\), compatible with the differential of \(C_I\) and with the cocompositions, remains to be constructed.\\
\bottomrule
\end{tabular}
\caption{Four bar constructions with their tensor products and units.}
\label{ogd:tab:bars}
\end{table}

The first three rows are the transverse, associative chiral and Lie chiral bars of Section~\textup{R12.1}.
The fourth keeps the collision coefficients \(C_I\) as a separate coefficient ring.
No identification among the four follows from an equality of graded dimensions.
Three further objects are also distinct.
The\phantomsection\label{ogd:def:bar-coalgebra-versus-cobar-resolution} bar coalgebra \(BA\) is the Koszul dual coalgebra of \(A\), and the cobar \(\Omega BA\) is a resolution of \(A\), not a dual.
The algebra \(A^{\mathrm{vac}}\) depends on the chosen module \(P\).
Verdier duality \(\mathbb D\), a contravariant functor on holonomic differential modules, exchanges the open direct image \(j_*\) with \(j_!=\mathbb Dj_*\mathbb D\).

\subsection{Coefficient duality and Koszul dual algebras}
\label{ogd:sec:polarity}

Over a field, the Koszul dual of an augmented algebra is the linear dual of its bar coalgebra.
On the ordered collision carrier, a dual of \(BA\) must retain the positive configurations, which carry the operations of \(A\).
The internal dual into \(K\) retains none of them: \(\underline{R\operatorname{Hom}}(BA,K)\simeq K\) by Proposition~\ref{ogd:prop:internal-empty}.
Verdier duality retains all of them.
On bounded holonomic coefficients, the duality is involutive.
It exchanges \(j_*\) with \(j_!\), so the dual of \(BA\) is an algebra
for a second convolution, assembled from \(j_!\)-extensions
(Theorem~\ref{ogd:thm:polarity}).
The Ran-space structure maps have coalgebra variance by the
open-restriction adjunction \ogdciteFG*{Remark~2.4.8}.
For an open immersion $j:U\hookrightarrow Y$, a morphism
$j^!A\to j^!B$ corresponds to a morphism $A\to j_*j^!B$.
An adjoint map from $j_!j^!A$ instead requires $j_!$ to be
defined on the chosen coefficient. Bounded holonomic coefficients
provide a uniform domain for $j_!=\mathbb D j_*\mathbb D$;
they are not necessary for every individual open immersion.

\begin{example}\label{ogd:ex:open-closed-shriek}
Let $X=\mathbb A^1_k$ and let $j:X\hookrightarrow X\sqcup X$
include the first component. The differential-module category of
the disjoint union is the product of the two component categories.
For arbitrary $M,N_1,N_2\in\ogdDMod(X)$,
\[
 R\operatorname{Hom}_{\ogdDMod(X\sqcup X)}
       ((M,0),(N_1,N_2))
 \simeq R\operatorname{Hom}_{\ogdDMod(X)}(M,N_1).
\]
Since $j^!(N_1,N_2)=N_1$, this proves that $j_!M=(M,0)$
is defined for every $M$. In particular, it is defined for
the free right module $M=\mathcal D_X$. Its associated graded
module is $\mathcal O_{T^*X}$, so its characteristic variety
is all of $T^*X$, of dimension two; it is not holonomic.
\end{example}

This example concerns an open-and-closed immersion. No ordinary
shriek functor on arbitrary coefficients for every open immersion
is assumed. The pro-perfect construction below instead specifies
its own coefficient category and transported functors.

\begin{definition}\label{ogd:def:admissible}
An object \(F\in\ogdcC\) is \emph{duality-admissible} if \(F_\varnothing\) is a perfect complex and each \(F_I\) is, on every chart of \(\mathscr R_I\), a bounded complex whose cohomology is holonomic: coherent, with characteristic variety of dimension at most \(|I|\).
The bounds may depend on \(I\).
Write \(\ogdcC_{\mathrm{exc}}\subset\ogdcC\) for the full subcategory of duality-admissible objects.
\end{definition}

External products, restriction to open subschemes, open direct images and extraordinary inverse images preserve bounded complexes with holonomic cohomology \cite{kd-HTT}*{Proposition~3.2.2 and Theorem~3.2.3}.
Each component of a convolution is a finite sum of objects obtained by these operations.
Hence \(\ogdcC_{\mathrm{exc}}\) is closed under \(\star\), \(\ogdotex\), shifts, retracts, finite limits and finite colimits, and it contains \(K\) and \(J\).
At a configuration \(I\), the normalized bar and cobar are finite iterated extensions of convolution products, by Lemma~\ref{ogd:lem:finite}.
Thus \(B\) and \(\Omega\) preserve \(\ogdcC_{\mathrm{exc}}\), and Theorem~\ref{ogd:thm:duality} restricts to it.

\begin{remark}\label{ogd:rem:current}
A coefficient whose singleton component contains the fields of a free boson is not duality-admissible.
Let \(X=\mathbb A^1\), let \(T\) be an endomorphism of a vector space, and let \(\varphi\) be a vector whose iterates \(T^n\varphi\) are linearly independent.
The current of a free boson, with \(T\) the translation, is an instance.
The right differential module \(\omega_X\otimes k[T]\varphi\), with
\[
 (g\,dz\otimes\psi)\,\partial_z=-(\partial_zg)\,dz\otimes\psi-g\,dz\otimes T\psi,
\]
is free of rank one on \(dz\otimes\varphi\).
Indeed, an operator in normal form \(\sum_n\partial_z^nf_n\) sends \(dz\otimes\varphi\) to \(\sum_n(-1)^nf_n\,dz\otimes T^n\varphi\).
Its characteristic variety is the whole cotangent bundle of \(X\), so it is not holonomic.
Truncation does not help: \(T\) raises the degree in \(T\), so no nonzero finite-dimensional subspace of \(k[T]\varphi\) is \(T\)-stable; for the free boson, \(T\) raises the energy by one.
The bounded holonomic construction therefore does not apply to full current fields.
An injective differential-module map from this free current to a
holonomic module cannot exist: its image would be a nonholonomic
submodule of a holonomic module.
The same argument excludes a map to a bounded holonomic complex
that is injective on $H^0$.
It also excludes an injection into a filtered colimit of holonomic
modules. The free current is finitely presented, so its map would
factor injectively through one holonomic stage
(Proposition~\ref{collar-prop:ind}).

Inverse systems of finite energy quotients have a different scope.
For the orbit $k[T]\varphi$, quotient by $T^{N+1}k[T]\varphi$.
The map to the inverse limit is injective: a nonzero polynomial
survives in every sufficiently large quotient.
An infinite-energy state may nevertheless have unbounded pole order
against a fixed field (\S\ref{collar-section-7}).
These limits therefore require their own field-operation domains.

The full-coefficient duality of Chapter~\ref{rc:full} takes another
route. It keeps finite complexes of finite-rank free modules and
their retracts in inverse diagrams, and transposes the geometric
functors with them.
It gives a duality in a different coefficient category, without
embedding the current in a holonomic module.
\end{remark}

\subsubsection{Pro-perfect coefficients}
\label{ogd:sec:full-coefficients}

The free current has a finite resolution over differential operators:
it is already free of rank one. Holonomicity and finite projective
resolution are different restrictions. The latter gives a dual of
the full current and leads to the coefficient category of
Chapter~\ref{rc:full}.

Let\phantomsection\label{ogd:def:perfect-differential-operator-coefficients} $U$ be a smooth affine coordinate chart of dimension $d$, and put
$R=\Gamma(U,\mathcal D_U)$. A perfect right $R$-complex is a retract
of a finite complex of finite-rank free right modules.
Write $\operatorname{Perf}_r(R)$ for their category.

Choose\phantomsection\label{ogd:def:perfect-coefficient-dual} a coordinate volume form. The formal transpose $t$ satisfies
$t(f)=f$, $t(\partial_i)=-\partial_i$ and $t(PQ)=t(Q)t(P)$.
For a perfect complex define
\[
 P^\dagger=
 \bigl(R\operatorname{Hom}_R(P,R)\bigr)^t[d].
\]

The superscript $t$ converts the left dual module to a right module.
Without a chosen volume, the canonical line makes this conversion.
Finite projective evaluation gives $(P^\dagger)^\dagger\simeq P$:
the second dual reverses the first shift.
On the affine line the free current therefore has dual $R[1]$.

Every derived right $R$-module has a filtered perfect presentation
$M\simeq\operatorname{colim}_\alpha P_\alpha$.
Its stages need not be submodules of $M$.
The dual retains the inverse diagram
\begin{equation}\label{ogd:eq:pro-coefficient}
 M^\sharp=\{P_\alpha^\dagger\}_\alpha
 \quad\text{in}\quad
 \operatorname{Pro}\bigl(\operatorname{Perf}_r(R)\bigr).
\end{equation}
The\phantomsection\label{ogd:def:inverse-pro-perfect-duality} notation $\operatorname{Pro}$ means that the inverse diagram
and its transition maps are kept. The inverse duality sends
$\{Q_\beta\}$ to $\operatorname{colim}_\beta Q_\beta^\dagger$.
Write this inverse as $\mathsf S$.

Write\phantomsection\label{ogd:def:derived-mapping-space} $\operatorname{Map}$ for the derived mapping space.
The duality also retains these spaces between coefficients.
For $N=\operatorname{colim}_\beta Q_\beta$, the pro mapping formula gives
\[
\begin{split}
 \operatorname{Map}_{\operatorname{Pro}}(N^\sharp,M^\sharp)
 &\simeq\lim_\alpha\operatorname{colim}_\beta
       \operatorname{Map}(Q_\beta^\dagger,P_\alpha^\dagger)\\
 &\simeq\lim_\alpha\operatorname{colim}_\beta
       \operatorname{Map}(P_\alpha,Q_\beta)
 \simeq\operatorname{Map}(M,N).
\end{split}
\]
The last equality uses compactness of the perfect complexes.
It proves that the construction is fully faithful and independent
of the presentation. Every pro-perfect diagram arises this way
(Theorem~\ref{full-thm:coeff}).

\begin{example}\label{ogd:ex:pro-localization}
Take $R=k\langle u,\partial\rangle/(\partial u-u\partial-1)$.
Let $L_u$ denote left multiplication by $u$, a right $R$-linear map.
As a right module, the localization is
\[
 R[u^{-1}]=\operatorname{colim}
   \bigl(R\xrightarrow{L_u}R\xrightarrow{L_u}R\longrightarrow\cdots\bigr).
\]
The map from stage $n$ sends $a$ to $u^{-n}a$.

In the left dual, $L_u$ acts by right multiplication by $u$.
Transpose converts it back to $L_u$, so
\[
 R[u^{-1}]^\sharp
 =\{\cdots\xrightarrow{L_u}R[1]\xrightarrow{L_u}R[1]\}.
\]
This pro-object is nonzero: every transition to a fixed stage
is multiplication by a nonzero power of $u$.

Its ordinary inverse limit is zero.
A compatible sequence has every entry in $\bigcap_{n\geq0}u^nR$.
Writing an operator as $\sum_j f_j(u)\partial^j$ shows that this
intersection is zero, since each polynomial $f_j$ would be divisible
by every power of $u$.
Thus the ordinary limit loses the nonzero diagram.
No vanishing of a derived inverse limit is asserted.
\end{example}

Apply this coefficient duality on every chart and retain the
coordinate-change comparisons. It gives a category of pro-perfect
differential-module systems. Denote its functor on $Y$ by $\sharp_Y$
and its inverse by $\mathsf S_Y$.

For an open immersion $j:U\hookrightarrow Y$
and a collision map $\delta:Z\to Y$, its geometric functors are
\begin{equation}\label{ogd:eq:pro-functors}
 j_{!,\mathrm{pro}}=\sharp_Y\circ j_*\circ\mathsf S_U,
 \qquad
 \delta^\dagger=\sharp_Z\circ\delta^!\circ\mathsf S_Y.
\end{equation}
The two contravariant functors in each composite make it covariant.

These are functors of pro-perfect coefficients.
An ordinary shriek extension on all differential modules is not
being assumed. On bounded holonomic coefficients they recover
the corresponding dualized functors.

Write\phantomsection\label{ogd:def:pro-perfect-collision-category} $\ogdcC^\sharp$ for systems with the collision comparisons
$\delta^\dagger$. Componentwise duality gives
$\ogdcC^{\mathrm{op}}\simeq\ogdcC^\sharp$.
The transported convolution is
\begin{equation}\label{ogd:eq:pro-convolution}
 E\star_\sharp F=(\mathsf S E\star\mathsf S F)^\sharp,
 \qquad K^\sharp=K.
\end{equation}
Its open-extension maps are those of \eqref{ogd:eq:pro-functors}.
Associativity and the coherent collision maps are transposed with
the coefficients, as in Proposition~\ref{full-prop:ops}.

For\phantomsection\label{ogd:def:pro-perfect-koszul-algebra} an augmented $\star$-algebra $A$ with $A_\varnothing=k$, put
$\mathsf K_\sharp A=(BA)^\sharp$.
The transpose of the bar coproduct is its multiplication:
\[
 (BA)^\sharp\star_\sharp(BA)^\sharp
 \simeq(BA\star BA)^\sharp
 \xrightarrow{\ \Delta_{BA}^\sharp\ }(BA)^\sharp.
\]
Its unit and augmentation transpose the bar counit and coaugmentation.

The inverse construction sends an augmented $\star_\sharp$-algebra $E$
with $E_\varnothing=k$ to $\mathsf S B_{\star_\sharp}E$.
Finite word length at each configuration gives the double comparison
\[
 \mathsf S B_{\star_\sharp}\bigl((BA)^\sharp\bigr)
 \simeq\Omega BA\xrightarrow{\epsilon_A}A.
\]
This is the full-coefficient algebra anti-equivalence of
Theorem~\ref{full-thm:koszul}.
It retains the current coefficients by changing their category and
transposing their operations. Realization of its output by a second
field system, with its local vacuum and singular products, is additional.

\subsubsection{The holonomic dual convolution}

On the holonomic coefficient category the same change of variance
can be expressed by ordinary shriek functors.

On\phantomsection\label{ogd:def:holonomic-verdier-duality} bounded complexes with holonomic cohomology, \(\mathbb D\) is an involutive contravariant equivalence \cite{kd-HTT}*{Proposition~3.2.1}.
It is local, so it is defined on \(\mathscr R_I\) by descent from the charts; on the empty component it is \(R\operatorname{Hom}_k(-,k)\).
It commutes with restriction to open subschemes.
It commutes with external products, because \(\mathcal D_{Y\times Z}\) is the external product of \(\mathcal D_Y\) and \(\mathcal D_Z\) and duals of coherent modules are computed from locally free resolutions.
It fixes \(\omega_Y\).
The functors \(j_!=\mathbb Dj_*\mathbb D\) and \(\delta_\alpha^*=\mathbb D\delta_\alpha^!\mathbb D\) are those of \cite{kd-HTT}*{Definition~3.2.13}.
These results are proved there over \(\mathbb C\) by algebraic arguments, which apply over any algebraically closed field of characteristic zero; extension of scalars to an algebraic closure gives them over \(k\).

Let\phantomsection\label{ogd:def:holonomic-dual-category} \(\ogdcC_{\mathrm{dual}}\) be defined as \(\ogdcC\) in~\eqref{ogd:eq:category}, with a perfect empty component, bounded components with holonomic cohomology, and the transitions \(\delta_\alpha^*\) in place of \(\delta_\alpha^!\).
Componentwise duality is then an equivalence
\[
 \mathbb D:\ogdcC_{\mathrm{exc}}^{\mathrm{op}}\xrightarrow{\ \sim\ }\ogdcC_{\mathrm{dual}},
 \qquad \mathbb DK=K.
\]
Define\phantomsection\label{ogd:def:holonomic-dual-convolution} \(F\star_!G=\mathbb D(\mathbb DF\star\mathbb DG)\) on \(\ogdcC_{\mathrm{dual}}\).
Since \(\mathbb D\) exchanges \(j_*\) with \(j_!\) and commutes with external products, formula~\eqref{ogd:eq:convolution} becomes
\begin{equation}\label{ogd:eq:shriek-convolution}
 (F\star_!G)_I=
 \bigoplus_{I=A\sqcup B}
 j_{A,B,!}(F_A\boxtimes G_B)|_{U_{A,B}}.
\end{equation}
By transport of structure, \(\star_!\) is symmetric monoidal with unit \(K\), and \(\mathbb D\) is a symmetric monoidal equivalence from \((\ogdcC_{\mathrm{exc}}^{\mathrm{op}},\star)\) to \((\ogdcC_{\mathrm{dual}},\star_!)\).
Write\phantomsection\label{ogd:def:shriek-bar-and-cobar} \(B_!\) and \(\Omega_!\) for the bar and cobar constructions of \(\star_!\); at each configuration they are finite, as for \(\star\).

\begin{construction}\label{ogd:con:koszul-dual}
Let \(A\) be an augmented associative algebra in \((\ogdcC_{\mathrm{exc}},\star)\) with \(A_\varnothing=k\) and augmentation ideal \(a\).
By the definition of coalgebras in Section~\ref{ogd:sec:enhanced-bar}, \(BA\) is an algebra in \((\ogdcC_{\mathrm{exc}}^{\mathrm{op}},\star)\).
Its image
\[
 \mathsf K_*A=\mathbb D(BA)
\]
is therefore an augmented algebra in \((\ogdcC_{\mathrm{dual}},\star_!)\), with \((\mathsf K_*A)_\varnothing=k\).
Its product is the transpose of deconcatenation,
\[
 \mathsf K_*A\star_!\mathsf K_*A\simeq\mathbb D(BA\star BA)
 \xrightarrow{\ \mathbb D\Delta_{BA}\ }\mathsf K_*A,
\]
and its unit and augmentation are the transposes of the counit and the coaugmentation of \(BA\).
At a nonempty configuration \(I\), the filtration of \(BA\) by word length dualizes to a finite filtration of \((\mathsf K_*A)_I\) with associated graded
\[
 \bigoplus_{r=1}^{|I|}\ \bigoplus_{P=(I_1|\cdots|I_r)}
 j_{P,!}\bigl(\mathbb Da_{I_1}\boxtimes\cdots\boxtimes\mathbb Da_{I_r}\bigr)|_{U_P}[-r],
\]
the sum running over ordered decompositions of \(I\) into nonempty parts.
Its differential is the transpose of the bar differential.
Thus the product of \(A\) enters the differential of \(\mathsf K_*A\), and deconcatenation enters its product.
Equivalently, \(\mathbb D\) carries the simplicial bar of \(A\) to the cosimplicial cobar of the \(\star_!\)-coalgebra \(\mathbb DA\), whose reduced coproduct is \(\mathbb D\mu_A\), so that \(\mathsf K_*A\simeq\Omega_!(\mathbb DA)\).
\end{construction}

\begin{theorem}\label{ogd:thm:polarity}
For an augmented algebra \(A'\) in \((\ogdcC_{\mathrm{dual}},\star_!)\) with \(A'_\varnothing=k\), put \(\mathsf K_!A'=\mathbb D(B_!A')\).
\begin{enumerate}
\item For a coaugmented coalgebra \(C\) in \((\ogdcC_{\mathrm{exc}},\star)\) with \(C_\varnothing=k\), there is a natural equivalence of algebras \(\mathbb DB_!(\mathbb DC)\simeq\Omega C\).
\item The functors \(\mathsf K_*\) and \(\mathsf K_!\) are inverse equivalences
\[
 \ogdAlg^{\mathrm{aug},\,A_\varnothing=k}_{E_1}(\ogdcC_{\mathrm{exc}},\star)^{\mathrm{op}}
 \simeq
 \ogdAlg^{\mathrm{aug},\,A'_\varnothing=k}_{E_1}(\ogdcC_{\mathrm{dual}},\star_!).
\]
\item The equivalence \(\mathsf K_!\mathsf K_*A\simeq A\) is the composite of~(i) for \(C=BA\) with the equivalence \(\Omega BA\simeq A\) of Theorem~\ref{ogd:thm:duality}.
\end{enumerate}
\end{theorem}
\begin{proof}
The functor \(\mathbb D\) carries the cosimplicial cobar of \(C\), with terms \(C^{\star q}\), to the simplicial bar of \(\mathbb DC\), with terms \((\mathbb DC)^{\star_!q}\), and it carries totalization to realization.
In normalized terms, cobar words become bar words and the shift \([-q]\) becomes \([q]\).
The terms of the differential that use the coproduct of \(C\) become those that use the product \(\mathbb D\Delta_C\) of \(\mathbb DC\), and concatenation becomes the transpose of deconcatenation.
Both constructions are finite at each configuration by Lemma~\ref{ogd:lem:finite}.
Hence \(B_!(\mathbb DC)\simeq\mathbb D(\Omega C)\) as coalgebras, and~(i) follows from \(\mathbb D^2\simeq1\).

Theorem~\ref{ogd:thm:duality}, restricted to \(\ogdcC_{\mathrm{exc}}\), makes \(B\) an equivalence from augmented algebras with \(A_\varnothing=k\) to coaugmented coalgebras with \(C_\varnothing=k\).
By Section~\ref{ogd:sec:enhanced-bar}, these coalgebras are the opposites of algebras in \((\ogdcC_{\mathrm{exc}}^{\mathrm{op}},\star)\), which \(\mathbb D\) identifies with algebras in \((\ogdcC_{\mathrm{dual}},\star_!)\).
Hence \(\mathsf K_*=\mathbb D\circ B\) is an equivalence onto the opposite category.
Taking \(C=BA\) in~(i) gives \(\mathsf K_!\mathsf K_*A\simeq\Omega BA\simeq A\), naturally in \(A\).
This proves~(iii), and it shows that \(\mathsf K_!\) is inverse to \(\mathsf K_*\), which proves~(ii).
\end{proof}

\begin{example}\label{ogd:ex:dual-unit}
The coalgebra \(J\) is duality-admissible, and \(B\ogdUU\simeq J\) by Theorem~\ref{ogd:thm:duality}.
Hence
\[
 \mathsf K_*\ogdUU\simeq\mathbb DJ,\qquad
 (\mathbb DJ)_I=\omega_{\mathscr R_I}[-|I|],\qquad(\mathbb DJ)_\varnothing=k.
\]
Its product is the transpose of the localization coproduct~\eqref{ogd:eq:J-coprod}.
On the summand of an ordered decomposition \(P\), it is the counit \(j_{P,!}j_P^!(\mathbb DJ)_I\to(\mathbb DJ)_I\), the identity on \(U_P\).
Take the chart of~\eqref{ogd:eq:opening-binary}, suppress densities and shifts, and restrict to the summand of \((\{1\}|\{2\})\).
There the coproduct of \(J\) is the localization \(R\hookrightarrow L\), and the product of \(\mathbb DJ\) is its transpose, the surjection in the exact sequence
\[
 0\longrightarrow L/R\longrightarrow\mathbb DL\longrightarrow R\longrightarrow0.
\]
Its kernel \(\mathbb D(L/R)\simeq L/R\) is supported on the diagonal, because \(L/R\) is the direct image of the functions on the diagonal and duality commutes with proper direct image \cite{kd-HTT}*{Theorem~2.7.2}.
Unlike \(\underline{R\operatorname{Hom}}(B\ogdUU,K)\simeq K\), the dual \(\mathsf K_*\ogdUU\) retains every positive configuration.
\end{example}

The exchange of \(j_*\) for \(j_!\) cannot be avoided.
On the underlying object of \(\mathsf K_*\ogdUU\), a product built from \(j_*\)-extensions would have to agree with the product of \(\mathsf K_*\ogdUU\) away from the diagonal, and the localization \(R\to L\) forbids this.

\begin{lemma}\label{ogd:lem:hom-zero}
For a noetherian domain \(R\) and a nonzero nonunit \(u\in R\), \(\operatorname{Hom}_R(R[u^{-1}],R)=0\).
\end{lemma}
\begin{proof}
Let \(f\) be such a map and put \(x_n=f(u^{-n})\), so that \(x_n=ux_{n+1}\).
Fix \(m\geq0\).
The ideals \(x_mR\subseteq x_{m+1}R\subseteq\cdots\) stabilize, so \(x_{n+1}=rx_n=rux_{n+1}\) for some \(n\geq m\) and \(r\in R\).
Since \(u\) is not a unit, \(1-ru\neq0\); hence \(x_{n+1}=0\), and \(x_m=u^{n+1-m}x_{n+1}=0\).
The elements \(u^{-m}\) generate \(R[u^{-1}]\) as an \(R\)-module, so \(f=0\).
\end{proof}

\begin{corollary}\label{ogd:cor:no-extension}
Let \(j:U\hookrightarrow\mathscr R_{\{1,2\}}\) be the disjoint locus, and let \(E\) be \(J_{\{1,2\}}\) or \((\mathbb DJ)_{\{1,2\}}\).
No map \(j_*j^!E\to E\) restricts to an isomorphism on \(U\).
Consequently, the counit \(m_!:j_!j^!E\to E\) does not factor through the canonical map \(\eta:j_!j^!E\to j_*j^!E\).
For the cone \(Q\) of \(\eta\), the class \(\partial[m_!]\in H^1R\operatorname{Hom}(Q,E)\) is nonzero.
\end{corollary}
\begin{proof}
The objects \(j_*j^!E\) and \(E\) are the same shift of the modules \(j_*\omega_U\) and \(\omega\), so a map between them is a map of differential modules, in particular \(\mathcal O\)-linear.
Let \(x\) be a point of a diagonal branch, and trivialize \(\omega\) on a neighborhood of \(x\) in a chart.
The map becomes an \(\mathcal O\)-linear map \(\mathcal O(*\Delta)\to\mathcal O\).
Its stalk at \(x\) is a map \(R[u^{-1}]\to R\), where \(R\) is the local ring at \(x\), a noetherian domain, and \(u\) is an equation of the diagonal.
By Lemma~\ref{ogd:lem:hom-zero} the stalk vanishes, so the image of the section \(1\) vanishes near \(x\).
An isomorphism on \(U\) would send \(1\) to a unit on \(U\), and \(U\) meets every neighborhood of \(x\).
If \(m_!=m\circ\eta\), then \(m\) restricts on \(U\) to \(m_!|_U\), the identity, since \(\eta\) is the identity on \(U\); this is excluded.
The long exact sequence of \(R\operatorname{Hom}(-,E)\) on the triangle \(j_!j^!E\xrightarrow{\eta}j_*j^!E\to Q\) shows that \(\partial[m_!]\) vanishes exactly when \(m_!\) factors through \(\eta\).
\end{proof}

The obstruction lies in the binary component of the product, which higher associativity data do not change.
A different carrier, or a realization that changes the coefficient object, is not excluded.

Verdier\phantomsection\label{ogd:def:verdier-diagonal-tensor} duality also transports the diagonal tensor of Section~\ref{ogd:sec:diamond}.
For \(F,G\in\ogdcC_{\mathrm{dual}}\), put
\[
 F\ogdotexd G=\mathbb D(\mathbb DF\ogdotex\mathbb DG),\qquad
 (F\ogdotexd G)_I=\Delta^*(F_I\boxtimes G_I),\qquad
 \Delta^*=\mathbb D\Delta^!\mathbb D.
\]
It is symmetric monoidal with unit \(\mathbb DJ\).

\begin{proposition}\label{ogd:prop:dual-tensor}
\begin{enumerate}
\item Verdier duality carries the interchange~\eqref{ogd:eq:interchange} to a natural map
\[
 \mathbb D\zeta:(F\ogdotexd H)\star_!(G\ogdotexd L)\longrightarrow(F\star_!G)\ogdotexd(H\star_!L).
\]
\item Augmented \(\star_!\)-algebras with empty component \(k\) form a symmetric monoidal category under \(\ogdotexd\), with unit \(\mathbb DJ\); the product of \(A\ogdotexd A'\) is \((\mu_A\ogdotexd\mu_{A'})\circ\mathbb D\zeta\).
\item The anti-equivalence \(\mathsf K_*\) carries \(\diamond\) and its unit \(\ogdUU\) to \(\ogdotexd\) and its unit \(\mathbb DJ\), with all coherences.
Hence it exchanges monoids for \(\diamond\) with comonoids for \(\ogdotexd\) in \(\star_!\)-algebras.
\end{enumerate}
\end{proposition}
\begin{proof}
By definition, componentwise duality exchanges \(\star\) with \(\star_!\) and \(\ogdotex\) with \(\ogdotexd\).
Applied to the map \(\zeta\) for the objects \(\mathbb DF,\mathbb DG,\mathbb DH,\mathbb DL\), it reverses the arrow and gives~(i).
The category \(\ogdcC_{\mathrm{exc}}\) contains \(J\) and is closed under \(\ogdotex\), so Proposition~\ref{ogd:prop:coalgtensor} holds there.
Duality carries its coalgebras, coproducts and constraints to the algebras, products and constraints of~(ii).
For~(iii), \(B(A\diamond A')\simeq BA\ogdotex BA'\) by~\eqref{ogd:eq:diamond}, so
\[
 \mathsf K_*(A\diamond A')\simeq\mathbb D(BA\ogdotex BA')=\mathsf K_*A\ogdotexd\mathsf K_*A',
\]
and \(\mathsf K_*\ogdUU\simeq\mathbb DJ\) by Example~\ref{ogd:ex:dual-unit}.
The coherences are transported in the same way.
A monoidal equivalence carries monoids to monoids, and the contravariance of \(\mathsf K_*\) turns them into comonoids.
\end{proof}

For\phantomsection\label{ogd:def:dual-unit-monoid-comonoid} example, the unit \(\ogdUU\) is a monoid for \(\diamond\), and \(\mathbb DJ\) is the corresponding comonoid.
The exchange of \(j_*\) for \(j_!\) is forced at the level of tensor structures as well.

\begin{proposition}\label{ogd:prop:no-interchange}
There are no maps
\[
 \chi:(F\ogdotex H)\star(G\ogdotex L)\longrightarrow(F\star G)\ogdotex(H\star L)
\]
that restrict on each disjoint locus to the identification of external products and for which \(\ogdotex\), with unit \(J\) and its componentwise unit constraints, lifts to a monoidal structure on a full subcategory of augmented \(\star\)-algebras containing \(\ogdUU\), with the product of \(A\ogdotex A'\) equal to \((\mu_A\ogdotex\mu_{A'})\circ\chi\).
\end{proposition}
\begin{proof}
Suppose such maps exist.
The unit \(J\) of the lifted structure is then a \(\star\)-algebra with a product \(\mu_J\), and the unit constraint \(J\ogdotex\ogdUU\simeq\ogdUU\) is a map of algebras.
Let \(U\subset\mathscr R_{\{1,2\}}\) be the disjoint locus, and restrict to the summand of the ordered decomposition \((\{1\}|\{2\})\).
There \(\chi\) is the identification of external products.
Over \(U\) the complex \(\ogdUU_{\{1,2\}}\) is \([\omega_U\xi\to\omega_Ue_{12}\oplus\omega_Ue_{21}]\) by~\eqref{ogd:eq:U-partitions}, so the product \(a_1\boxtimes a_2\mapsto e_{12}\) of \(\ogdUU\) is an equivalence there.
The component of \(\mu_J\) is an endomorphism of \(j^!J_{\{1,2\}}=\omega_U[2]\), namely multiplication by a function \(c\) with \(dc=0\).
Compatibility of the unit constraint with the products multiplies the equivalence by \(c\) on one side and leaves it unchanged on the other, so \(c=1\).
Hence the component \(j_*j^!J_{\{1,2\}}\to J_{\{1,2\}}\) of \(\mu_J\) restricts to the identity on \(U\), which Corollary~\ref{ogd:cor:no-extension} excludes.
\end{proof}

A diagonal tensor with unit \(J\) is therefore not carried to \(\star\)-algebras by an interchange that is the identity away from collisions.
Its Verdier dual, with unit \(\mathbb DJ\), is carried to \(\star_!\)-algebras by Proposition~\ref{ogd:prop:dual-tensor}, and the unit components are the counits of Example~\ref{ogd:ex:dual-unit}.

\begin{openproblem}\label{ogd:op:realization}
Theorem~\ref{ogd:thm:polarity} is a duality of associative algebras between \(\ogdcC_{\mathrm{exc}}\) and \(\ogdcC_{\mathrm{dual}}\).
Two maps are missing.
\begin{enumerate}
\item An operation-preserving realization of \(\mathsf K_*A\) by singular associative fields, which retains the vacuum and the common expansion domains of the fields and carries \(\ogdotexd\) to an independently specified tensor product of fields.
\item A comparison of \(\diamond\), which is defined by transport along \(B\), with an independently specified tensor product \(\odot\) of \(\star\)-algebras.
The specification of \(\odot\) includes an exterior product of modules \(\ogdMod_A(\ogdcC)\times\ogdMod_{A'}(\ogdcC)\to\ogdMod_{A\odot A'}(\ogdcC)\).
Under the comparison, a coproduct \(A\to A\odot A\) must satisfy the bialgebra identity, and the resulting tensor product of \(A\)-modules must satisfy the intertwiner identities.
By Proposition~\ref{ogd:prop:dual-tensor}(iii), \(\mathsf K_*\) carries a \(\diamond\)-comonoid structure on \(A\) to an \(\ogdotexd\)-monoid structure on \(\mathsf K_*A\).
Without the comparison, a coproduct \(A\to A\odot A\) and the intertwiners between tensor products of \(A\)-modules are not known to be carried by \(\mathsf K_*\).
\end{enumerate}
\end{openproblem}

\subsection{The envelope sector}
\label{ogd:sec:envelope-sector}

The algebra \(\ogdUU=\Omega J\), the envelopes of Section~\ref{ogd:sec:CE}, and the tensor algebras on positive objects lie in the Lie-envelope sector.
Their bars are equivalent to cocommutative coalgebras.
The next proposition proves these assertions and gives a conditional closure statement for \(\diamond\), with a compatible strict cocommutative model as an explicit hypothesis.

\begin{proposition}\label{ogd:prop:envelope-sector}
\begin{enumerate}
\item Let \(C\) be a connected coaugmented \(\star\)-coalgebra with a strict model whose coproduct is cocommutative.
Let \(L_C\subset T_\star(tc)\) be the free Lie object on \(tc\).
The cobar differential preserves \(L_C\), and \(\Omega C=U_\star(L_C)\).
Hence \(C\simeq B(\Omega C)\simeq\ogdCE_\star(L_C)\) as coalgebras.
\item The\phantomsection\label{ogd:def:dualizing-lie-cobar} algebra \(\ogdUU=\Omega J\) is the envelope \(U_\star(L_J)\), where \(L_J\) is the free Lie object on the generators \(\xi_B\) and
\begin{equation}\label{ogd:eq:J-Lie}
 d\xi_B=\sum_{\substack{B=A\sqcup C,\ A,C\ne\varnothing\\ \min B\in A}}
 (-1)^{|A|+1}\varepsilon(A,C)\,[\xi_A,\xi_C].
\end{equation}
In particular \([a_1,a_2]=d\xi_{12}\), and the singular class~\eqref{ogd:eq:singular} is the localized image of this bracket.
\item The tensor algebra \(T_\star V\) on a positive object \(V\) is the envelope of the free Lie object on \(V\).
\item Let \(A=U_\star L\) and \(A'=U_\star L'\), where \(L,L'\) are positive Lie objects, and put
\[
 E=\ogdCE_\star(L)\ogdotex\ogdCE_\star(L').
\]
Give \(E\) the coaugmented counital \(\star\)-coalgebra structure induced by~\eqref{ogd:eq:interchange}.
Assume that this derived coalgebra admits a strict cocommutative representative \(E_{\mathrm{str}}\).
The representative and its equivalence to \(E\) must include the ordered collision maps and the induced coproduct on one presentation by differential modules, not only an equivalence of underlying tensor complexes.
Then there are equivalences of augmented \(\star\)-algebras
\[
 A\diamond A'\simeq\Omega E\simeq\Omega E_{\mathrm{str}}
       =U_\star(L_{E_{\mathrm{str}}}),
\]
where \(L_{E_{\mathrm{str}}}\) is the positive Lie cobar object of part~\textup{(i)}.
Thus closure under \(\diamond\) holds under the stated strict-model hypothesis.
\end{enumerate}
\end{proposition}

\begin{proof}
(i) Write \(\bar\Delta c=\sum c'\otimes c''\).
Cocommutativity pairs the term \(c'\otimes c''\) with \((-1)^{|c'||c''|}c''\otimes c'\).
In~\eqref{ogd:eq:cobar-signs} the two terms of such a pair have coefficients \((-1)^{|c'|+1}\) and \((-1)^{|c''|+1+|c'||c''|}\).
Their ratio is \((-1)^{|c'|+|c''|+|c'||c''|}=-(-1)^{|tc'||tc''|}\), since \(|tc|=|c|+1\).
Hence
\[
 d(tc)=-t(d_Cc)+\tfrac12\sum(-1)^{|c'|+1}[tc',tc''],
\]
a sum of graded commutators.

A derivation of \(T_\star(tc)\) that maps the generators into \(L_C\) preserves \(L_C\), because derivations preserve graded commutators.
Thus \((L_C,d)\) is a positive differential graded Lie object.
The envelope of a free Lie object is the tensor algebra on its generators, since both represent maps of the generators into associative algebras.
The differentials agree on generators.
Hence \(\Omega C\) is the strict envelope of \(L_C\): the tensor algebra modulo the relations \(xy-(-1)^{|x||y|}yx=[x,y]\), formed without cofibrant replacement.

The Poincar\'e--Birkhoff--Witt identity~\eqref{ogd:eq:PBW} holds for strict envelopes, with a filtration that is finite at each configuration.
In characteristic zero, \(\ogdSym_\star\) preserves quasi-isomorphisms.
Hence the strict envelope preserves quasi-isomorphisms and computes the derived envelope, so \(\Omega C=U_\star(L_C)\).
Theorem~\ref{ogd:thm:duality} gives \(C\simeq B\Omega C\), and Theorem~\ref{ogd:thm:CE} gives \(B(U_\star L_C)\simeq\ogdCE_\star(L_C)\).

(ii) The coproduct~\eqref{ogd:eq:J-coprod} of \(J\) is cocommutative, so (i) applies.
Pairing the splittings \((A,C)\) and \((C,A)\) of~\eqref{ogd:eq:partition-sign} as in (i) gives~\eqref{ogd:eq:J-Lie}.
Formula~\eqref{ogd:eq:binary-sign} is its first case.

(iii) Both \(T_\star V\) and the envelope of the free Lie object on \(V\) represent maps from \(V\) into associative algebras.

(iv) Theorem~\ref{ogd:thm:CE} gives coalgebra equivalences
\(BA\simeq\ogdCE_\star(L)\) and \(BA'\simeq\ogdCE_\star(L')\).
By the definition of \(\diamond\) and preservation of equivalences,
\[
 A\diamond A'\simeq\Omega E\simeq\Omega E_{\mathrm{str}}.
\]
Part~\textup{(i)} applies to the specified strict cocommutative representative and identifies the last algebra with \(U_\star(L_{E_{\mathrm{str}}})\).
On a compatible strict coefficient presentation the common-decomposition interchange preserves cocommutativity: exchanging the outputs exchanges the same ordered decomposition in both factors.
This verifies the coproduct identity on such a presentation; the existence of that presentation is the hypothesis of part~\textup{(iv)}.
\end{proof}

Corollary~\ref{lde:cor:diamond} proves closure under \(\diamond\) for arbitrary positive \(L,L'\) using coherent cocommutative coalgebras.
Part~\textup{(iv)} gives the explicit Lie-cobar description when a compatible strict representative is supplied; constructing such representatives is a separate problem.
If the bar of an augmented algebra admits a compatible strict cocommutative representative, part~\textup{(i)} and bar--cobar duality identify that algebra with an envelope.
An invariant that vanishes on all envelopes can therefore detect associative algebras outside this sector.
For coalgebras that satisfy factorization~\eqref{ogd:eq:factorization}, there is such an invariant.
At two points, factorization makes both ordered binary components of the coproduct equivalences on the disjoint locus.
Their comparison is then an automorphism there.

Let\phantomsection\label{ogd:def:factorizing-exchange-input} \(C\) be a coaugmented \(\star\)-coalgebra that satisfies factorization, with reduced part \(c\) and singleton part \(N=c_{\{1\}}\).
Let \(U\subset\mathscr R_{\{1,2\}}\) be the disjoint locus.
Write\phantomsection\label{ogd:def:ordered-binary-coproduct-components} \(\Delta_{(1|2)}\) and \(\Delta_{(2|1)}\) for the two components of the reduced coproduct on \(j^!c_{\{1,2\}}\).
Compose the second with the symmetry of \(\boxtimes\), so that both components take values in \((N\boxtimes N)|_U\).
The \emph{exchange} of \(C\) is the automorphism
\begin{equation}\label{ogd:eq:exchange}
 R_C=\Delta_{(2|1)}\Delta_{(1|2)}^{-1}\in\operatorname{Aut}\bigl((N\boxtimes N)|_U\bigr).
\end{equation}
\begin{proposition}\label{ogd:prop:exchange}
Let \(C\) be a coaugmented \(\star\)-coalgebra that satisfies factorization.
\begin{enumerate}
\item \(R_C=\mathrm{id}\) if and only if the binary reduced coproduct is cocommutative on \(U\).
In particular \(R_C=\mathrm{id}\) when \(C\) is cocommutative.
\item Let \(f:C\to C'\) be an equivalence of coalgebras.
Then \(C'\) satisfies factorization, and \(R_{C'}(f_1\boxtimes f_1)=(f_1\boxtimes f_1)R_C\) in the homotopy category of differential modules on \(U\), where \(f_1\) is the singleton component of \(f\).
The condition \(R_C=\mathrm{id}\) is invariant under equivalence.
\end{enumerate}
\end{proposition}

\begin{proof}
(i) By definition, \(R_C=\mathrm{id}\) says that the two binary components agree after the symmetry.
This is cocommutativity of the binary component on \(U\).

(ii) The components of \(f\) are equivalences, so the factorization maps of \(C'\) are equivalences by the two-out-of-three property.
The equivalence \(f\) commutes with the coproducts up to coherent homotopy.
Restriction to \(U\) and projection to a component are functors, so \(f\) intertwines \(\Delta_{(1|2)}\) and \(\Delta_{(2|1)}\) in the homotopy category.
Composing one intertwining relation with the inverse of the other gives the identity.
\end{proof}

\begin{corollary}\label{ogd:cor:envelope-exchange}
If \(A\) is equivalent to an envelope \(U_\star L\) and \(BA\) satisfies factorization, then \(R_{BA}=\mathrm{id}\).
An augmented algebra whose bar satisfies factorization with \(R_{BA}\ne\mathrm{id}\) is not equivalent to an envelope.
The exchange \(R_{BA}\) depends only on the multiplication of \(A\), and a twist of a coproduct on \(A\) does not change it.
\end{corollary}

\begin{proof}
Theorem~\ref{ogd:thm:CE} gives \(BA\simeq\ogdCE_\star(L)\), whose unshuffle coproduct is cocommutative.
Proposition~\ref{ogd:prop:exchange} gives \(R_{BA}=\mathrm{id}\).
The bar is defined from \((A,\mu,\eta)\) alone.
\end{proof}

The proved examples in the envelope sector of \((\ogdcC,\star)\) include \(\ogdUU=U_\star(L_J)\), the envelopes of Section~\ref{ogd:sec:CE}, and the tensor algebras \(T_\star V\).
Products under \(\diamond\) of such envelopes again lie in this sector by Corollary~\ref{lde:cor:diamond}.
The same holds for the current envelope of Theorem~\textup{R23.4.1}, and, in the \(k[\kappa]\)-linear convolution category, for the Heisenberg envelope of Section~\textup{R23.3}.
The triangular Hopf algebra of Theorem~\textup{R23.4.2} keeps the multiplication of an envelope and twists only its coproduct.
Its \(R\)-matrix is supported at the empty configuration, so it does not depend on positions.

Some algebras whose bars satisfy factorization lie outside the envelope sector.
Let\phantomsection\label{ogd:def:pulled-back-main-diagonal-object} \(N\) be a right differential module on \(X\), and let \(V=q^!\iota_*N\) be the pullback of its main-diagonal family.
For \(|I|\geq2\) and \(\tau\in\ogdOrd(I)\), let \(\delta_\tau:X\to\mathscr R_I\) be the small diagonal of the chart \(\tau\).
For \(\rho\ne\tau\), the charts \(\rho\) and \(\tau\) overlap in the complement of the diagonals \(\Delta_{ab}\) of the pairs that they order differently.
At least one such pair exists, so this overlap misses the image of \(\delta_\tau\).

Hence \(\delta_\tau\) is a closed immersion, and distinct orders have disjoint images.
Since \(q_I\) is a local isomorphism, the preimage of the small diagonal of \(X^I\) is the disjoint union of these images.
Base change therefore gives \(V_\varnothing=0\), \(V_{\{1\}}=N\) and
\begin{equation}\label{ogd:eq:diagonal-object}
 V_I=\bigoplus_{\tau\in\ogdOrd(I)}\delta_{\tau,*}N\qquad(|I|\geq2).
\end{equation}

Let\phantomsection\label{ogd:def:signed-invariant-coalgebra} \(\Sigma_r\) act on \((sV)^{\star r}\) by the symmetry of \(\star\) multiplied by the sign character.
In characteristic zero the invariants are the images of the averaging idempotents.
The product \(\star\) carries two such idempotents to the averaging idempotent of \(\Sigma_p\times\Sigma_q\).
Deconcatenation carries \(\Sigma_{p+q}\)-invariants to \(\Sigma_p\times\Sigma_q\)-invariants.
The invariants therefore form a connected coaugmented subcoalgebra \(C\) of the tensor coalgebra \(K\oplus\bigoplus_{r\geq1}(sV)^{\star r}\) with deconcatenation.

\begin{proposition}\label{ogd:prop:sign-exchange}
Let \(C\) be the coalgebra of these invariants for the object \(V\) of~\eqref{ogd:eq:diagonal-object}.
Then \(C\) satisfies factorization, and \(R_C=-\mathrm{id}\) on \((sN\boxtimes sN)|_U\).
Suppose that \((N\boxtimes N)|_U\ne0\), for instance \(N=\omega_X\).
Then \(A=\Omega C\) is an augmented algebra with \(A_\varnothing=k\) whose bar satisfies factorization and has \(R_{BA}=-\mathrm{id}\ne\mathrm{id}\).
In particular, \(A\) is not equivalent to an envelope.
\end{proposition}

\begin{proof}
The group \(\Sigma_r\) permutes the summands of~\eqref{ogd:eq:many-convolution} for \((sV)^{\star r}\) freely.
Choose one ordering of the blocks of each set partition \(\pi\) of \(I\) into \(r\) blocks.
Projection to the corresponding summands \(j_{\pi,*}(\boxtimes_{B\in\pi}sV_B)|_{U_\pi}\) identifies the invariants at \(I\) with their sum.

Let \(P=(I_1|\cdots|I_s)\) be an ordered partition of \(I\).
If a block \(B\) of \(\pi\) meets two parts of \(P\), then \(V_B\) is supported where the points of \(B\) coincide.
Lemma~\ref{ogd:lem:basechange} then shows that \(j_P^!\) kills the summand of \(\pi\).
The remaining partitions refine \(P\), and they are the unions of partitions of the parts \(I_j\).
The \(P\)-component of the iterated deconcatenation maps each remaining summand, up to sign, onto the corresponding summand of \((C_{I_1}\boxtimes\cdots\boxtimes C_{I_s})|_{U_P}\), which proves factorization.

Now take \(I=\{1,2\}\).
The length-one term \(sV_{\{1,2\}}\) is supported on the diagonal branches, so \(j^!\) kills it.
The map \(u\mapsto u-\sigma u\), with \(\sigma\) the symmetry of \(\star\), identifies the summand of the ordered decomposition \((\{1\},\{2\})\) in \(j^!(sV\star sV)_{\{1,2\}}\) with \(j^!C_{\{1,2\}}\).
The two components of the coproduct send \(u-\sigma u\) to \(u\) and to \(-\sigma u\).
After the symmetry of \(\boxtimes\) the second becomes \(-u\), so \(R_C=-\mathrm{id}\).

Theorem~\ref{ogd:thm:duality} gives an equivalence \(C\simeq BA\), and Proposition~\ref{ogd:prop:exchange}(ii) shows that \(BA\) satisfies factorization with \(R_{BA}=-\mathrm{id}\).
This differs from the identity because \((sN\boxtimes sN)|_U\ne0\) and \(k\) has characteristic zero.
Corollary~\ref{ogd:cor:envelope-exchange} gives the last assertion.
\end{proof}

The coalgebra \(C\) is cocommutative for a different symmetry.
Grade the objects \((sV)^{\star r}\) by the word length \(r\).
On objects of word lengths \(p\) and \(q\), multiply the symmetry of \(\star\) by \((-1)^{pq}\).
This gives a second symmetric monoidal structure on graded objects, whose permutation action on \((sV)^{\star r}\) is the action used above.
For this modified symmetry, the symmetrization map~\eqref{ogd:eq:CE-map} identifies \(C\) with the Chevalley--Eilenberg coalgebra of \(V\) with zero bracket.
It matches the unshuffle coproduct with deconcatenation.

The proof of Theorem~\ref{ogd:thm:CE} uses only the symmetric monoidal structure of the convolution, its exactness and the positive configuration bound.
It therefore applies to the modified symmetry.
It identifies \(C\) with the bar of the envelope, for that symmetry, of \(V\) with zero bracket.
Theorem~\ref{ogd:thm:duality} then identifies \(A\) with this envelope.
The exchange \(-\mathrm{id}\) records the sign by which the two symmetries differ on two factors.

More\phantomsection\label{ogd:def:bicharacter-graded-coefficients} generally, let \(N=\bigoplus_{\gamma\in\Gamma}N_\gamma\) be graded by an abelian group \(\Gamma\), so that \(V\) is graded as well.
Let \(\chi\) be a bicharacter of \(\Gamma\) with \(\chi(\gamma,\delta)\chi(\delta,\gamma)=1\).
Multiply the symmetry by \(\chi\), and take invariants in the same way.
The same proof shows that the resulting coalgebra satisfies factorization.
Its exchange acts by \(\chi(\gamma,\delta)\) on the component in which the label \(1\) has degree \(\gamma\) and the label \(2\) has degree \(\delta\).
The coalgebra \(C\) is the case \(\Gamma=\mathbb Z\), with \(N\) in degree one and \(\chi(p,q)=(-1)^{pq}\).

The same proof with the unmodified symmetry shows that the cofree cocommutative coalgebra on \(sV\) also satisfies factorization.
Its exchange is the identity.
A set partition of \(I\) with an order on each block determines a collision stratum of \(\mathscr R_I\).
On this stratum, the points of each block coincide on the branch of its order, and points of different blocks are distinct.
For \(I=\{1,2\}\), the strata are the disjoint locus \(U\) and the two diagonal branches.
In both coalgebras, each summand of a chart \(C_I\) is the open direct image of an object supported on one stratum.

Let\phantomsection\label{ogd:def:closed-collision-strata-union} \(Z\) be a closed union of strata.
By Lemma~\ref{ogd:lem:basechange}, each summand is supported in \(Z\) or has zero exceptional restriction to \(Z\).
Every recollement triangle~\eqref{ogd:eq:recollement} along \(Z\) therefore splits, and its connecting map vanishes.
Suppose that \((N\boxtimes N)|_U\ne0\).
Then the two coalgebras have the same singleton part and vanishing connecting maps, but their exchanges are \(-\mathrm{id}\) and \(\mathrm{id}\).
The connecting maps therefore do not determine the exchange.

It is an open problem to determine which automorphisms of \((c_{\{1\}}\boxtimes c_{\{1\}})|_U\) occur as exchanges \(R_C\) of coalgebras \(C\) that satisfy factorization.
